%% Volume I, Chapter 19: Eigenvalues and Eigenvectors
\chapter{Eigenvalues and Eigenvectors}
\label{ch:eigenvalues}

\abstract{What happens when a matrix is applied again and again? Population models, Markov chains, linear recurrences, iterative methods and the stability of differential equations all ask this question. Most directions are turned as well as stretched, but some are only stretched, and in a basis of such directions the matrix becomes diagonal and its powers become powers of numbers. The stretching factors are the eigenvalues. They are the roots of a polynomial, they can be complex even for a real matrix, and the largest of them in modulus decides the behaviour over long times. Finding them is harder than it looks: beyond degree four there is no formula, and the obvious route through the polynomial can lose every digit.}

Eigenvalues have been in view for a long time. \Cref{sec:extrema-n} defined them in passing, to state the criteria for definiteness in \cref{thm:definiteness-criteria}, and the method of Lagrange multipliers in \cref{sec:lagrange} found them as the extreme values of a quadratic form on the sphere. \Cref{ch:vectors-matrices} met vectors that a matrix maps to multiples of themselves, the sine vectors of the second difference matrix $T_n$, saw the powers of the Fibonacci matrix grow like $\phi^k$, and observed that the Neumann series can converge although $\|A\|\ge1$. \Cref{ch:linear-maps} asked in which basis the matrix of an operator is as simple as possible, found that a diagonal matrix needs a basis of vectors that the operator merely stretches, and showed that the rotation $R_\varphi$ has no such basis over $\mathbb R$ but one over $\mathbb C$. \Cref{ch:determinants} supplied the determinant, which decides when $A-\lambda I$ is singular, and \cref{ch:inner-products} found the axis of a rotation of space as a vector fixed by it. This chapter takes up all these threads.

We begin with the definitions for operators and matrices and with examples from geometry and from spaces of functions. The eigenvalues of a matrix are the roots of its characteristic polynomial $\det(\lambda I-A)$, whose coefficients contain the trace and the determinant; by the fundamental theorem of algebra a complex $n\times n$ matrix has $n$ eigenvalues counted with multiplicity, and conversely every polynomial is the characteristic polynomial of a matrix. Eigenvectors for different eigenvalues are independent, and a matrix can be diagonalised exactly when it has enough of them; then its powers, and every polynomial in it, are computed from powers of numbers, and the largest modulus of an eigenvalue, the spectral radius, decides whether the powers grow or decay. The Cayley--Hamilton theorem says that every matrix satisfies its own characteristic polynomial, and the minimal polynomial gives a second criterion for diagonalisability. Complex eigenvalues of real matrices describe rotations, the eigenvalues of unitary matrices lie on the unit circle and those of Hermitian matrices on the real line. A section of structured examples computes the eigenvalues of the second difference matrix, of circulant matrices and of rank-one matrices exactly, and the Markov matrices of probability lead to Perron's theorem and to the localisation of eigenvalues by Gershgorin's discs. The last section is a first look at computation: the power method, the reason why eigenvalues are never computed from the characteristic polynomial, and the sensitivity of eigenvalues to perturbations.

As in the preceding chapters, $\mathbb F$ denotes $\mathbb R$ or $\mathbb C$. A real matrix is also a complex matrix, and its complex eigenvalues are often needed even when the matrix is real; we say explicitly when only real eigenvalues are meant. The spectral theorem for Hermitian matrices is the subject of \cref{ch:symmetric-matrices}, the singular value decomposition of \cref{ch:svd}, and the Schur and Jordan forms of Chapter~22; none of them is used here, and we indicate where they complete the results.

The chapter is organised as follows. \Cref{sec:eigen-definitions} defines eigenvalues, eigenvectors and eigenspaces and collects examples. \Cref{sec:characteristic-polynomial} introduces the characteristic polynomial, its coefficients, the companion matrix and the existence of complex eigenvalues. \Cref{sec:eigenspaces} studies eigenspaces and the two multiplicities, and \cref{sec:diagonalisation} characterises the matrices that can be diagonalised. \Cref{sec:matrix-powers} applies the results to powers of matrices and to the spectral radius, \cref{sec:cayley-hamilton} proves the Cayley--Hamilton theorem and introduces the minimal polynomial, and \cref{sec:complex-eigenvalues} treats real matrices with complex eigenvalues, unitary and Hermitian matrices. \Cref{sec:structured-eigen} computes the eigenvalues of structured matrices, \cref{sec:markov} treats Markov matrices and Gershgorin's theorem, and \cref{sec:computing-eigenvalues} looks at computation.

\section{Eigenvalues and Eigenvectors}
\label{sec:eigen-definitions}

\subsection{Definitions}

An operator on a vector space maps the space into itself, so that a vector and its image can be compared. In general the image points in a different direction; the vectors whose direction does not change are the subject of this chapter.

\begin{definition}[Eigenvalue, eigenvector, eigenspace]
\label{def:eigenvalue}
Let $T$ be an operator on a vector space $V$ over $\mathbb F$. A number $\lambda\in\mathbb F$ is an \term{eigenvalue} of $T$ if there is a vector $\mathbf v\ne\mathbf 0$ with
\[
T\mathbf v=\lambda\mathbf v ,
\]
and every such $\mathbf v$ is an \term{eigenvector} of $T$ for the eigenvalue $\lambda$. The subspace
\[
E_\lambda(T)\coloneqq\ker(T-\lambda\operatorname{id})=\{\mathbf v\in V: T\mathbf v=\lambda\mathbf v\}
\]
is the \term{eigenspace} of $\lambda$, and the set of all eigenvalues is the \term{spectrum} $\operatorname{spec}T$. The eigenvalues, eigenvectors and eigenspaces of a matrix $A\in\mathbb F^{n\times n}$ are those of the operator $\mathbf x\mapsto A\mathbf x$ on $\mathbb F^n$; in particular $E_\lambda(A)=\mathcal N(A-\lambda I)$.
\end{definition}

For matrices this is the definition of \cref{sec:extrema-n}. The eigenvector must be nonzero, since $T\mathbf 0=\lambda\mathbf 0$ holds for every $\lambda$; the eigenvalue may be zero. The eigenspace $E_\lambda$ is defined for every $\lambda\in\mathbb F$, and $\lambda$ is an eigenvalue exactly when $E_\lambda\ne\{\mathbf 0\}$; its nonzero vectors are the eigenvectors for $\lambda$. We write $E_\lambda$ when the operator is clear from the context. The words come from German, where \emph{eigen} means \emph{own} or \emph{characteristic}: an eigenvector is a vector that has its own direction under $T$, and older English texts speak of characteristic values and proper values. The first properties follow directly from the definition.

\begin{proposition}[First properties]
\label{prop:eigen-basic}
Let $T$ be an operator on $V$ and $\lambda\in\mathbb F$.
\begin{enumerate}
\item $\lambda$ is an eigenvalue of $T$ if and only if $T-\lambda\operatorname{id}$ is not injective. If $V$ is finite-dimensional, this holds if and only if $T-\lambda\operatorname{id}$ is not invertible, and for a matrix $A\in\mathbb F^{n\times n}$ if and only if $\det(A-\lambda I)=0$.
\item $E_\lambda(T)$ is a subspace of $V$, and $T$ maps it into itself. A line $\operatorname{span}\{\mathbf v\}$, $\mathbf v\ne\mathbf 0$, is invariant under $T$ if and only if $\mathbf v$ is an eigenvector.
\item $0$ is an eigenvalue of $T$ if and only if $\ker T\ne\{\mathbf 0\}$; then $E_0(T)=\ker T$.
\item If $T\mathbf v=\lambda\mathbf v$, then $T^k\mathbf v=\lambda^k\mathbf v$ for every $k\in\mathbb N$ and $p(T)\mathbf v=p(\lambda)\mathbf v$ for every polynomial $p$. If $T$ is invertible, then $\lambda\ne0$ and $T^{-1}\mathbf v=\lambda^{-1}\mathbf v$.
\end{enumerate}
\end{proposition}

\begin{proof}
(i) $T\mathbf v=\lambda\mathbf v$ says $(T-\lambda\operatorname{id})\mathbf v=\mathbf 0$, so an eigenvector is a nonzero vector of $\ker(T-\lambda\operatorname{id})$, which exists exactly when this map is not injective (\cref{prop:kernel-image}). On a finite-dimensional space an operator is injective if and only if it is invertible (\cref{cor:injective-surjective}), and a square matrix is invertible if and only if its determinant is nonzero (\cref{thm:det-invertible}). (ii) A kernel is a subspace, and if $T\mathbf v=\lambda\mathbf v$, then $T(T\mathbf v)=\lambda T\mathbf v$, so $T\mathbf v\in E_\lambda$. The line spanned by $\mathbf v$ is invariant exactly when $T\mathbf v=\lambda\mathbf v$ for some $\lambda\in\mathbb F$. (iii) is the case $\lambda=0$ of the definition. (iv) By induction, $T^{k+1}\mathbf v=T(\lambda^k\mathbf v)=\lambda^kT\mathbf v=\lambda^{k+1}\mathbf v$, and for $p(x)=\sum_kc_kx^k$, $p(T)\mathbf v=\sum_kc_kT^k\mathbf v=\sum_kc_k\lambda^k\mathbf v=p(\lambda)\mathbf v$. If $T$ is invertible, $\lambda\ne0$ by (iii), and applying $T^{-1}$ to $T\mathbf v=\lambda\mathbf v$ gives $\mathbf v=\lambda T^{-1}\mathbf v$.
\end{proof}

By (i), the eigenvalues of a matrix are the numbers $\lambda$ that make $A-\lambda I$ singular, which is how \cref{ch:linear-maps} anticipated them: in \cref{prop:similarity-invariants}\,(iv) they are the numbers with $\dim\mathcal N(A-\lambda I)>0$. Criterion (i) does not hold in infinite dimensions in its second form: the operator that multiplies a polynomial by $x$ on the space of all polynomials is injective, so $0$ is not an eigenvalue, but it is not invertible, since the constant $1$ is not in its image. In fact it has no eigenvalues at all, since $xp(x)=\lambda p(x)$ forces $p=0$ by comparing the highest powers.

\subsection{Geometric Examples}

The maps of the plane in \cref{fig:linear-maps} illustrate the possibilities, and \cref{tab:plane-eigen} collects them.

The orthogonal projection $P_\alpha=\mathbf u\mathbf u^{\mathsf T}$ onto the line spanned by the unit vector $\mathbf u=(\cos\alpha,\sin\alpha)$ keeps every vector of the line and maps every vector of the perpendicular line to $\mathbf 0$. So $P_\alpha\mathbf u=\mathbf u$ and $P_\alpha\mathbf u^\perp=\mathbf 0$ for $\mathbf u^\perp=(-\sin\alpha,\cos\alpha)$: the eigenvalues are $1$ and $0$, with the eigenspaces $\operatorname{span}\{\mathbf u\}$ and $\operatorname{span}\{\mathbf u^\perp\}$. Every projection $P$, orthogonal or oblique, behaves in this way, by \cref{prop:projection}: its image is $E_1(P)$, its kernel is $E_0(P)$, and $V=E_1\oplus E_0$. The reflection $S_\alpha=2P_\alpha-I$ in the same line has the eigenvalues $1$ and $-1$, with the same two eigenspaces. A scaling $\operatorname{diag}(s_1,s_2)$ with $s_1\ne s_2$ has the eigenvalues $s_1$ and $s_2$, with the coordinate axes as eigenspaces, and a multiple $sI$ of the identity has the single eigenvalue $s$, for which every nonzero vector is an eigenvector.

The shear $K$ with the rows $(1,k)$ and $(0,1)$, $k\ne0$, maps $(x_1,x_2)$ to $(x_1+kx_2,x_2)$. If this is $\lambda(x_1,x_2)$, the second component gives $x_2=\lambda x_2$ and the first $x_1+kx_2=\lambda x_1$. For $x_2\ne0$ the first equation forces $\lambda=1$, and then the second gives $kx_2=0$, which is impossible; so $x_2=0$, $x_1\ne0$ and $\lambda=1$. The shear has the single eigenvalue $1$, and its eigenspace is the $x_1$-axis only: unlike the identity, which also has the single eigenvalue $1$, it has no basis of eigenvectors. \Cref{sec:eigenspaces} gives such matrices a name.

The rotation $R_\varphi$ with $\varphi\notin\pi\mathbb Z$ has no real eigenvalue, since it turns every nonzero vector of the plane through the angle $\varphi$, and a turned vector is not a real multiple of the original; \cref{sec:similarity} proved this by computation. Regarded as a complex matrix it has the eigenvalues $e^{i\varphi}$ and $e^{-i\varphi}$, with the eigenvectors $(1,-i)$ and $(1,i)$ found there. For $\varphi=\pi$ the rotation is $-I$, with the single eigenvalue $-1$.

\begin{table}
\sdfloat{%
\begin{sdtable}{colspec={lll}, row{1}={font=\small\bfseries}}
\toprule
map & eigenvalues & eigenspaces \\
\midrule
projection $P_\alpha=\mathbf u\mathbf u^{\mathsf T}$ & $1$, $0$ & $\operatorname{span}\{\mathbf u\}$, $\operatorname{span}\{\mathbf u^\perp\}$ \\
reflection $S_\alpha=2P_\alpha-I$ & $1$, $-1$ & $\operatorname{span}\{\mathbf u\}$, $\operatorname{span}\{\mathbf u^\perp\}$ \\
scaling $\operatorname{diag}(s_1,s_2)$, $s_1\ne s_2$ & $s_1$, $s_2$ & the two axes \\
multiple $sI$ of the identity & $s$ & the whole plane \\
shear $K$, $k\ne0$ & $1$ & the $x_1$-axis \\
rotation $R_\varphi$, $\varphi\notin\pi\mathbb Z$ & none real; $e^{\pm i\varphi}$ in $\mathbb C$ & spanned by $(1,\mp i)$ in $\mathbb C^2$ \\
\bottomrule
\end{sdtable}}{Eigenvalues and eigenspaces of the linear maps of the plane in \cref{fig:linear-maps}, with $\mathbf u=(\cos\alpha,\sin\alpha)$ and $\mathbf u^\perp=(-\sin\alpha,\cos\alpha)$.\label{tab:plane-eigen}}
\end{table}

For a general matrix the eigenvectors are not visible at a glance, but a picture still shows what they are. \Cref{fig:unit-circle-image} shows the images of the unit vectors under two matrices,
\begin{equation}
A=\begin{bmatrix}2&1\\1&2\end{bmatrix},\qquad B=\begin{bmatrix}1&1\\0&2\end{bmatrix}.
\label{eq:two-by-two-examples}
\end{equation}
The image of the unit circle is an ellipse in both cases. A unit vector $\mathbf x$ and its image $A\mathbf x$ usually point in different directions, but along two lines the direction is kept. For $A$ these are the lines spanned by $(1,1)$ and $(1,-1)$, on which $A$ multiplies by $3$ and by $1$, since $A(1,1)=(3,3)$ and $A(1,-1)=(1,-1)$; they are the axes of the ellipse, at right angles to each other. For $B$ they are the $x_1$-axis, with $B\mathbf e_1=\mathbf e_1$, and the line spanned by $(1,1)$, with $B(1,1)=(2,2)$; they are not perpendicular, and they are not the axes of the ellipse. That the eigenvectors of the symmetric matrix $A$ are orthogonal is no accident (\cref{prop:hermitian-eigen}), and the axes of the ellipse are the subject of \cref{ch:svd}.

\begin{figure}
\newcommand\sdcirclepanel[5]{\begin{tikzpicture}[scale=0.78, font=\scriptsize]
\draw[black!40, line width=0.5pt] (0,0) circle (1);
\draw[line width=0.7pt, smooth, samples=120, domain=0:360] plot ({#1*cos(\x)+#2*sin(\x)}, {#3*cos(\x)+#4*sin(\x)});
\foreach \t in {0,30,...,330} {\draw[black!45, line width=0.3pt] ({cos(\t)},{sin(\t)}) -- ({#1*cos(\t)+#2*sin(\t)}, {#3*cos(\t)+#4*sin(\t)});
  \fill[black!45] ({cos(\t)},{sin(\t)}) circle (0.9pt);}
#5
\draw[line width=0.4pt, -{Latex[length=1.5mm]}] (-3.3,0) -- (3.5,0) node[below] {$x_1$};
\draw[line width=0.4pt, -{Latex[length=1.5mm]}] (0,-3.0) -- (0,3.2) node[left] {$x_2$};
\end{tikzpicture}}
\sdpanel{A}{\sdcirclepanel{2}{1}{1}{2}{%
\draw[dashed, line width=0.5pt] (-2.9,-2.9) -- (2.9,2.9);
\draw[dashed, line width=0.5pt] (-2.2,2.2) -- (2.2,-2.2);
\draw[line width=1pt, -{Latex[length=2mm]}] (0,0) -- (2.1213,2.1213) node[right, fill=white, inner sep=1pt] {$3\mathbf x_1$};
\draw[line width=1pt, -{Latex[length=2mm]}] (0,0) -- (0.7071,-0.7071) node[below left, fill=white, inner sep=1pt] {$\mathbf x_2$};}}\hfill
\sdpanel{B}{\sdcirclepanel{1}{1}{0}{2}{%
\draw[dashed, line width=0.5pt] (-2.4,-2.4) -- (2.4,2.4);
\draw[line width=1pt, -{Latex[length=2mm]}] (0,0) -- (1.4142,1.4142);
\node[fill=white, inner sep=1pt] at (2.05,1.2) {$2\mathbf x_2$};
\draw[line width=1pt, -{Latex[length=2mm]}] (0,0) -- (1,0);
\node[fill=white, inner sep=1pt] at (0.5,-0.28) {$\mathbf x_1$};}}
\caption{The unit circle (grey), the images of twelve unit vectors (each joined to its image by a thin line) and the image of the circle (black ellipse). (A) Under the symmetric matrix $A$ of \eqref{eq:two-by-two-examples}, the unit vectors $\mathbf x_1=(1,1)/\sqrt2$ and $\mathbf x_2=(1,-1)/\sqrt2$ are mapped to $3\mathbf x_1$ and $\mathbf x_2$; the two eigenvector lines (dashed) are perpendicular and are the axes of the ellipse. (B) Under the matrix $B$ of \eqref{eq:two-by-two-examples}, $\mathbf x_1=\mathbf e_1$ is fixed and $\mathbf x_2=(1,1)/\sqrt2$ is doubled; the eigenvector lines (the $x_1$-axis and the dashed line) are not perpendicular, and they are not the axes of the ellipse.\label{fig:unit-circle-image}}
\end{figure}

\subsection{Operators on Spaces of Functions}

Eigenvalues are not confined to matrices. The differentiation operator $D$ maps $C^\infty(\mathbb R)$, the smooth real functions, into itself, and $Df=\lambda f$ is the differential equation $f'=\lambda f$. By \cref{thm:exp-ode}, its solutions are exactly the functions $ce^{\lambda x}$. So every real number $\lambda$ is an eigenvalue of $D$, and its eigenspace is the line spanned by $e^{\lambda x}$; on the complex-valued smooth functions, every complex $\lambda$ is an eigenvalue in the same way, with the eigenfunction $e^{\lambda x}$ of \cref{sec:complex-exp}. Eigenvectors in spaces of functions are called \term{eigenfunctions}. This is the reason why exponential functions solve linear differential equations with constant coefficients, the subject of Chapter~24: such an equation is $p(D)f=0$ for a polynomial $p$, and $p(D)e^{\lambda x}=p(\lambda)e^{\lambda x}$ by \cref{prop:eigen-basic}\,(iv), which vanishes when $\lambda$ is a root of $p$.

On the finite-dimensional space $\mathcal P_n$ of polynomials of degree at most $n$, differentiation behaves quite differently. If $p\ne0$ has degree $d$, then $p'$ has degree less than $d$, so $p'=\lambda p$ is possible only for $\lambda=0$ and constant $p$. The only eigenvalue is $0$, with the constants as eigenspace, in accordance with the Jordan block that \cref{sec:similarity} found as the matrix of $D$. The second derivative with a sign, $-D^2$, on the smooth functions on $[0,\pi]$ that vanish at both ends, maps $\sin kx$ to $k^2\sin kx$ for every integer $k\ge1$; these are the eigenfunctions of the self-adjoint operator of \cref{sec:adjoints}, and \cref{sec:structured-eigen} finds their discrete counterparts in the eigenvectors of $T_n$. On sequences, the \term{shift} $(x_0,x_1,x_2,\ldots)\mapsto(x_1,x_2,x_3,\ldots)$ has every number $r$ as an eigenvalue, with the geometric sequence $(1,r,r^2,\ldots)$ as eigenvector; this is why the solutions of linear recurrences in \cref{thm:linear-recurrence} are combinations of powers $r^n$.

\subsection{Similarity and Triangular Matrices}

The eigenvalues of an operator on a finite-dimensional space can be computed from any of its matrices, because similar matrices have the same eigenvalues.

\begin{proposition}[Eigenvalues are similarity invariants]
\label{prop:similar-eigen}
Let $B=P^{-1}AP$ with $A,B,P\in\mathbb F^{n\times n}$ and $P$ invertible. Then $A$ and $B$ have the same eigenvalues, $\dim E_\lambda(B)=\dim E_\lambda(A)$ for every $\lambda$, and $\mathbf x$ is an eigenvector of $B$ for $\lambda$ if and only if $P\mathbf x$ is an eigenvector of $A$ for $\lambda$; that is, $E_\lambda(A)=P\,E_\lambda(B)$. Consequently, if $T$ is an operator on a finite-dimensional space $V$ with a basis $\mathcal B$, then $T$ and $[T]_{\mathcal B}$ have the same eigenvalues, and $\mathbf v$ is an eigenvector of $T$ if and only if its coordinate vector $[\mathbf v]_{\mathcal B}$ is an eigenvector of $[T]_{\mathcal B}$.
\end{proposition}

\begin{proof}
$B\mathbf x=\lambda\mathbf x$ is equivalent to $AP\mathbf x=\lambda P\mathbf x$, by multiplying with $P$ or $P^{-1}$ from the left, and $\mathbf x\ne\mathbf 0$ if and only if $P\mathbf x\ne\mathbf 0$. So $P$ maps $E_\lambda(B)$ onto $E_\lambda(A)$, injectively, and the dimensions agree; this is also \cref{prop:similarity-invariants}\,(iv). For the operator, $[T\mathbf v]_{\mathcal B}=[T]_{\mathcal B}[\mathbf v]_{\mathcal B}$ (\cref{thm:matrix-of-map}), and the coordinate map is an isomorphism (\cref{thm:coordinate-map}), so $T\mathbf v=\lambda\mathbf v$ if and only if $[T]_{\mathcal B}[\mathbf v]_{\mathcal B}=\lambda[\mathbf v]_{\mathcal B}$.
\end{proof}

The simplest matrices whose eigenvalues can be read off are the triangular ones.

\begin{proposition}[Eigenvalues of triangular matrices]
\label{prop:triangular-eigen}
The eigenvalues of a triangular matrix are its diagonal entries.
\end{proposition}

\begin{proof}
If $U$ is upper triangular with the diagonal entries $u_{ii}$, then $U-\lambda I$ is upper triangular with the diagonal entries $u_{ii}-\lambda$, and by \cref{thm:triangular}\,(ii) it is singular if and only if one of them vanishes, that is, if and only if $\lambda=u_{ii}$ for some $i$. The lower triangular case is the same.
\end{proof}

The eigenvalues of the matrix $B$ of \eqref{eq:two-by-two-examples} are therefore $1$ and $2$, and those of the shear $K$ and of every Jordan block are the single diagonal entry. The proposition fails for matrices that are triangular only in blocks or nearly triangular: the eigenvalues of a matrix are not its diagonal entries in general, as $A$ in \eqref{eq:two-by-two-examples}, with diagonal entries $2$ and eigenvalues $3$ and $1$, shows. They do lie near the diagonal entries when the entries off the diagonal are small, as \cref{sec:markov} makes precise.

\section{The Characteristic Polynomial}
\label{sec:characteristic-polynomial}

\subsection{Definition and Coefficients}

By \cref{prop:eigen-basic}\,(i), the eigenvalues of $A\in\mathbb F^{n\times n}$ are the numbers $\lambda$ with $\det(A-\lambda I)=0$. The entries of $\lambda I-A$ are $\lambda\delta_{ij}-a_{ij}$, and each term of the Leibniz formula \eqref{eq:leibniz} is a product of $n$ such entries, one from each column; expanding the products with $\lambda$ treated as a variable turns $\det(\lambda I-A)$ into a polynomial in $\lambda$, whose value at every number $\lambda\in\mathbb F$ is the determinant of the matrix $\lambda I-A$. We use $\lambda I-A$ rather than $A-\lambda I$ because its determinant has the leading coefficient $1$; the two determinants differ by the factor $(-1)^n$.

\begin{definition}[Characteristic polynomial]
\label{def:characteristic-polynomial}
The \term{characteristic polynomial} of $A\in\mathbb F^{n\times n}$ is $p_A(\lambda)\coloneqq\det(\lambda I-A)$, and $p_A(\lambda)=0$ is its \term{characteristic equation}.
\end{definition}

Since $\mathbb F$ is infinite, a polynomial is determined by its values: two polynomials that agree at infinitely many points are equal, by \cref{cor:factor}\,(ii) applied to their difference. Identities between characteristic polynomials may therefore be proved by evaluating both sides at every number $\lambda$, or at every $\lambda$ except finitely many, and we shall use this repeatedly. The coefficients of $p_A$ are sums of minors. For a set $S\subseteq\{1,\ldots,n\}$ we write $A_S$ for the \term{principal submatrix} of $A$ formed by the rows and the columns with indices in $S$, and its determinant is a \term{principal minor} of order $|S|$.

\begin{theorem}[Coefficients of the characteristic polynomial]
\label{thm:charpoly-coefficients}
For $A\in\mathbb F^{n\times n}$, $p_A$ is a polynomial of degree $n$ with the leading coefficient $1$,
\[
p_A(\lambda)=\lambda^n+c_{n-1}\lambda^{n-1}+\cdots+c_1\lambda+c_0,\qquad c_{n-k}=(-1)^k\sum_{|S|=k}\det A_S\quad(k=1,\ldots,n),
\]
where the sum runs over all subsets $S$ of $\{1,\ldots,n\}$ with $k$ elements. In particular
\[
c_{n-1}=-\operatorname{tr}A,\qquad c_0=(-1)^n\det A .
\]
The eigenvalues of $A$ in $\mathbb F$ are the roots of $p_A$ in $\mathbb F$.
\end{theorem}

\begin{proof}
The $j$th column of $\lambda I-A$ is $\lambda\mathbf e_j-\mathbf a_j$, where $\mathbf a_j$ is the $j$th column of $A$. The determinant is linear in each column (\cref{thm:det-existence}), so splitting every column into its two parts expresses $\det(\lambda I-A)$ as a sum of $2^n$ determinants, one for each choice of the part taken from each column. Let $S$ be the set of the columns from which $-\mathbf a_j$ is taken. Taking the factor $\lambda$ out of each of the other $n-|S|$ columns, the term for $S$ is $\lambda^{n-|S|}\det N_S$, where $N_S$ has the column $-\mathbf a_j$ for $j\in S$ and the column $\mathbf e_j$ for $j\notin S$. Expanding $\det N_S$ along a column $\mathbf e_j$ with $j\notin S$ (\cref{thm:laplace-expansion}) leaves the single term with the sign $(-1)^{j+j}=1$ and the minor obtained by deleting row $j$ and column $j$; in the smaller matrix the remaining columns $\mathbf e_{j'}$ still have their entry $1$ on the diagonal. Deleting all rows and columns with indices outside $S$ in this way leaves the matrix $-A_S$, so $\det N_S=\det(-A_S)=(-1)^{|S|}\det A_S$, with $\det N_\varnothing=\det I=1$. Collecting the terms with $|S|=k$ gives the coefficient of $\lambda^{n-k}$. For $k=1$ the principal minors are the diagonal entries $a_{jj}$, whose sum is the trace, and for $k=n$ the only set is $S=\{1,\ldots,n\}$. Finally $\det(A-\lambda I)=(-1)^np_A(\lambda)$, so the eigenvalues are the roots of $p_A$ by \cref{prop:eigen-basic}\,(i).
\end{proof}

For a matrix $A$ of order $2$ with the rows $(a,b)$ and $(c,d)$,
\[
p_A(\lambda)=\lambda^2-(a+d)\lambda+(ad-bc)=\lambda^2-(\operatorname{tr}A)\lambda+\det A ,
\]
and the quadratic formula gives the eigenvalues
\begin{equation}
\lambda_{1,2}=\frac{\operatorname{tr}A\pm\sqrt{(\operatorname{tr}A)^2-4\det A}}{2},\qquad(\operatorname{tr}A)^2-4\det A=(a-d)^2+4bc .
\label{eq:eigen-2x2}
\end{equation}
For a real matrix there are three cases: if the discriminant $(a-d)^2+4bc$ is positive, there are two real eigenvalues; if it is zero, a single real one; and if it is negative, two complex conjugate ones, computed with the complex square root of \cref{sec:complex}. A real symmetric matrix, with $b=c$, always has real eigenvalues, since then the discriminant is $(a-d)^2+4b^2\ge0$; \cref{prop:hermitian-eigen} extends this to all orders. The examples met so far fit the formula. The matrix $A$ of \eqref{eq:two-by-two-examples} has $p_A(\lambda)=\lambda^2-4\lambda+3=(\lambda-1)(\lambda-3)$. The Fibonacci matrix $F$ with the rows $(1,1)$ and $(1,0)$ has the trace $1$ and the determinant $-1$, so
\[
p_F(\lambda)=\lambda^2-\lambda-1,
\]
the characteristic polynomial of the Fibonacci recurrence in \cref{thm:linear-recurrence}, with the roots $\phi=(1+\sqrt5)/2$ and $\hat\phi=(1-\sqrt5)/2$. The rotation $R_\varphi$ has the trace $2\cos\varphi$ and the determinant $1$, so $p(\lambda)=\lambda^2-2\lambda\cos\varphi+1$, whose roots are $\cos\varphi\pm\sqrt{\cos^2\varphi-1}=\cos\varphi\pm i\sin\varphi=e^{\pm i\varphi}$. The shear $K$ has $p_K(\lambda)=(\lambda-1)^2$.

For a matrix of order $3$ the theorem gives $p_A(\lambda)=\lambda^3-(\operatorname{tr}A)\lambda^2+E_2\lambda-\det A$, where $E_2$ is the sum of the three principal minors of order $2$. For the second difference matrix $T_3$ of \cref{sec:special-matrices}, the trace is $6$, the principal minors of order $2$ are $3$, $4$ and $3$, and the determinant is $4$, so
\[
p_{T_3}(\lambda)=\lambda^3-6\lambda^2+10\lambda-4=(\lambda-2)(\lambda^2-4\lambda+2),
\]
with the eigenvalues $2-\sqrt2$, $2$ and $2+\sqrt2$. For matrices of higher order the coefficients are no longer a practical way to compute the polynomial, let alone the eigenvalues, and \cref{sec:computing-eigenvalues} explains why the polynomial is avoided altogether in computation.

\subsection{Similarity Invariance}

\begin{proposition}[Invariance of the characteristic polynomial]
\label{prop:charpoly-similar}
Similar matrices have the same characteristic polynomial. Moreover $p_{A^{\mathsf T}}=p_A$, and the coefficients of $p_{A^{\mathsf H}}$ are the complex conjugates of those of $p_A$, so that the eigenvalues of $A^{\mathsf H}$ are the complex conjugates of those of $A$.
\end{proposition}

\begin{proof}
If $B=P^{-1}AP$, then $\lambda I-B=P^{-1}(\lambda I-A)P$ for every $\lambda$, and by the product rule (\cref{thm:det-product}) $\det(\lambda I-B)=\det P^{-1}\det(\lambda I-A)\det P=\det(\lambda I-A)$. The polynomials agree at every point and are therefore equal. Next, $\det(\lambda I-A^{\mathsf T})=\det\bigl((\lambda I-A)^{\mathsf T}\bigr)=\det(\lambda I-A)$ by \cref{thm:det-transpose}. Finally, conjugating every entry conjugates every term of the Leibniz formula, so $\det\bar M=\overline{\det M}$ for every complex matrix $M$, and $\det(\lambda I-A^{\mathsf H})=\overline{\det(\bar\lambda I-A)}=\overline{p_A(\bar\lambda)}$, which is the polynomial with the conjugate coefficients.
\end{proof}

By the proposition, the characteristic polynomial of an operator $T$ on a finite-dimensional space is well defined as that of any of its matrices $[T]_{\mathcal B}$, and so are all its coefficients, among them the trace and the determinant of \cref{prop:similarity-invariants,prop:det-operator}. The converse of the proposition is false: the identity $I_2$ and the shear $K$ both have the characteristic polynomial $(\lambda-1)^2$, and they are not similar, as \cref{sec:similarity} observed. The characteristic polynomial does not see the difference between them; \cref{sec:eigenspaces} introduces a number that does.

\subsection{Complex Eigenvalues Exist}

Over $\mathbb R$ a matrix need not have any eigenvalue, as the rotations show. Over $\mathbb C$ the fundamental theorem of algebra (\cref{thm:fta}) guarantees one, and with it a full set of $n$ eigenvalues counted with multiplicity.

\begin{theorem}[Existence of eigenvalues]
\label{thm:eigenvalues-exist}
\leavevmode
\begin{enumerate}
\item Every $A\in\mathbb C^{n\times n}$, $n\ge1$, has an eigenvalue. More precisely, there are numbers $\lambda_1,\ldots,\lambda_n\in\mathbb C$, unique up to their order, with
\[
p_A(\lambda)=(\lambda-\lambda_1)(\lambda-\lambda_2)\cdots(\lambda-\lambda_n),
\]
and the eigenvalues of $A$ are exactly $\lambda_1,\ldots,\lambda_n$. Every operator on a complex vector space of finite dimension $n\ge1$ has an eigenvector.
\item If $A\in\mathbb R^{n\times n}$ and $A\mathbf x=\lambda\mathbf x$ with $\mathbf x\in\mathbb C^n$, then $A\bar{\mathbf x}=\bar\lambda\bar{\mathbf x}$. The nonreal eigenvalues of a real matrix come in pairs $\lambda,\bar\lambda$, and $\lambda$ and $\bar\lambda$ occur equally often among $\lambda_1,\ldots,\lambda_n$.
\item A real eigenvalue of a real matrix has a real eigenvector, and a real matrix of odd order has a real eigenvalue.
\end{enumerate}
\end{theorem}

\begin{proof}
(i) By \cref{thm:charpoly-coefficients}, $p_A$ has degree $n\ge1$ and the leading coefficient $1$, so it factors as stated by \cref{cor:factor}\,(ii). The factorisation is unique up to order, by induction on $n$: if $\prod_i(\lambda-\lambda_i)=\prod_i(\lambda-\mu_i)$, then $\lambda_1$ is a root of the right side, hence equal to some $\mu_j$, and after renumbering $\mu_1=\lambda_1$; the two products of the remaining $n-1$ factors agree for every $\lambda\ne\lambda_1$, hence everywhere, and the induction hypothesis applies to them. The roots of $p_A$ are the $\lambda_i$, and they are the eigenvalues by \cref{thm:charpoly-coefficients}. For an operator, apply this to its matrix in a basis and use \cref{prop:similar-eigen}. (ii) Conjugating $A\mathbf x=\lambda\mathbf x$ entry by entry gives $\bar A\bar{\mathbf x}=\bar\lambda\bar{\mathbf x}$, and $\bar A=A$. The coefficients of $p_A$ are real, so conjugating the factorisation of (i) gives $p_A(\lambda)=\overline{p_A(\bar\lambda)}=\prod_i(\lambda-\bar\lambda_i)$; by uniqueness, the list $\bar\lambda_1,\ldots,\bar\lambda_n$ is a rearrangement of $\lambda_1,\ldots,\lambda_n$, so each number occurs in it as often as its conjugate. (iii) If $\lambda$ is real, $A-\lambda I$ is a real singular matrix, and elimination in real arithmetic finds a nonzero real vector in its null space (\cref{thm:invertible}). By (ii), the nonreal $\lambda_i$ can be grouped into pairs, so the number of real ones has the same parity as $n$, and is at least $1$ when $n$ is odd.
\end{proof}

The numbers $\lambda_1,\ldots,\lambda_n$ are the \term{eigenvalues of $A$ counted with multiplicity}, and the number of indices $i$ with $\lambda_i=\lambda$ is the \term{algebraic multiplicity} $\operatorname{am}(\lambda)$ of the eigenvalue $\lambda$: the exponent of the factor $\lambda-\lambda_i$ in $p_A$. A real matrix is regarded as a complex one for this purpose. For the shear $K$ the eigenvalue $1$ has algebraic multiplicity $2$, and so has the eigenvalue $1$ of the identity $I_2$. The theorem turns the trace and the determinant into functions of the eigenvalues, which explains why they are similarity invariants.

\begin{theorem}[Trace and determinant]
\label{thm:trace-det-eigen}
If $\lambda_1,\ldots,\lambda_n$ are the eigenvalues of $A\in\mathbb C^{n\times n}$ counted with multiplicity, then
\[
\operatorname{tr}A=\lambda_1+\lambda_2+\cdots+\lambda_n,\qquad\det A=\lambda_1\lambda_2\cdots\lambda_n .
\]
\end{theorem}

\begin{proof}
Multiplying out $\prod_i(\lambda-\lambda_i)$, the coefficient of $\lambda^{n-1}$ is $-\sum_i\lambda_i$ and the constant term is $(-1)^n\prod_i\lambda_i$. Comparing with \cref{thm:charpoly-coefficients} gives $-\operatorname{tr}A=-\sum_i\lambda_i$ and $(-1)^n\det A=(-1)^n\prod_i\lambda_i$.
\end{proof}

In the same way, the sum of the principal minors of order $k$ is the sum of all products of $k$ of the eigenvalues with different indices. For the Fibonacci matrix, $\phi+\hat\phi=1$ and $\phi\hat\phi=-1$; for the rotation, $e^{i\varphi}+e^{-i\varphi}=2\cos\varphi$ and $e^{i\varphi}e^{-i\varphi}=1$; for $T_3$, the eigenvalues add up to $6$ and multiply to $2\cdot(4-2)=4$. For a projection $P$ of rank $r$ on $\mathbb F^n$, the eigenvalue $1$ occurs $r$ times and $0$ occurs $n-r$ times, as the block form of \cref{prop:projection}\,(iii) shows with \cref{prop:triangular-eigen}, and the trace $r$ is their sum. The theorem also explains case (ii) after Jacobi's formula (\cref{thm:jacobi-formula}). If $B\in\mathbb C^{n\times n}$ has the eigenvalues $\mu_1,\ldots,\mu_n$ and $t\ne0$, then
\[
p_{I+tB}(\lambda)=\det\bigl((\lambda-1)I-tB\bigr)=t^np_B\Bigl(\frac{\lambda-1}t\Bigr)=\prod_{i=1}^n\bigl(\lambda-1-t\mu_i\bigr),
\]
so $I+tB$ has the eigenvalues $1+t\mu_i$, and
\[
\det(I+tB)=\prod_{i=1}^n(1+t\mu_i)=1+t\sum_{i=1}^n\mu_i+O(t^2)=1+t\operatorname{tr}B+O(t^2)\qquad(t\to0):
\]
to first order the determinant adds up the small changes $t\mu_i$ of the eigenvalues, and their sum is the trace.

\subsection{Products in Either Order}

The products $AB$ and $BA$ of two matrices are in general different, and of different sizes if the matrices are rectangular. Their traces are equal (\cref{prop:trace}\,(ii)), and so are all their nonzero eigenvalues.

\begin{theorem}[Eigenvalues of \texorpdfstring{$AB$}{AB} and \texorpdfstring{$BA$}{BA}]
\label{thm:ab-ba}
Let $A\in\mathbb F^{m\times n}$ and $B\in\mathbb F^{n\times m}$ with $m\ge n$. Then
\[
p_{AB}(\lambda)=\lambda^{m-n}p_{BA}(\lambda).
\]
Consequently $AB$ and $BA$ have the same nonzero eigenvalues with the same algebraic multiplicities, and if $m=n$, then $p_{AB}=p_{BA}$.
\end{theorem}

\begin{proof}
Let $\lambda\ne0$. Then $\lambda I_m-AB=\lambda(I_m-\lambda^{-1}AB)$, and by Sylvester's determinant identity (\cref{thm:det-lemma}\,(i)), applied to the matrices $-\lambda^{-1}A$ and $B$,
\[
\det(\lambda I_m-AB)=\lambda^m\det(I_m-\lambda^{-1}AB)=\lambda^m\det(I_n-\lambda^{-1}BA)=\lambda^{m-n}\det(\lambda I_n-BA).
\]
The two sides are polynomials in $\lambda$ that agree for all $\lambda\ne0$, so they are equal. The factor $\lambda^{m-n}$ contributes only the root $0$.
\end{proof}

The eigenvectors correspond as well: if $AB\mathbf x=\lambda\mathbf x$ with $\lambda\ne0$ and $\mathbf x\ne\mathbf 0$, then $\mathbf y\coloneqq B\mathbf x$ is not zero, since $A\mathbf y=\lambda\mathbf x\ne\mathbf 0$, and $BA\mathbf y=B(AB\mathbf x)=\lambda\mathbf y$. For a column $\mathbf u$ and a row $\mathbf v^{\mathsf T}$, both with $n$ entries, the theorem with $A=\mathbf u$ and $B=\mathbf v^{\mathsf T}$ gives the characteristic polynomial of the rank-one matrix $\mathbf u\mathbf v^{\mathsf T}$:
\[
p_{\mathbf u\mathbf v^{\mathsf T}}(\lambda)=\lambda^{n-1}(\lambda-\mathbf v^{\mathsf T}\mathbf u),
\]
since $\mathbf v^{\mathsf T}\mathbf u$ is a $1\times1$ matrix. Its eigenvalues are $\mathbf v^{\mathsf T}\mathbf u$ and $0$, the latter with algebraic multiplicity at least $n-1$. Even for square matrices, $AB$ and $BA$ need not be similar: with $N$ the matrix with the rows $(0,1)$ and $(0,0)$ and $E_{11}=\operatorname{diag}(1,0)$, $NE_{11}=0$ and $E_{11}N=N$, two matrices with the same characteristic polynomial $\lambda^2$ that are not similar, since only $0$ is similar to $0$.

\subsection{The Companion Matrix}

Every matrix has a characteristic polynomial, and every polynomial with the leading coefficient $1$, a \term{monic} polynomial, arises in this way.

\begin{proposition}[Companion matrix]
\label{prop:companion}
Let $q(\lambda)=\lambda^n+c_{n-1}\lambda^{n-1}+\cdots+c_1\lambda+c_0$ with $c_k\in\mathbb F$ and $n\ge1$, and let
\[
C_q\coloneqq\begin{bmatrix}0&0&\cdots&0&-c_0\\1&0&\cdots&0&-c_1\\0&1&\cdots&0&-c_2\\\vdots&&\ddots&&\vdots\\0&0&\cdots&1&-c_{n-1}\end{bmatrix}\in\mathbb F^{n\times n},
\]
with ones below the diagonal, the negated coefficients in the last column and zeros elsewhere. Then $p_{C_q}=q$. If $r$ is a root of $q$, then $(1,r,r^2,\ldots,r^{n-1})$ is an eigenvector of $C_q^{\mathsf T}$ for the eigenvalue $r$.
\end{proposition}

\begin{proof}
By induction on $n$. For $n=1$, $C_q=[-c_0]$ and $\det(\lambda+c_0)=q(\lambda)$. For $n\ge2$, expand $\det(\lambda I-C_q)$ along its first row, which is $(\lambda,0,\ldots,0,c_0)$. The minor of the entry $\lambda$ is $\det(\lambda I-C_{q'})$ for the companion matrix of $q'(\lambda)=\lambda^{n-1}+c_{n-1}\lambda^{n-2}+\cdots+c_1$, which is $q'(\lambda)$ by the induction hypothesis. The minor of the entry $c_0$ in position $(1,n)$ is the determinant of the matrix formed by the rows $2,\ldots,n$ and the columns $1,\ldots,n-1$ of $\lambda I-C_q$; this matrix is upper triangular with $-1$ on the diagonal, so its determinant is $(-1)^{n-1}$. Hence
\[
p_{C_q}(\lambda)=\lambda q'(\lambda)+(-1)^{1+n}c_0(-1)^{n-1}=\lambda^n+c_{n-1}\lambda^{n-1}+\cdots+c_1\lambda+c_0=q(\lambda).
\]
For the eigenvector, the rows of $C_q^{\mathsf T}$ are $\mathbf e_2^{\mathsf T},\ldots,\mathbf e_n^{\mathsf T}$ and $(-c_0,\ldots,-c_{n-1})$, so $C_q^{\mathsf T}$ maps $\mathbf x=(1,r,\ldots,r^{n-1})$ to $(r,r^2,\ldots,r^{n-1},-\sum_kc_kr^k)$, and $-\sum_{k<n}c_kr^k=r^n$ because $q(r)=0$. So $C_q^{\mathsf T}\mathbf x=r\mathbf x$.
\end{proof}

The transpose $C_q^{\mathsf T}$ is the matrix of a linear recurrence. If $a_{k+n}=-c_{n-1}a_{k+n-1}-\cdots-c_0a_k$, the vectors $\mathbf s_k\coloneqq(a_k,a_{k+1},\ldots,a_{k+n-1})$ satisfy $\mathbf s_{k+1}=C_q^{\mathsf T}\mathbf s_k$, and the eigenvectors of the proposition are the solutions $a_k=r^k$ of \cref{thm:linear-recurrence}. The Fibonacci recurrence has $q(\lambda)=\lambda^2-\lambda-1$, and $C_q^{\mathsf T}$ is the matrix with the rows $(0,1)$ and $(1,1)$, which is the Fibonacci matrix $F$ with the order of the coordinates reversed. Chapter~27 develops linear recurrences in general from this point of view, together with the stability of their numerical solution.

The proposition has a sobering consequence. By \cref{thm:charpoly-coefficients} and \cref{prop:companion}, finding the eigenvalues of matrices and finding the roots of polynomials are the same problem: the eigenvalues of $A$ are the roots of $p_A$, and the roots of $q$ are the eigenvalues of $C_q$. For polynomials of degree $2$, $3$ and $4$ there are formulas for the roots in terms of the coefficients, built from the arithmetic operations and roots of numbers, but for degree $5$ and higher there is none, as the note below explains. It follows that no finite formula of this kind can compute the eigenvalues of every matrix of order $5$ or more, and every general method for eigenvalues is \emph{iterative}: it produces a sequence of approximations that converges to the eigenvalues and stops when they are accurate enough. \Cref{sec:computing-eigenvalues} meets the simplest such method, and also shows why the characteristic polynomial, although it defines the eigenvalues, is a bad intermediate step for computing them. In the other direction, computing the eigenvalues of the companion matrix is a reliable way to find the roots of a polynomial~\cite{edelman1995}.

\begin{note}{Solving equations by radicals}
The quadratic formula was known in antiquity, and Italian mathematicians of the sixteenth century, Scipione del Ferro, Tartaglia, Cardano and Ferrari, found formulas for the roots of cubic and quartic equations in terms of the coefficients, the four arithmetic operations and the extraction of roots~\cite{katz2009}. For the quintic, the search failed. Ruffini gave an incomplete proof in 1799 that no such formula exists, and Abel a complete one in 1826~\cite{abel1826quintic,stewart2015galois}: there is no expression in the coefficients of the general equation of degree $n\ge5$, built from the arithmetic operations and radicals, that yields its roots. Galois then characterised the particular equations that can be solved by radicals, through the group of permutations of the roots that his theory attaches to each equation; the equation $x^5-6x+3=0$, for instance, cannot be solved by radicals. Stewart's book~\cite{stewart2015galois} gives a readable account of the theory and its history. For computation the theorem means that the eigenvalues of a general matrix of order $5$ or more cannot be obtained by a finite sequence of arithmetic operations and root extractions; every eigenvalue algorithm is iterative, as the QR algorithm of Volume~VIII is.
\end{note}

\section{Eigenspaces and Multiplicity}
\label{sec:eigenspaces}

\subsection{Eigenvectors for Different Eigenvalues}

The eigenvectors that belong to one eigenvalue form a subspace. Eigenvectors that belong to different eigenvalues have nothing in common.

\begin{theorem}[Eigenvectors for distinct eigenvalues]
\label{thm:distinct-independent}
Let $T$ be an operator on $V$, and let $\mathbf x_1,\ldots,\mathbf x_k$ be eigenvectors of $T$ for pairwise different eigenvalues $\lambda_1,\ldots,\lambda_k$. Then $\mathbf x_1,\ldots,\mathbf x_k$ are linearly independent.
\end{theorem}

\begin{proof}
By induction on $k$. For $k=1$, $\mathbf x_1\ne\mathbf 0$ is independent. Let $k\ge2$ and $c_1\mathbf x_1+\cdots+c_k\mathbf x_k=\mathbf 0$. Applying $T-\lambda_k\operatorname{id}$, which maps $\mathbf x_i$ to $(\lambda_i-\lambda_k)\mathbf x_i$, gives
\[
c_1(\lambda_1-\lambda_k)\mathbf x_1+\cdots+c_{k-1}(\lambda_{k-1}-\lambda_k)\mathbf x_{k-1}=\mathbf 0 .
\]
By the induction hypothesis all coefficients vanish, and since $\lambda_i\ne\lambda_k$ for $i<k$, $c_1=\cdots=c_{k-1}=0$. Then $c_k\mathbf x_k=\mathbf 0$, and $c_k=0$.
\end{proof}

The proof shows the mechanism: the factor $T-\lambda_k\operatorname{id}$ annihilates one of the vectors and only rescales the others, so a dependence among them would survive with one vector fewer. The theorem extends from vectors to subspaces.

\begin{corollary}[Eigenspaces form a direct sum]
\label{cor:eigenspaces-direct}
Let $\lambda_1,\ldots,\lambda_k$ be pairwise different eigenvalues of an operator $T$. Then the sum $E_{\lambda_1}+\cdots+E_{\lambda_k}$ is direct. If $V$ has dimension $n$ and $\mathcal B_i$ is a basis of $E_{\lambda_i}$, the bases $\mathcal B_1,\ldots,\mathcal B_k$ together form a basis of $E_{\lambda_1}\oplus\cdots\oplus E_{\lambda_k}$, so that
\[
\dim E_{\lambda_1}+\cdots+\dim E_{\lambda_k}\le n ;
\]
in particular $T$ has at most $n$ different eigenvalues.
\end{corollary}

\begin{proof}
Let $\mathbf v_1+\cdots+\mathbf v_k=\mathbf 0$ with $\mathbf v_i\in E_{\lambda_i}$. If some $\mathbf v_i$ were nonzero, the nonzero ones would be eigenvectors for different eigenvalues with a linear dependence among them, contradicting \cref{thm:distinct-independent}. So all $\mathbf v_i=\mathbf 0$, and the sum is direct (\cref{def:sum}). The vectors of $\mathcal B_1,\ldots,\mathcal B_k$ span the sum. If a combination of them vanishes, group its terms by $i$: the group of $\mathcal B_i$ is a vector $\mathbf v_i\in E_{\lambda_i}$, and $\sum_i\mathbf v_i=\mathbf 0$, so every $\mathbf v_i=\mathbf 0$, and since $\mathcal B_i$ is independent, all coefficients of the group vanish. Hence the bases together form a basis of the sum, whose dimension is at most $n$ (\cref{prop:subspace-dimension}). Each eigenspace has dimension at least $1$, which gives the last statement.
\end{proof}

This is the direct sum of several subspaces promised in \cref{sec:subspaces}. The three lines in the plane spanned by $(1,0)$, $(0,1)$ and $(1,1)$, whose sum is not direct, cannot therefore be eigenspaces of a single operator for three different eigenvalues; indeed an operator on the plane has at most two eigenvalues.

\subsection{Algebraic and Geometric Multiplicity}

\Cref{sec:characteristic-polynomial} counted an eigenvalue by its multiplicity as a root of the characteristic polynomial. The dimension of its eigenspace is a second count.

\begin{definition}[Geometric multiplicity, defective matrix]
\label{def:multiplicities}
The \term{geometric multiplicity} of an eigenvalue $\lambda$ of $A\in\mathbb F^{n\times n}$ is $\operatorname{gm}(\lambda)\coloneqq\dim E_\lambda(A)=n-\operatorname{rank}(A-\lambda I)$. A matrix is \term{defective} if, regarded as a complex matrix, it has an eigenvalue with $\operatorname{gm}(\lambda)<\operatorname{am}(\lambda)$.
\end{definition}

The rank of $A-\lambda I$ is the number of pivots in elimination, which is computed in the same way over $\mathbb R$ and over $\mathbb C$. For a real matrix and a real eigenvalue, the real eigenspace in $\mathbb R^n$ and the complex eigenspace in $\mathbb C^n$ therefore have the same dimension, and the geometric multiplicity does not depend on the field.

\begin{theorem}[The two multiplicities]
\label{thm:gm-am}
Every eigenvalue $\lambda$ of $A\in\mathbb C^{n\times n}$ satisfies $1\le\operatorname{gm}(\lambda)\le\operatorname{am}(\lambda)$. Consequently $\sum_\lambda\operatorname{gm}(\lambda)\le\sum_\lambda\operatorname{am}(\lambda)=n$, the sums running over the different eigenvalues, with equality if and only if $\operatorname{gm}(\lambda)=\operatorname{am}(\lambda)$ for every eigenvalue.
\end{theorem}

\begin{proof}
The eigenspace of an eigenvalue is not zero, so $g\coloneqq\operatorname{gm}(\lambda)\ge1$. Let $\mathbf x_1,\ldots,\mathbf x_g$ be a basis of $E_\lambda$, extend it to a basis of $\mathbb C^n$ (\cref{thm:basis-extension}), and let $S$ be the invertible matrix with these vectors as columns. The first $g$ columns of $S^{-1}AS$ are the coordinates of $A\mathbf x_j=\lambda\mathbf x_j$, that is, $\lambda\mathbf e_j$, so
\[
S^{-1}AS=\begin{bmatrix}\lambda I_g&B\\0&C\end{bmatrix}
\]
for some matrices $B$ and $C$, as in \cref{prop:block-triangular}. By \cref{prop:charpoly-similar} and the determinant of a block triangular matrix (\cref{thm:det-block-triangular}), applied to $xI-S^{-1}AS$,
\[
p_A(x)=\det\bigl((x-\lambda)I_g\bigr)\,p_C(x)=(x-\lambda)^gp_C(x).
\]
Factoring $p_C$ into linear factors and using the uniqueness of the factorisation in \cref{thm:eigenvalues-exist}, the factor $x-\lambda$ occurs at least $g$ times in $p_A$, so $\operatorname{am}(\lambda)\ge g$. The algebraic multiplicities add up to the degree $n$ of $p_A$, and a sum of inequalities $\operatorname{gm}\le\operatorname{am}$ is an equality only if each of them is.
\end{proof}

The simplest matrices with $\operatorname{gm}<\operatorname{am}$ are the Jordan blocks of \cref{sec:similarity}. For $\lambda\in\mathbb F$ and $n\ge2$ let
\[
J_n(\lambda)\coloneqq\begin{bmatrix}\lambda&1&&\\&\lambda&\ddots&\\&&\ddots&1\\&&&\lambda\end{bmatrix}\in\mathbb F^{n\times n},
\]
with $\lambda$ on the diagonal, $1$ just above it and zeros elsewhere. It is triangular, so $p(x)=(x-\lambda)^n$ and $\operatorname{am}(\lambda)=n$; but $J_n(\lambda)-\lambda I$ has the rank $n-1$, since its nonzero columns $\mathbf e_1,\ldots,\mathbf e_{n-1}$ are independent, so $\operatorname{gm}(\lambda)=1$, and the eigenspace is spanned by $\mathbf e_1$. The shear $K$ is $J_2(1)$, and the matrix of differentiation on $\mathcal P_n$ is $J_{n+1}(0)$ in the basis $1,x,x^2/2!,\ldots$, with the single eigenvalue $0$ of algebraic multiplicity $n+1$ and geometric multiplicity $1$. The two matrices
\[
\begin{bmatrix}2&0&0\\0&2&0\\0&0&3\end{bmatrix},\qquad\begin{bmatrix}2&1&0\\0&2&0\\0&0&3\end{bmatrix}
\]
have the same characteristic polynomial $(x-2)^2(x-3)$, but the eigenvalue $2$ has the geometric multiplicity $2$ in the first and $1$ in the second. By \cref{prop:similar-eigen} they are not similar, and the geometric multiplicity is the number that separates $I_2$ from the shear $K$ at the end of \cref{sec:characteristic-polynomial}. In the language of \cref{prop:kernel-chain}, the eigenspace $E_\lambda=\ker(A-\lambda I)$ is the first of the increasing kernels $\ker(A-\lambda I)^k$; for $J_n(\lambda)$ they grow by one dimension at each step until they fill the space. Chapter~22 shows that for every matrix the last of these kernels has the dimension $\operatorname{am}(\lambda)$, and builds the Jordan form from them.

\section{Diagonalisation}
\label{sec:diagonalisation}

\subsection{When a Matrix Can Be Diagonalised}

\begin{definition}[Diagonalisable matrix and operator]
\label{def:diagonalisable}
A matrix $A\in\mathbb F^{n\times n}$ is \term{diagonalisable} over $\mathbb F$ if there are an invertible $S\in\mathbb F^{n\times n}$ and a diagonal $\Lambda\in\mathbb F^{n\times n}$ with
\[
A=S\Lambda S^{-1}.
\]
An operator $T$ on a finite-dimensional space $V$ is \term{diagonalisable} if $V$ has a basis $\mathcal B$ for which $[T]_{\mathcal B}$ is diagonal.
\end{definition}

The definition says that $A$ is similar to a diagonal matrix. In the language of operators, the matrix $[T]_{\mathcal B}=\operatorname{diag}(\lambda_1,\ldots,\lambda_n)$ is diagonal exactly when $T\mathbf b_j=\lambda_j\mathbf b_j$ for every vector of the basis, since the $j$th column of $[T]_{\mathcal B}$ holds the coordinates of $T\mathbf b_j$: \emph{a matrix is diagonal in a basis if and only if the basis consists of eigenvectors}, which was the observation of \cref{sec:similarity}. The next theorem turns it into criteria.

\begin{theorem}[Criteria for diagonalisability]
\label{thm:diagonalisable}
Let $A\in\mathbb F^{n\times n}$ have the pairwise different eigenvalues $\mu_1,\ldots,\mu_k$ in $\mathbb F$. The following are equivalent:
\begin{enumerate}
\item $A$ is diagonalisable over $\mathbb F$;
\item $\mathbb F^n$ has a basis consisting of eigenvectors of $A$;
\item $\operatorname{gm}(\mu_1)+\cdots+\operatorname{gm}(\mu_k)=n$;
\item $\mathbb F^n=E_{\mu_1}\oplus E_{\mu_2}\oplus\cdots\oplus E_{\mu_k}$.
\end{enumerate}
If $p_A$ is a product of linear factors over $\mathbb F$, which is always the case for $\mathbb F=\mathbb C$, they are also equivalent to
\begin{enumerate}[start=5]
\item $\operatorname{gm}(\mu)=\operatorname{am}(\mu)$ for every eigenvalue $\mu$.
\end{enumerate}
If $A=S\Lambda S^{-1}$ with $\Lambda=\operatorname{diag}(\lambda_1,\ldots,\lambda_n)$, then the $j$th column of $S$ is an eigenvector for $\lambda_j$, and each eigenvalue $\mu$ occurs $\operatorname{am}(\mu)$ times among $\lambda_1,\ldots,\lambda_n$.
\end{theorem}

\begin{proof}
(i)$\Leftrightarrow$(ii) Let $S$ have the columns $\mathbf s_1,\ldots,\mathbf s_n$. The equation $A=S\Lambda S^{-1}$ is equivalent to $AS=S\Lambda$, and comparing columns, to $A\mathbf s_j=\lambda_j\mathbf s_j$ for every $j$. The matrix $S$ is invertible if and only if its columns form a basis of $\mathbb F^n$ (\cref{thm:invertible}). So a factorisation $A=S\Lambda S^{-1}$ is the same thing as a basis of eigenvectors, written as the columns of $S$, together with their eigenvalues. (ii)$\Rightarrow$(iv) Each vector of the basis lies in one of the eigenspaces, so their sum contains a basis and is $\mathbb F^n$; it is direct by \cref{cor:eigenspaces-direct}. (iv)$\Rightarrow$(iii) By \cref{cor:eigenspaces-direct}, the dimension of the direct sum is the sum of the dimensions. (iii)$\Rightarrow$(ii) By the same corollary, bases of the eigenspaces together form an independent list of $\operatorname{gm}(\mu_1)+\cdots+\operatorname{gm}(\mu_k)=n$ eigenvectors, which is a basis of $\mathbb F^n$ (\cref{thm:basis-extension}). (iii)$\Leftrightarrow$(v) If $p_A$ factors into linear factors over $\mathbb F$, all its roots are eigenvalues in $\mathbb F$, the algebraic multiplicities add up to $n$, and the proof of \cref{thm:gm-am}, which works over every $\mathbb F$ once $p_C$ factors, gives $\operatorname{gm}\le\operatorname{am}$; so the sum in (iii) is $n$ if and only if every inequality is an equality. Finally, $p_A=p_\Lambda=\prod_j(x-\lambda_j)$ by \cref{prop:charpoly-similar}, which gives the multiplicities.
\end{proof}

For a real matrix the theorem can be applied over $\mathbb R$ or over $\mathbb C$. Over $\mathbb R$, the sum in (iii) can reach $n$ only if all eigenvalues are real, so a real matrix is diagonalisable over $\mathbb R$ if and only if all its eigenvalues are real and each has equal multiplicities. The rotation $R_\varphi$ with $\varphi\notin\pi\mathbb Z$ is diagonalisable over $\mathbb C$, with two different eigenvalues, but not over $\mathbb R$, where it has none: with the eigenvectors of \cref{sec:similarity} as columns,
\[
R_\varphi=S\begin{bmatrix}e^{i\varphi}&0\\0&e^{-i\varphi}\end{bmatrix}S^{-1},\qquad S=\begin{bmatrix}1&1\\-i&i\end{bmatrix},\qquad S^{-1}=\frac12\begin{bmatrix}1&i\\1&-i\end{bmatrix}.
\]
Whether a matrix can be diagonalised thus depends on the field, and complex numbers enter for the same reason as in the fundamental theorem of algebra: over $\mathbb R$ the characteristic polynomial need not have roots. \Cref{sec:complex-eigenvalues} shows what can be achieved with real matrices when the eigenvalues are complex. A defective matrix cannot be diagonalised over any field, by (v).

\begin{corollary}[Distinct eigenvalues]
\label{cor:distinct-diagonalisable}
A matrix $A\in\mathbb F^{n\times n}$ with $n$ pairwise different eigenvalues in $\mathbb F$ is diagonalisable over $\mathbb F$.
\end{corollary}

\begin{proof}
Each of the $n$ eigenvalues has $\operatorname{gm}\ge1$, and the sum of the $n$ geometric multiplicities is at most $n$ by \cref{cor:eigenspaces-direct}, so each is $1$ and the sum is $n$; apply \cref{thm:diagonalisable}\,(iii).
\end{proof}

The converse is false: the identity has a single eigenvalue and is diagonal. A repeated eigenvalue is no obstacle in itself; what matters is whether its eigenspace is large enough.

\subsection{A Complete Example}

For the matrix
\[
A=\begin{bmatrix}4&0&1\\2&3&2\\1&0&4\end{bmatrix},
\]
the second column of $\lambda I-A$ has the single nonzero entry $\lambda-3$, on the diagonal, and the expansion along this column (\cref{thm:laplace-expansion}) gives
\[
p_A(\lambda)=(\lambda-3)\det\begin{bmatrix}\lambda-4&-1\\-1&\lambda-4\end{bmatrix}=(\lambda-3)\bigl((\lambda-4)^2-1\bigr)=(\lambda-3)^2(\lambda-5).
\]
The eigenvalue $3$ has algebraic multiplicity $2$, and $5$ has $1$. The eigenspace of $3$ is the null space of
\[
A-3I=\begin{bmatrix}1&0&1\\2&0&2\\1&0&1\end{bmatrix},
\]
whose rows are multiples of $(1,0,1)$; its rank is $1$, so $\operatorname{gm}(3)=2$, and the equation $x_1+x_3=0$, with $x_2$ free, gives the basis $(0,1,0)$, $(1,0,-1)$. For the eigenvalue $5$, the first and third rows of $A-5I$, $(-1,0,1)$ and $(1,0,-1)$, give $x_3=x_1$, and the second, $(2,-2,2)$, then gives $x_2=2x_1$, so $E_5$ is spanned by $(1,2,1)$. The geometric multiplicities add up to $3$, and with these eigenvectors as columns,
\[
S=\begin{bmatrix}0&1&1\\1&0&2\\0&-1&1\end{bmatrix},\qquad S^{-1}=\frac12\begin{bmatrix}-2&2&-2\\1&0&-1\\1&0&1\end{bmatrix},\qquad S^{-1}AS=\begin{bmatrix}3&0&0\\0&3&0\\0&0&5\end{bmatrix},
\]
as one confirms by multiplying out. A small change of one entry destroys the diagonalisability. The matrix $A'$ that agrees with $A$ except for the entry $a'_{23}=1$ has the same characteristic polynomial, by the same expansion, but
\[
A'-3I=\begin{bmatrix}1&0&1\\2&0&1\\1&0&1\end{bmatrix}
\]
has the rank $2$, so $\operatorname{gm}(3)=1<2=\operatorname{am}(3)$, and $A'$ is defective. Its eigenvectors are the multiples of $(0,1,0)$ and of $(2,3,2)$, and they span only a plane. The characteristic polynomial cannot tell the two matrices apart; the ranks of $A-\mu I$ can.

\begin{note}{Diagonalisable matrices are dense}
Defective matrices are exceptional. If the eigenvalues of $A\in\mathbb C^{n\times n}$ are not all different, an arbitrarily small change of the entries makes them different, and the changed matrix is diagonalisable by \cref{cor:distinct-diagonalisable}. The proof uses the Schur form of Chapter~22: $A=UTU^{\mathsf H}$ with a unitary $U$ and an upper triangular $T$, whose diagonal entries are the eigenvalues (\cref{prop:triangular-eigen}); adding small, suitably chosen numbers to the diagonal of $T$ makes them distinct and changes $A$ by as little as one wishes~\cite[Sect.~2.4]{horn2013}. In the other direction, a matrix has a repeated eigenvalue exactly when the discriminant of its characteristic polynomial vanishes, and the discriminant is a polynomial in the entries of the matrix; for $n=2$ it is $(a-d)^2+4bc$, as in \eqref{eq:eigen-2x2}. The matrices with a repeated eigenvalue thus form the zero set of a single nonzero polynomial, a set without interior points, and the defective matrices are a part of it. The same argument applies to real matrices and real changes: the discriminant is a nonzero polynomial in the real entries, so the real matrices with $n$ different complex eigenvalues, which are diagonalisable over $\mathbb C$, are dense in $\mathbb R^{n\times n}$. Diagonalisability over $\mathbb R$ is a different matter. A real $2\times2$ matrix with $(a-d)^2+4bc<0$ has nonreal eigenvalues, and so has every real matrix whose entries differ from it by little enough, since the discriminant depends continuously on the entries; such matrices fill a nonempty open set, so the matrices that can be diagonalised over $\mathbb R$ are not dense in $\mathbb R^{2\times2}$. Near a defective matrix, eigenvalues are extremely sensitive to perturbations, as \cref{sec:computing-eigenvalues} shows; this, and not their rarity, is why defective matrices matter in computation.
\end{note}

\subsection{Spectral Projections and Commuting Matrices}

A diagonalisable matrix splits $\mathbb F^n$ into the direct sum of its eigenspaces, and the components of a vector in this sum can be computed by polynomials in the matrix. For pairwise different numbers $\mu_1,\ldots,\mu_k$, let
\begin{equation}
\ell_j(x)\coloneqq\prod_{i\ne j}\frac{x-\mu_i}{\mu_j-\mu_i}\qquad(j=1,\ldots,k),
\label{eq:lagrange-basis}
\end{equation}
the Lagrange polynomials of degree $k-1$ with $\ell_j(\mu_i)=\delta_{ij}$ (\cref{thm:interpolation}); for $k=1$, $\ell_1=1$.

\begin{lemma}[Spectral projections]
\label{lem:spectral-projections}
Let $A\in\mathbb F^{n\times n}$ be diagonalisable with the pairwise different eigenvalues $\mu_1,\ldots,\mu_k$, and let $P_j\coloneqq\ell_j(A)$.
\begin{enumerate}
\item $P_j$ is the projection onto $E_{\mu_j}$ along the sum of the other eigenspaces: if $\mathbf v=\mathbf v_1+\cdots+\mathbf v_k$ with $\mathbf v_i\in E_{\mu_i}$, then $P_j\mathbf v=\mathbf v_j$.
\item $P_1+\cdots+P_k=I$, $P_iP_j=0$ for $i\ne j$, $P_j^2=P_j$, and $q(A)=\sum_jq(\mu_j)P_j$ for every polynomial $q$; in particular $A=\sum_j\mu_jP_j$.
\item If a subspace $U$ is invariant under $A$, then $U=(U\cap E_{\mu_1})\oplus\cdots\oplus(U\cap E_{\mu_k})$, so $U$ has a basis of eigenvectors of $A$.
\end{enumerate}
\end{lemma}

\begin{proof}
(i) By \cref{thm:diagonalisable}\,(iv) the decomposition of $\mathbf v$ exists and is unique, and by \cref{prop:eigen-basic}\,(iv), $\ell_j(A)\mathbf v=\sum_i\ell_j(\mu_i)\mathbf v_i=\mathbf v_j$. (ii) For every $\mathbf v$, $\sum_jP_j\mathbf v=\sum_j\mathbf v_j=\mathbf v$, $P_iP_j\mathbf v=P_i\mathbf v_j=\mathbf 0$ for $i\ne j$, $P_j^2\mathbf v=P_j\mathbf v_j=\mathbf v_j$, and $q(A)\mathbf v=\sum_jq(\mu_j)\mathbf v_j=\sum_jq(\mu_j)P_j\mathbf v$. (iii) If $\mathbf u\in U$, then $P_j\mathbf u=\ell_j(A)\mathbf u\in U$, since $U$ is invariant under $A$ and therefore under every polynomial in $A$, and $P_j\mathbf u\in E_{\mu_j}$; so $\mathbf u=\sum_jP_j\mathbf u$ lies in the sum of the subspaces $U\cap E_{\mu_j}$, which is direct by \cref{cor:eigenspaces-direct}. Bases of these subspaces together form a basis of $U$, and they consist of eigenvectors.
\end{proof}

The matrices $P_j$ are the \term{spectral projections} of $A$, and (ii) is the \term{spectral decomposition} $A=\sum_j\mu_jP_j$: on the $j$th eigenspace, $A$ acts as multiplication by $\mu_j$. Part (iii) shows that the restriction of a diagonalisable matrix to an invariant subspace is again diagonalisable, and this gives a criterion for diagonalising two matrices with the same $S$.

\begin{theorem}[Simultaneous diagonalisation]
\label{thm:simultaneous}
Let $A,B\in\mathbb F^{n\times n}$ be diagonalisable over $\mathbb F$. Then $AB=BA$ if and only if there is an invertible $S\in\mathbb F^{n\times n}$ for which $S^{-1}AS$ and $S^{-1}BS$ are both diagonal, that is, if and only if $\mathbb F^n$ has a basis of vectors that are eigenvectors of both $A$ and $B$.
\end{theorem}

\begin{proof}
If $A=S\Lambda S^{-1}$ and $B=SMS^{-1}$ with diagonal $\Lambda$ and $M$, then $AB=S\Lambda MS^{-1}=SM\Lambda S^{-1}=BA$, since diagonal matrices commute (\cref{prop:diagonal}). Conversely, let $AB=BA$, and let $\mu$ be an eigenvalue of $A$. If $\mathbf x\in E_\mu(A)$, then $A(B\mathbf x)=B(A\mathbf x)=\mu B\mathbf x$, so $B\mathbf x\in E_\mu(A)$: the eigenspace $E_\mu(A)$ is invariant under $B$. By \cref{lem:spectral-projections}\,(iii), applied to $B$, it has a basis of eigenvectors of $B$, and these are eigenvectors of $A$ as well. Doing this for each eigenvalue of $A$ and using \cref{thm:diagonalisable}\,(iv) for $A$, we obtain a basis of $\mathbb F^n$ of common eigenvectors, the columns of the required $S$.
\end{proof}

Every polynomial in $A$ commutes with $A$, and if $A$ is diagonalisable, so is $q(A)$, with the same $S$: $q(A)=Sq(\Lambda)S^{-1}$. The hypothesis that both matrices are diagonalisable cannot be dropped, since the shear $K$ commutes with $I$ but is not diagonalisable. The circulant matrices of \cref{sec:structured-eigen} are an example: they all commute, and they are all diagonalised by the Fourier matrix.

\subsection{Left Eigenvectors}

A nonzero row vector $\mathbf y^{\mathsf H}$ with $\mathbf y^{\mathsf H}A=\lambda\mathbf y^{\mathsf H}$ is a \term{left eigenvector} of $A$ for $\lambda$; taking the conjugate transpose, the condition is $A^{\mathsf H}\mathbf y=\bar\lambda\mathbf y$, so the left eigenvectors of $A$ are the eigenvectors of $A^{\mathsf H}$, and by \cref{prop:charpoly-similar} the eigenvalues that have left eigenvectors are those of $A$. When the distinction matters, the eigenvectors of the definition are called \term{right eigenvectors}. Left and right eigenvectors for different eigenvalues are orthogonal, and in a diagonalisation they come in matching pairs.

\begin{proposition}[Left and right eigenvectors]
\label{prop:left-eigenvectors}
\leavevmode
\begin{enumerate}
\item If $A\mathbf x=\lambda\mathbf x$ and $\mathbf y^{\mathsf H}A=\mu\mathbf y^{\mathsf H}$ with $\lambda\ne\mu$, then $\mathbf y^{\mathsf H}\mathbf x=0$.
\item Let $A=S\Lambda S^{-1}$ with $\Lambda=\operatorname{diag}(\lambda_1,\ldots,\lambda_n)$, let $\mathbf x_1,\ldots,\mathbf x_n$ be the columns of $S$ and $\mathbf y_1^{\mathsf H},\ldots,\mathbf y_n^{\mathsf H}$ the rows of $S^{-1}$. Then $\mathbf y_i^{\mathsf H}A=\lambda_i\mathbf y_i^{\mathsf H}$, $\mathbf y_i^{\mathsf H}\mathbf x_j=\delta_{ij}$, and
\[
A=\sum_{i=1}^n\lambda_i\,\mathbf x_i\mathbf y_i^{\mathsf H},\qquad q(A)=\sum_{i=1}^nq(\lambda_i)\,\mathbf x_i\mathbf y_i^{\mathsf H}
\]
for every polynomial $q$. For each eigenvalue $\mu$, the sum of the matrices $\mathbf x_i\mathbf y_i^{\mathsf H}$ over the indices with $\lambda_i=\mu$ is the spectral projection of \cref{lem:spectral-projections}.
\end{enumerate}
\end{proposition}

\begin{proof}
(i) $\mu\mathbf y^{\mathsf H}\mathbf x=(\mathbf y^{\mathsf H}A)\mathbf x=\mathbf y^{\mathsf H}(A\mathbf x)=\lambda\mathbf y^{\mathsf H}\mathbf x$, and $\mu\ne\lambda$. (ii) From $A=S\Lambda S^{-1}$, $S^{-1}A=\Lambda S^{-1}$, whose $i$th row says $\mathbf y_i^{\mathsf H}A=\lambda_i\mathbf y_i^{\mathsf H}$, and $S^{-1}S=I$ says $\mathbf y_i^{\mathsf H}\mathbf x_j=\delta_{ij}$. The product $S\Lambda S^{-1}$ is the sum of the outer products of the columns of $S\Lambda$, which are $\lambda_i\mathbf x_i$, with the rows of $S^{-1}$ (\cref{thm:four-readings}), and $q(A)=Sq(\Lambda)S^{-1}$ by \cref{prop:similarity-invariants}\,(ii). The sum $Q_\mu$ of the matrices $\mathbf x_i\mathbf y_i^{\mathsf H}$ with $\lambda_i=\mu$ maps $\mathbf x_j$ to $\mathbf x_j$ if $\lambda_j=\mu$ and to $\mathbf 0$ otherwise, by the biorthogonality; so it agrees with the spectral projection of $\mu$ on the basis $\mathbf x_1,\ldots,\mathbf x_n$, and the two are equal.
\end{proof}

For the matrix $B$ of \eqref{eq:two-by-two-examples}, the eigenvectors $\mathbf x_1=\mathbf e_1$ and $\mathbf x_2=(1,1)$ for the eigenvalues $1$ and $2$ form the matrix $S$ with the rows $(1,1)$ and $(0,1)$, whose inverse has the rows $(1,-1)$ and $(0,1)$. So the left eigenvectors are $\mathbf y_1=(1,-1)$ and $\mathbf y_2=(0,1)$; indeed $\mathbf y_1^{\mathsf T}B=(1,-1)$ and $\mathbf y_2^{\mathsf T}B=(0,2)$. The spectral projections and the decomposition of $B$ are
\[
P_1=\mathbf x_1\mathbf y_1^{\mathsf T}=\begin{bmatrix}1&-1\\0&0\end{bmatrix},\qquad P_2=\mathbf x_2\mathbf y_2^{\mathsf T}=\begin{bmatrix}0&1\\0&1\end{bmatrix},\qquad B=P_1+2P_2 ,
\]
and $P_1$ is the oblique projection onto the $x_1$-axis along the line spanned by $(1,1)$. The left eigenvector $\mathbf y_1$ is orthogonal to the right eigenvector $\mathbf x_2$ and not to $\mathbf x_1$, as (i) requires. For a symmetric matrix the left and the right eigenvectors coincide, and when they are normalised, as for the matrix $A$ of \eqref{eq:two-by-two-examples}, the spectral projections are orthogonal projections.

\section{Powers of a Matrix}
\label{sec:matrix-powers}

\subsection{Powers of Diagonalisable Matrices}

If $A=S\Lambda S^{-1}$, the factors $S^{-1}S$ in the middle of a product cancel, and
\[
A^k=S\Lambda^kS^{-1}=\sum_{i=1}^n\lambda_i^k\,\mathbf x_i\mathbf y_i^{\mathsf H}\qquad(k\in\mathbb N),
\]
in the notation of \cref{prop:left-eigenvectors}; if $A$ is invertible, the same holds for negative $k$. The $k$th power of a diagonal matrix is the diagonal matrix of the $k$th powers, so in the basis of eigenvectors the $k$-fold application of $A$ multiplies the $i$th coordinate by $\lambda_i^k$. This is the formula that explains the growth of matrix powers.

The Fibonacci matrix $F$ with the rows $(1,1)$ and $(1,0)$ has the eigenvalues $\phi$ and $\hat\phi$, the roots of $\lambda^2=\lambda+1$, and the eigenvectors $(\phi,1)$ and $(\hat\phi,1)$, since $F(\phi,1)=(\phi+1,\phi)=(\phi^2,\phi)$. With
\[
S=\begin{bmatrix}\phi&\hat\phi\\1&1\end{bmatrix},\qquad S^{-1}=\frac1{\sqrt5}\begin{bmatrix}1&-\hat\phi\\-1&\phi\end{bmatrix},
\]
where $\det S=\phi-\hat\phi=\sqrt5$, the powers are
\[
F^k=S\begin{bmatrix}\phi^k&0\\0&\hat\phi^k\end{bmatrix}S^{-1}=\frac{\phi^k}{\sqrt5}\begin{bmatrix}\phi&-\phi\hat\phi\\1&-\hat\phi\end{bmatrix}-\frac{\hat\phi^k}{\sqrt5}\begin{bmatrix}\hat\phi&-\phi\hat\phi\\1&-\phi\end{bmatrix}.
\]
By \cref{sec:matrix-multiplication}, the entry $(2,1)$ of $F^k$ is $F_k$, and the formula gives $F_k=(\phi^k-\hat\phi^k)/\sqrt5$, which is Binet's formula \eqref{eq:binet}, now obtained from the eigenvalues. Since $|\hat\phi|<1$, the second term tends to zero, and since $\phi\hat\phi=-1$ and $-\hat\phi=1/\phi$,
\[
F^k=\frac{\phi^k}{\sqrt5}\begin{bmatrix}\phi&1\\1&1/\phi\end{bmatrix}+O(|\hat\phi|^k)\qquad(k\to\infty):
\]
all entries of $F^k$ grow like $\phi^k$, by the factor $\phi\approx1.618$ per step, and the columns of $F^k$ turn towards the eigenvector $(\phi,1)$. This is the explanation promised in \cref{sec:matrix-multiplication}.

\subsection{Linear Recurrences and the Dominant Eigenvalue}

The sequence $\mathbf x_{k+1}=A\mathbf x_k$, $\mathbf x_k=A^k\mathbf x_0$, is the simplest dynamical system: a state is mapped to the next state by a fixed linear rule. If $A$ is diagonalisable with the eigenvectors $\mathbf x_i$ and the left eigenvectors $\mathbf y_i$ of \cref{prop:left-eigenvectors}, the starting vector has the coordinates $c_i\coloneqq\mathbf y_i^{\mathsf H}\mathbf x_0$ in the basis of eigenvectors, $\mathbf x_0=\sum_ic_i\mathbf x_i$, and
\begin{equation}
A^k\mathbf x_0=\sum_{i=1}^nc_i\lambda_i^k\mathbf x_i=c_1\lambda_1^k\mathbf x_1+c_2\lambda_2^k\mathbf x_2+\cdots+c_n\lambda_n^k\mathbf x_n .
\label{eq:power-expansion}
\end{equation}
Each component evolves on its own, multiplied by its eigenvalue at every step. The continuous counterpart, the system of differential equations $\mathbf x'=A\mathbf x$, is solved by eigenvalues in the same way in Chapter~25, where the components grow like $e^{\lambda_it}$ and the real parts of the eigenvalues decide stability. Over a long time the components with the largest $|\lambda_i|$ dominate. An eigenvalue $\lambda_1$ is \term{dominant} if $|\lambda_1|>|\lambda_i|$ for all other eigenvalues counted with multiplicity, so that in particular $\operatorname{am}(\lambda_1)=1$.

\begin{proposition}[Dominant eigenvalue]
\label{prop:dominant}
Let $A\in\mathbb C^{n\times n}$ be diagonalisable with the eigenvalues $|\lambda_1|>|\lambda_2|\ge\cdots\ge|\lambda_n|$, and let $\mathbf x_0$ have the coordinate $c_1\ne0$ in \eqref{eq:power-expansion}. Then
\[
\Bigl\|\frac{A^k\mathbf x_0}{\lambda_1^k}-c_1\mathbf x_1\Bigr\|\le C\,\Bigl|\frac{\lambda_2}{\lambda_1}\Bigr|^k\qquad(k\in\mathbb N),\qquad C\coloneqq\sum_{i=2}^n|c_i|\,\|\mathbf x_i\| .
\]
In particular $A^k\mathbf x_0/\lambda_1^k\to c_1\mathbf x_1$, and the direction of $A^k\mathbf x_0$ approaches that of the eigenvector $\mathbf x_1$.
\end{proposition}

\begin{proof}
Divide \eqref{eq:power-expansion} by $\lambda_1^k$ and use the triangle inequality with $|\lambda_i/\lambda_1|^k\le|\lambda_2/\lambda_1|^k$ for $i\ge2$.
\end{proof}

The length of $A^k\mathbf x_0$ grows or decays like $|\lambda_1|^k$, and the angle between $A^k\mathbf x_0$ and $\mathbf x_1$ shrinks like $|\lambda_2/\lambda_1|^k$. \Cref{fig:powers-plane}A shows the iterates of six starting vectors under the matrix
\[
A=\begin{bmatrix}1&0.2\\0.4&0.8\end{bmatrix},
\]
with the eigenvalues $1.2$ and $0.6$, by \eqref{eq:eigen-2x2}, and the eigenvectors $(1,1)$ and $(1,-2)$. The component along $(1,-2)$ halves at every step relative to the other, since $0.6/1.2=1/2$, and every path bends into the direction of $(1,1)$, on which it then grows by the factor $1.2$. A starting vector on the line spanned by $(1,-2)$, for which $c_1=0$, would stay on that line and shrink to $\mathbf 0$; any other start, however close to that line, eventually follows $(1,1)$. This observation is the basis of the power method of \cref{sec:computing-eigenvalues}. In a population model in which $\mathbf x_k$ counts the individuals of each age group in year $k$ and the entries of $A$ are birth and survival rates, the proposition says that the proportions of the age groups settle into those of the dominant eigenvector, while the size of the population grows or decays like $\lambda_1^k$. Panel~B shows a real matrix with complex eigenvalues, which \cref{sec:complex-eigenvalues} explains.

\begin{figure}
\sdpanel{A}{\begin{tikzpicture}[scale=0.92, font=\scriptsize]
\draw[black!40, line width=0.5pt] (0,0) circle (1);
\draw[dashed, line width=0.5pt] (-2.4,-2.4) -- (2.4,2.4);
\draw[dashed, line width=0.5pt] (-1.2,2.4) -- (1.2,-2.4);
\foreach \p in {%
{(0.985,0.174) (1.020,0.533) (1.126,0.834) (1.293,1.118) (1.516,1.411) (1.799,1.736) (2.146,2.108)},
{(0.500,0.866) (0.673,0.893) (0.852,0.984) (1.048,1.128) (1.274,1.321) (1.538,1.567) (1.852,1.869)},
{(-0.707,0.707) (-0.566,0.283) (-0.509,0.000) (-0.509,-0.204) (-0.550,-0.367) (-0.623,-0.513) (-0.726,-0.660)},
{(-0.985,-0.174) (-1.020,-0.533) (-1.126,-0.834) (-1.293,-1.118) (-1.516,-1.411) (-1.799,-1.736) (-2.146,-2.108)},
{(-0.500,-0.866) (-0.673,-0.893) (-0.852,-0.984) (-1.048,-1.128) (-1.274,-1.321) (-1.538,-1.567) (-1.852,-1.869)},
{(0.707,-0.707) (0.566,-0.283) (0.509,-0.000) (0.509,0.204) (0.550,0.367) (0.623,0.513) (0.726,0.660)}}
{\draw[line width=0.6pt] plot[mark=*, mark size=1.1pt] coordinates {\p};}
\draw[line width=0.4pt, -{Latex[length=1.5mm]}] (-2.6,0) -- (2.7,0) node[below] {$x_1$};
\draw[line width=0.4pt, -{Latex[length=1.5mm]}] (0,-2.6) -- (0,2.7) node[left] {$x_2$};
\node[right] at (2.4,2.4) {$(1,1)$};
\node[right] at (1.15,-2.3) {$(1,-2)$};
\end{tikzpicture}}\hfill
\sdpanel{B}{\begin{tikzpicture}[scale=2.8, font=\scriptsize]
\draw[line width=0.6pt] plot coordinates {(1.000,0.000) (0.900,0.400) (0.610,0.640) (0.229,0.692) (-0.140,0.576) (-0.414,0.347) (-0.546,0.078) (-0.530,-0.164) (-0.395,-0.327) (-0.192,-0.387) (0.021,-0.348) (0.192,-0.235) (0.291,-0.088) (0.306,0.055) (0.247,0.161) (0.142,0.211) (0.022,0.205) (-0.082,0.152) (-0.150,0.074) (-0.172,-0.008) (-0.151,-0.075) (-0.098,-0.113) (-0.032,-0.118) (0.030,-0.096) (0.075,-0.055) (0.095,-0.008) (0.090,0.032) (0.065,0.058) (0.029,0.067) (-0.007,0.058) (-0.036,0.038) (-0.051,0.012) (-0.052,-0.012) (-0.041,-0.029) (-0.022,-0.037) (-0.002,-0.035) (0.016,-0.025) (0.027,-0.011) (0.030,0.003) (0.025,0.014) (0.016,0.020)};
\draw plot[only marks, mark=*, mark size=0.4pt] coordinates {(1.000,0.000) (0.900,0.400) (0.610,0.640) (0.229,0.692) (-0.140,0.576) (-0.414,0.347) (-0.546,0.078) (-0.530,-0.164) (-0.395,-0.327) (-0.192,-0.387) (0.021,-0.348) (0.192,-0.235) (0.291,-0.088) (0.306,0.055) (0.247,0.161) (0.142,0.211) (0.022,0.205) (-0.082,0.152) (-0.150,0.074) (-0.172,-0.008) (-0.151,-0.075)};
\node[above right=-1pt] at (1.0,0) {$\mathbf x_0$};
\draw[line width=0.4pt, -{Latex[length=1.5mm]}] (-0.75,0) -- (1.2,0) node[below] {$x_1$};
\draw[line width=0.4pt, -{Latex[length=1.5mm]}] (0,-0.5) -- (0,0.85) node[left] {$x_2$};
\end{tikzpicture}}
\caption{Iterates $\mathbf x_k=A^k\mathbf x_0$ in the plane, joined by lines. (A) The matrix with the rows $(1,0.2)$ and $(0.4,0.8)$ and the eigenvalues $1.2$ and $0.6$, for $k=0,\ldots,6$ and six starting vectors on the unit circle (grey); the dashed lines are the eigenvector lines. Every path turns into the direction of the dominant eigenvector $(1,1)$. (B) The matrix with the rows $(0.9,-0.5)$ and $(0.4,0.7)$ and the complex eigenvalues $0.8\pm0.436i$, of modulus $0.911$, for $k=0,\ldots,40$ and $\mathbf x_0=\mathbf e_1$ (dots for $k\le20$): the iterates turn by about $29$ degrees per step on average and spiral into the origin.\label{fig:powers-plane}}
\end{figure}

\subsection{The Spectral Radius}

\begin{definition}[Spectral radius]
\label{def:spectral-radius}
The \term{spectral radius} of $A\in\mathbb C^{n\times n}$ is $\rho(A)\coloneqq\max\{|\lambda|: \lambda\text{ an eigenvalue of }A\}$; for a real matrix the complex eigenvalues count.
\end{definition}

The spectral radius is the radius of the smallest disc about $0$ in the complex plane that contains all eigenvalues. It decides whether the powers of a matrix decay.

\begin{theorem}[Convergent matrices]
\label{thm:powers-to-zero}
For $A\in\mathbb C^{n\times n}$, $A^k\to0$ entrywise as $k\to\infty$ if and only if $\rho(A)<1$.
\end{theorem}

\begin{proof}
Suppose that $A^k\to0$, and let $A\mathbf x=\lambda\mathbf x$ with $\mathbf x\ne\mathbf 0$. Then $A^k\mathbf x=\lambda^k\mathbf x$, and each component of $A^k\mathbf x$ is a fixed combination of the entries of $A^k$, so $\lambda^k\mathbf x\to\mathbf 0$; for a component $x_j\ne0$ this gives $\lambda^k\to0$ and $|\lambda|<1$. So $\rho(A)<1$. Conversely, let $\rho(A)<1$. If $A=S\Lambda S^{-1}$ is diagonalisable, every entry of $A^k=S\Lambda^kS^{-1}$ is a fixed combination of the numbers $\lambda_i^k$, which tend to zero. For a general matrix the converse is proved in \cref{cor:powers-general}, with the help of the Cayley--Hamilton theorem of the next section.
\end{proof}

A matrix with $A^k\to0$ is called \term{convergent}. \Cref{fig:neumann} showed a convergent matrix that is not diagonalisable,
\[
A_2=\begin{bmatrix}4/5&2\\0&4/5\end{bmatrix},\qquad A_2^k=\begin{bmatrix}(4/5)^k&2k(4/5)^{k-1}\\0&(4/5)^k\end{bmatrix},
\]
the sum $(4/5)I+2N$ with the matrix $N$ of the rows $(0,1)$ and $(0,0)$: the single eigenvalue $4/5$ has algebraic multiplicity $2$ and geometric multiplicity $1$. Its powers tend to zero because the geometric factor $(4/5)^k$ eventually beats the linear factor $k$, after an initial growth of the entry $2k(4/5)^{k-1}$ that comes from this linear factor. The theorem settles the question left open after the Neumann series (\cref{thm:neumann}).

\begin{corollary}[Convergence of the Neumann series]
\label{cor:neumann-spectral}
For $A\in\mathbb C^{n\times n}$, the series $\sum_{k=0}^\infty A^k$ converges entrywise if and only if $\rho(A)<1$, and its sum is then $(I-A)^{-1}$.
\end{corollary}

\begin{proof}
If the series converges, its terms tend to zero, and $\rho(A)<1$ by \cref{thm:powers-to-zero}. Conversely, let $\rho(A)<1$. Since $A^k\to0$, there is an $m\ge1$ with $\|A^m\|_\infty<1$ (\cref{prop:operator-norms}\,(iv)), and by \cref{thm:neumann}, applied to $A^m$ (its proof holds verbatim for complex matrices), the series $\sum_jA^{mj}$ converges. For $N=mj+r$ with $0\le r<m$, the partial sum is
\[
\sum_{k=0}^NA^k=\Bigl(\sum_{i=0}^{j-1}A^{mi}\Bigr)(I+A+\cdots+A^{m-1})+A^{mj}(I+A+\cdots+A^r),
\]
and as $N\to\infty$, so that $j\to\infty$, the first term converges and the second tends to zero, since $A^{mj}\to0$ and the factor in parentheses takes only $m$ values. So the series converges, to a matrix $S$, and letting $N\to\infty$ in $(I-A)\sum_{k\le N}A^k=I-A^{N+1}$ gives $(I-A)S=I$.
\end{proof}

The Neumann series of $A_2$ converges although $\|A_2\|_\infty=14/5$, as \cref{sec:trace-norms} observed. The norm of a matrix, on the other hand, is always an upper bound for the spectral radius.

\begin{proposition}[Spectral radius and norms]
\label{prop:spectral-radius-norm}
Let $\|\cdot\|$ be one of the norms $\|\cdot\|_1$, $\|\cdot\|_2$, $\|\cdot\|_\infty$ or $\|\cdot\|_{\mathrm F}$ on $\mathbb C^{n\times n}$. Then $\rho(A)\le\|A\|$ and $\rho(A)\le\|A^k\|^{1/k}$ for every $A\in\mathbb C^{n\times n}$ and every $k\ge1$. For a real matrix the four norms have the same values whether they are computed over $\mathbb R$ or over $\mathbb C$.
\end{proposition}

\begin{proof}
Each of these norms satisfies $\|A\mathbf x\|_p\le\|A\|\|\mathbf x\|_p$ for complex vectors, with $p=1,2,\infty$ for the operator norms and $p=2$ for the Frobenius norm (\cref{prop:operator-norms} and \cref{sec:complex-matrices}). If $A\mathbf x=\lambda\mathbf x$ with $\mathbf x\ne\mathbf 0$, then $|\lambda|\|\mathbf x\|_p=\|A\mathbf x\|_p\le\|A\|\|\mathbf x\|_p$, so $|\lambda|\le\|A\|$. Applied to $A^k$, whose eigenvalues include $\lambda^k$ by \cref{prop:eigen-basic}\,(iv), this gives $|\lambda|^k\le\|A^k\|$. For a real matrix, the formulas of \cref{thm:norm-formulas} and the Frobenius norm involve only the moduli of the entries, and they do not depend on the field. For the $2$-norm, write a complex vector as $\mathbf z=\mathbf x+i\mathbf y$ with real $\mathbf x$ and $\mathbf y$; then $A\mathbf z=A\mathbf x+iA\mathbf y$ with real parts, and $\|A\mathbf z\|^2=\|A\mathbf x\|^2+\|A\mathbf y\|^2\le\|A\|_2^2(\|\mathbf x\|^2+\|\mathbf y\|^2)=\|A\|_2^2\|\mathbf z\|^2$, while real vectors are complex vectors.
\end{proof}

The bound can be far from sharp: $\rho(A_2)=4/5$, while $\|A_2\|_\infty=14/5$ and $\|A_2\|_2>2$, since $\|A_2\mathbf e_2\|>2$. For Hermitian matrices, \cref{ch:symmetric-matrices} shows that $\rho(A)=\|A\|_2$. The second bound becomes an equality in the limit.

\begin{theorem}[Gelfand's formula]
\label{thm:gelfand}
For every $A\in\mathbb C^{n\times n}$ and each norm of \cref{prop:spectral-radius-norm}, $\rho(A)=\lim_{k\to\infty}\|A^k\|^{1/k}$.
\end{theorem}

\begin{proof}
By \cref{prop:spectral-radius-norm}, $\rho(A)\le\|A^k\|^{1/k}$ for every $k$. Let $\varepsilon>0$ and $B\coloneqq A/(\rho(A)+\varepsilon)$. The eigenvalues of $B$ are those of $A$ divided by $\rho(A)+\varepsilon$, so $\rho(B)<1$, and $B^k\to0$ by \cref{thm:powers-to-zero}. Every entry of $B^k$ tends to zero and each of the four norms is at most $n$ times the largest modulus of an entry (\cref{prop:operator-norms}\,(iv)), so $\|B^k\|\le1$ for all $k$ from some $k_0$ on, that is, $\|A^k\|^{1/k}\le\rho(A)+\varepsilon$ for $k\ge k_0$. So $\rho(A)\le\|A^k\|^{1/k}\le\rho(A)+\varepsilon$ for all large $k$, and $\varepsilon$ was arbitrary.
\end{proof}

For $A_2$, $\|A_2^k\|_\infty=(4/5)^k+2k(4/5)^{k-1}$, and its $k$th root tends to $4/5$, because $(1+5k/2)^{1/k}\to1$. Gelfand's formula says that over many steps a matrix multiplies lengths by $\rho(A)$ per step on average, whatever it does in a single step; the transient growth of \cref{fig:neumann} is the difference between the single step and the average.

\subsection{Polynomials in a Matrix}

If $A\mathbf x=\lambda\mathbf x$, then $q(A)\mathbf x=q(\lambda)\mathbf x$, by \cref{prop:eigen-basic}\,(iv). The eigenvalues of $q(A)$ are therefore among the numbers $q(\lambda_i)$, and in fact they are exactly these numbers, with the right multiplicities.

\begin{theorem}[Eigenvalues of polynomials in a matrix]
\label{thm:spectral-mapping}
Let $A\in\mathbb C^{n\times n}$ have the eigenvalues $\lambda_1,\ldots,\lambda_n$, counted with multiplicity.
\begin{enumerate}
\item For every polynomial $q$, $p_{q(A)}(x)=\prod_{i=1}^n\bigl(x-q(\lambda_i)\bigr)$: the eigenvalues of $q(A)$, counted with multiplicity, are $q(\lambda_1),\ldots,q(\lambda_n)$.
\item If $A$ is invertible, then $p_{A^{-1}}(x)=\prod_{i=1}^n(x-1/\lambda_i)$: the eigenvalues of $A^{-1}$ are the numbers $1/\lambda_i$.
\end{enumerate}
\end{theorem}

\begin{proof}
(i) If $q=c$ is constant, $q(A)=cI$ and both sides are $(x-c)^n$. Let $q$ have degree $d\ge1$ and the leading coefficient $c$, and fix $x\in\mathbb C$. By \cref{cor:factor}, the polynomial $t\mapsto x-q(t)$, of degree $d$ in $t$ with the leading coefficient $-c$, factors as $x-q(t)=-c\prod_{j=1}^d(t-\tau_j)$ with some $\tau_j\in\mathbb C$ that depend on $x$. Substituting $A$ for $t$, which is allowed because polynomials in $A$ multiply like polynomials (\cref{sec:matrix-multiplication}), gives $xI-q(A)=-c\prod_j(A-\tau_jI)$. By the product rule and $\det(A-\tau I)=(-1)^n\prod_i(\tau-\lambda_i)$,
\[
\det\bigl(xI-q(A)\bigr)=(-c)^n\prod_{j=1}^d(-1)^n\prod_{i=1}^n(\tau_j-\lambda_i)=(-c)^n\prod_{i=1}^n\prod_{j=1}^d(\lambda_i-\tau_j)=\prod_{i=1}^n\bigl(x-q(\lambda_i)\bigr),
\]
where the middle step uses $(-1)^{nd}\prod_{i,j}(\tau_j-\lambda_i)=\prod_{i,j}(\lambda_i-\tau_j)$, and the last uses $-c\prod_j(\lambda_i-\tau_j)=x-q(\lambda_i)$, the factorisation evaluated at $t=\lambda_i$. Both sides are polynomials in $x$, and they agree for every $x$. (ii) The eigenvalues are nonzero, since $\det A=\prod_i\lambda_i\ne0$. For $x\ne0$, $xI-A^{-1}=xA^{-1}(A-x^{-1}I)$, so
\[
\det(xI-A^{-1})=\frac{x^n}{\det A}\det(A-x^{-1}I)=\frac{x^n}{\prod_i\lambda_i}\prod_{i=1}^n(\lambda_i-x^{-1})=\prod_{i=1}^n\Bigl(x-\frac1{\lambda_i}\Bigr),
\]
and the two polynomials agree for every $x\ne0$.
\end{proof}

The theorem is known as the \term{spectral mapping theorem} for polynomials. Its first special cases are used constantly: $A^k$ has the eigenvalues $\lambda_i^k$, so that $\operatorname{tr}A^k=\sum_i\lambda_i^k$ by \cref{thm:trace-det-eigen}; the \term{shifted} matrix $A-\sigma I$ has the eigenvalues $\lambda_i-\sigma$, with the same eigenvectors; and, combining the two parts, $(A-\sigma I)^{-1}$ has the eigenvalues $1/(\lambda_i-\sigma)$ when $\sigma$ is not an eigenvalue. The last matrix turns the eigenvalue nearest to $\sigma$ into the eigenvalue of largest modulus, and \cref{sec:computing-eigenvalues} exploits this. The eigenvectors of $A$ are eigenvectors of $q(A)$, but $q(A)$ may have more: for the Jordan block $J_2(0)$ and $q(x)=x^2$, $q(J_2(0))=0$, for which every nonzero vector is an eigenvector.

\section{The Cayley--Hamilton Theorem}
\label{sec:cayley-hamilton}

\subsection{Every Matrix Satisfies Its Characteristic Polynomial}

The $n^2+1$ matrices $I,A,A^2,\ldots,A^{n^2}$ lie in the space $\mathbb F^{n\times n}$ of dimension $n^2$, so they are linearly dependent (\cref{thm:dimension}), and some nonzero polynomial $q$ of degree at most $n^2$ satisfies $q(A)=0$. A far better polynomial is available: one of degree $n$, which we already know.

\begin{theorem}[Cayley--Hamilton theorem]
\label{thm:cayley-hamilton}
For every $A\in\mathbb F^{n\times n}$, $p_A(A)=0$: if $p_A(\lambda)=\lambda^n+c_{n-1}\lambda^{n-1}+\cdots+c_0$, then
\[
A^n+c_{n-1}A^{n-1}+\cdots+c_1A+c_0I=0 .
\]
\end{theorem}

\begin{proof}
Let $B(\lambda)\coloneqq\operatorname{adj}(\lambda I-A)$. Its entries are cofactors of $\lambda I-A$, determinants of order $n-1$ whose entries are polynomials of degree at most $1$ in $\lambda$, so they are polynomials of degree at most $n-1$. Collecting the coefficients of each power of $\lambda$ entry by entry, we can write
\[
B(\lambda)=B_0+\lambda B_1+\cdots+\lambda^{n-1}B_{n-1}
\]
with matrices $B_k\in\mathbb F^{n\times n}$ that do not depend on $\lambda$. By the adjugate formula (\cref{thm:adjugate}), for every $\lambda\in\mathbb F$,
\[
(\lambda I-A)B(\lambda)=\det(\lambda I-A)\,I=p_A(\lambda)\,I .
\]
With $B_{-1}\coloneqq0$, $B_n\coloneqq0$ and $c_n\coloneqq1$, the left side is $\sum_{k=0}^n\lambda^k(B_{k-1}-AB_k)$ and the right side is $\sum_{k=0}^n\lambda^kc_kI$. Each entry of the two sides is a polynomial in $\lambda$, and the two agree for every $\lambda$, so their coefficients agree:
\[
B_{k-1}-AB_k=c_kI\qquad(k=0,1,\ldots,n).
\]
Multiply the $k$th equation by $A^k$ from the left and add:
\[
\sum_{k=0}^n\bigl(A^kB_{k-1}-A^{k+1}B_k\bigr)=\sum_{k=0}^nc_kA^k=p_A(A).
\]
On the left, the sum telescopes: the term $-A^{k+1}B_k$ cancels the term $A^{k+1}B_k$ of the next index, and what remains is $A^0B_{-1}-A^{n+1}B_n=0$.
\end{proof}

A tempting one-line argument is wrong: ``substituting $\lambda=A$ in $p_A(\lambda)=\det(\lambda I-A)$ gives $\det(AI-A)=\det0=0$''. The left side $p_A(A)$ is a matrix, the right side a number, and the substitution is meaningless: in $\lambda I-A$ the variable $\lambda$ multiplies the identity matrix, and replacing it by a matrix does not produce the matrix $A-A$. The theorem is a statement about the polynomial $p_A$, computed first with $\lambda$ as a number, and then evaluated at $A$. For $n=2$ it can be checked directly: for $A$ with the rows $(a,b)$ and $(c,d)$,
\[
A^2-(a+d)A=\begin{bmatrix}a^2+bc&ab+bd\\ac+cd&bc+d^2\end{bmatrix}-\begin{bmatrix}a^2+ad&ab+bd\\ac+cd&ad+d^2\end{bmatrix}=(bc-ad)I ,
\]
so $A^2-(\operatorname{tr}A)A+(\det A)I=0$. For a diagonalisable matrix the theorem is also easy: $p_A(A)=Sp_A(\Lambda)S^{-1}$, and $p_A(\Lambda)$ is the diagonal matrix of the numbers $p_A(\lambda_i)=0$. Together with the density of diagonalisable matrices in the note of \cref{sec:diagonalisation} and the continuity of $A\mapsto p_A(A)$ in the entries, this gives another proof for complex matrices; the proof above needs neither.

\subsection{Consequences}

If $\det A\ne0$, the constant coefficient $c_0=(-1)^n\det A$ is nonzero, and the theorem can be written as $A(A^{n-1}+c_{n-1}A^{n-2}+\cdots+c_1I)=-c_0I$. Hence
\[
A^{-1}=-\frac1{c_0}\bigl(A^{n-1}+c_{n-1}A^{n-2}+\cdots+c_1I\bigr):
\]
the inverse of a matrix is a polynomial in the matrix, of degree at most $n-1$. For $T_3$, with $p_{T_3}(\lambda)=\lambda^3-6\lambda^2+10\lambda-4$, this gives $T_3^{-1}=(T_3^2-6T_3+10I)/4$, and indeed
\[
T_3^2=\begin{bmatrix}5&-4&1\\-4&6&-4\\1&-4&5\end{bmatrix},\qquad T_3^2-6T_3+10I=\begin{bmatrix}3&2&1\\2&4&2\\1&2&3\end{bmatrix}=4\,T_3^{-1},
\]
in agreement with \cref{prop:second-difference-inverse}. In the same way, every power of $A$ reduces to a combination of $I,A,\ldots,A^{n-1}$. Polynomials can be divided with remainder, like integers (\cref{thm:division}): for polynomials $a$ and $b\ne0$ there are polynomials $s$ and $r$ with $a=sb+r$ and $\deg r<\deg b$, obtained by repeatedly subtracting a multiple $cx^jb$ of $b$ that cancels the highest term of what is left. Dividing $x^k$ by $p_A$ gives $x^k=s(x)p_A(x)+r(x)$ with $\deg r<n$, and then $A^k=s(A)p_A(A)+r(A)=r(A)$. For matrices of order $2$, $A^2=(\operatorname{tr}A)A-(\det A)I$. For the Fibonacci matrix this is $F^2=F+I$, and induction gives
\[
F^k=F_kF+F_{k-1}I\qquad(k\ge1),
\]
since $F^{k+1}=F_kF^2+F_{k-1}F=(F_k+F_{k-1})F+F_kI$; the right side is the matrix with the rows $(F_{k+1},F_k)$ and $(F_k,F_{k-1})$ found in \cref{sec:matrix-multiplication}.

The Cayley--Hamilton theorem also completes the proof of \cref{thm:powers-to-zero}. Besides the theorem, the argument needs one fact about sequences: a geometric sequence with ratio of modulus less than $1$ still tends to zero when a vanishing disturbance is added at every step.

\begin{lemma}[A disturbed geometric sequence]
\label{lem:disturbed-geometric}
Let $\mu\in\mathbb C$ with $|\mu|<1$, and let $(t_k)$ and $(u_k)$ be sequences of complex numbers with $t_{k+1}=\mu t_k+u_k$ for all $k\in\mathbb N$. If $u_k\to0$, then $t_k\to0$.
\end{lemma}

\begin{proof}
By induction on $k$, $t_k=\mu^kt_0+\sum_{j=0}^{k-1}\mu^{k-1-j}u_j$. The convergent sequence $(u_k)$ is bounded, say by $U$. Let $\varepsilon>0$, and choose $J$ with $|u_j|\le\varepsilon$ for all $j\ge J$. For $k>J$, split the sum at $J$. In the first part $k-1-j\ge k-J$, and the second part is at most $\varepsilon$ times a geometric sum, so
\[
|t_k|\le|\mu|^k|t_0|+\sum_{j=0}^{J-1}|\mu|^{k-1-j}U+\varepsilon\sum_{j=J}^{k-1}|\mu|^{k-1-j}\le|\mu|^k|t_0|+JU|\mu|^{k-J}+\frac{\varepsilon}{1-|\mu|}.
\]
The first two terms tend to zero as $k\to\infty$ (\cref{thm:standard-limits}), so $|t_k|\le2\varepsilon/(1-|\mu|)$ for all large $k$. Since $\varepsilon$ was arbitrary, $t_k\to0$.
\end{proof}

\begin{corollary}[Powers of a general matrix]
\label{cor:powers-general}
If $A\in\mathbb C^{n\times n}$ and $\rho(A)<1$, then $A^k\to0$.
\end{corollary}

\begin{proof}
Let $\lambda_1,\ldots,\lambda_n$ be the eigenvalues of $A$ counted with multiplicity, so that $p_A(x)=(x-\lambda_1)\cdots(x-\lambda_n)$ (\cref{thm:eigenvalues-exist}). For $r=0,1,\ldots,n$ let
\[
B_r\coloneqq(A-\lambda_{r+1}I)(A-\lambda_{r+2}I)\cdots(A-\lambda_nI),\qquad B_n\coloneqq I .
\]
Polynomials in $A$ multiply like polynomials, so $B_0=p_A(A)$, which is $0$ by \cref{thm:cayley-hamilton}, and $B_{r-1}=(A-\lambda_rI)B_r$ for $r\ge1$. Consider the sequences of matrices $X^{(r)}_k\coloneqq A^kB_r$, $k\in\mathbb N$. Since $A^{k+1}=\lambda_rA^k+A^k(A-\lambda_rI)$,
\[
X^{(r)}_{k+1}=\lambda_rX^{(r)}_k+A^k(A-\lambda_rI)B_r=\lambda_rX^{(r)}_k+X^{(r-1)}_k\qquad(r\ge1).
\]
We show by induction on $r$ that $X^{(r)}_k\to0$ entrywise as $k\to\infty$. For $r=0$, $X^{(0)}_k=A^kB_0=0$ for every $k$. If $X^{(r-1)}_k\to0$, then for each position $(i,j)$ the entries $t_k\coloneqq(X^{(r)}_k)_{ij}$ satisfy $t_{k+1}=\lambda_rt_k+u_k$ with $u_k\coloneqq(X^{(r-1)}_k)_{ij}\to0$, and $|\lambda_r|\le\rho(A)<1$, so $t_k\to0$ by \cref{lem:disturbed-geometric}. For $r=n$ the sequence is $X^{(n)}_k=A^k$.
\end{proof}

The proof shows that $A^k\to0$ without describing the entries of $A^k$, and more can be said about them. For each position $(i,j)$, the sequence $s_k\coloneqq(A^k)_{ij}$ satisfies the linear recurrence whose characteristic polynomial is $p_A$, since $A^kp_A(A)=0$ for every $k$. Let $m\coloneqq\operatorname{am}(0)$ and $q(x)\coloneqq p_A(x)/x^m$, a polynomial with a nonzero constant term. Since $A^{k+m}q(A)=A^kp_A(A)=0$, the sequence $(A^{k+m})_{ij}$ satisfies the recurrence with the characteristic polynomial $q$, and \cref{prop:linear-recurrence-k} gives it the form $\sum_iP_i(k)\mu_i^k$, where $\mu_i$ runs through the different nonzero eigenvalues and $P_i$ is a polynomial of degree less than $\operatorname{am}(\mu_i)$. This is the form of the example $A_2$, whose entries are $(4/5)^k$ and $2k(4/5)^{k-1}$: a polynomial factor of degree less than the multiplicity, times the power of the eigenvalue. For a diagonalisable matrix the factors are constants, by the formula for $A^k$ in \cref{sec:matrix-powers}. Chapter~22 obtains the same form from the Jordan form, with the binomial coefficients of the powers of Jordan blocks as the polynomials.

\subsection{The Minimal Polynomial}

The characteristic polynomial is one polynomial that annihilates $A$; there may be others of lower degree. The identity $I_n$ satisfies $I-I=0$, a polynomial of degree $1$, while $p_I(x)=(x-1)^n$.

\begin{definition}[Minimal polynomial]
\label{def:minimal-polynomial}
The \term{minimal polynomial} $m_A$ of $A\in\mathbb F^{n\times n}$ is the monic polynomial of least degree with $m_A(A)=0$.
\end{definition}

By \cref{thm:cayley-hamilton} there is a monic polynomial that annihilates $A$, namely $p_A$, so there is one of least degree, and the next proposition shows that it is unique.

\begin{proposition}[Properties of the minimal polynomial]
\label{prop:minimal-polynomial}
Let $A\in\mathbb F^{n\times n}$.
\begin{enumerate}
\item If $q(A)=0$ for a polynomial $q$, then $m_A$ divides $q$. In particular $m_A$ divides $p_A$, and $m_A$ is unique.
\item The roots of $m_A$ in $\mathbb F$ are exactly the eigenvalues of $A$ in $\mathbb F$.
\item Similar matrices have the same minimal polynomial.
\end{enumerate}
\end{proposition}

\begin{proof}
(i) Divide with remainder, $q=sm_A+r$ with $\deg r<\deg m_A$. Then $r(A)=q(A)-s(A)m_A(A)=0$. If $r\ne0$, dividing $r$ by its leading coefficient would give a monic polynomial of smaller degree than $m_A$ that annihilates $A$; so $r=0$. If $m$ and $m'$ are both monic of least degree with $m(A)=m'(A)=0$, each divides the other, and since both are monic of the same degree, $m=m'$. (ii) If $A\mathbf x=\lambda\mathbf x$ with $\mathbf x\ne\mathbf 0$, then $\mathbf 0=m_A(A)\mathbf x=m_A(\lambda)\mathbf x$ by \cref{prop:eigen-basic}\,(iv), so $m_A(\lambda)=0$. Conversely, if $m_A(\mu)=0$, then $m_A(x)=(x-\mu)r(x)$ by \cref{cor:factor}\,(i), with a monic $r$ of smaller degree, so $r(A)\ne0$ and there is a vector $\mathbf v$ with $\mathbf w\coloneqq r(A)\mathbf v\ne\mathbf 0$. Then $(A-\mu I)\mathbf w=m_A(A)\mathbf v=\mathbf 0$, and $\mu$ is an eigenvalue with the eigenvector $\mathbf w$. (iii) By \cref{prop:similarity-invariants}\,(ii), $q(B)=0$ if and only if $q(A)=0$ for similar $A$ and $B$.
\end{proof}

By (i) and (ii), the minimal polynomial has the same roots as the characteristic polynomial, with multiplicities that are at least $1$ and at most the algebraic multiplicities. The shear $K$ has $m_K=(x-1)^2$, since $K-I\ne0$, while $m_I=x-1$; the Jordan block $J_n(\lambda)$ has $m=(x-\lambda)^n$, since $N\coloneqq J_n(\lambda)-\lambda I$ maps $\mathbf e_k$ to $\mathbf e_{k-1}$ for $k\ge2$, so that $N^{n-1}\mathbf e_n=\mathbf e_1$ and $N^{n-1}\ne0$; and $\operatorname{diag}(2,2,3)$ has $m=(x-2)(x-3)$. For the matrix $A$ of the complete example in \cref{sec:diagonalisation}, the product $(A-3I)(A-5I)$ is zero, as one checks by multiplying, so $m_A=(x-3)(x-5)$, while for the defective matrix $A'$ the corresponding product has the entry $-1$ in position $(2,1)$, and $m_{A'}=(x-3)^2(x-5)=p_{A'}$. The pattern is general.

\begin{theorem}[Diagonalisability and the minimal polynomial]
\label{thm:diagonalisable-minimal}
For $A\in\mathbb F^{n\times n}$ the following are equivalent:
\begin{enumerate}
\item $A$ is diagonalisable over $\mathbb F$;
\item $m_A$ is a product of pairwise different linear factors $x-\mu_j$ with $\mu_j\in\mathbb F$;
\item $q(A)=0$ for some polynomial $q$ that is a product of pairwise different linear factors over $\mathbb F$.
\end{enumerate}
\end{theorem}

\begin{proof}
(i)$\Rightarrow$(ii) Let $A=S\Lambda S^{-1}$ with the pairwise different eigenvalues $\mu_1,\ldots,\mu_k$ and $q(x)\coloneqq\prod_j(x-\mu_j)$. The diagonal matrix $q(\Lambda)$ has the entries $q(\lambda_i)=0$, so $q(A)=Sq(\Lambda)S^{-1}=0$, and $m_A$ divides $q$ by \cref{prop:minimal-polynomial}; a monic divisor of a product of different linear factors is a product of some of them, by the uniqueness of the factorisation. (ii)$\Rightarrow$(iii) Take $q=m_A$. (iii)$\Rightarrow$(i) Let $q(x)=\prod_{j=1}^k(x-\mu_j)$ with pairwise different $\mu_j$, and let $\ell_j$ be the Lagrange polynomials \eqref{eq:lagrange-basis}. The polynomial $\ell_1+\cdots+\ell_k-1$ has degree less than $k$ and vanishes at the $k$ points $\mu_i$, so it is zero (\cref{cor:factor}), and $\sum_j\ell_j(A)=I$. Since $(x-\mu_j)\ell_j(x)$ is a constant multiple of $q(x)$, $(A-\mu_jI)\ell_j(A)=0$, so $\ell_j(A)\mathbf v\in E_{\mu_j}$ for every $\mathbf v$. Hence every vector $\mathbf v=\sum_j\ell_j(A)\mathbf v$ is a sum of vectors from eigenspaces; the eigenspaces of the $\mu_j$ that are not eigenvalues are zero, and the others form a direct sum (\cref{cor:eigenspaces-direct}) that is all of $\mathbb F^n$. By \cref{thm:diagonalisable}\,(iv), $A$ is diagonalisable.
\end{proof}

The theorem decides diagonalisability from a single polynomial identity, often without computing any eigenvector.
\begin{itemize}
\item A projection, $P^2=P$, satisfies $q(P)=0$ for $q(x)=x(x-1)$, so it is diagonalisable with eigenvalues in $\{0,1\}$, as \cref{prop:projection} found by other means.
\item An involution, $R^2=I$, satisfies $q(R)=0$ for $q(x)=(x-1)(x+1)$, so it is diagonalisable with eigenvalues in $\{1,-1\}$; the reflections are examples.
\item A matrix with $A^m=I$ for some $m\ge1$, such as a permutation matrix (\cref{sec:structured-eigen}), is annihilated by $x^m-1$, which has the $m$ different roots of unity of \cref{prop:roots-of-unity} as roots; so it is diagonalisable over $\mathbb C$, and its eigenvalues are $m$th roots of unity.
\item A nilpotent matrix, $N^k=0$, has $m_N=x^j$ for some $j\le k$, by (i) of the proposition, so its only eigenvalue is $0$, its characteristic polynomial is $x^n$ (by \cref{thm:eigenvalues-exist}), and $N^n=0$ by the Cayley--Hamilton theorem, which is the statement of \cref{sec:rank-inequalities} once more. A nilpotent $N\ne0$ has $j\ge2$ and is not diagonalisable. Its trace and its determinant are $0$.
\end{itemize}
If $U$ is a subspace invariant under $A$, the minimal polynomial of the restriction $A|_U$ divides $m_A$, since $m_A(A|_U)=m_A(A)|_U=0$; with the theorem this gives a second proof that the restriction of a diagonalisable matrix to an invariant subspace is diagonalisable (\cref{lem:spectral-projections}\,(iii)). Chapter~22 shows that the exponent of $x-\mu$ in $m_A$ is the size of the largest Jordan block for the eigenvalue $\mu$.

\section{Complex Eigenvalues of Real Matrices}
\label{sec:complex-eigenvalues}

\subsection{Rotation and Scaling in an Invariant Plane}

A real matrix with a nonreal eigenvalue has no real eigenvector for it, but it still has a real invariant subspace: a plane on which it acts as a rotation combined with a scaling. For real $a$ and $b$ let
\[
C_{a,b}\coloneqq\begin{bmatrix}a&-b\\b&a\end{bmatrix}=r\begin{bmatrix}\cos\theta&-\sin\theta\\\sin\theta&\cos\theta\end{bmatrix}=rR_\theta\qquad\text{where }a+ib=re^{i\theta}.
\]
This is the matrix of multiplication by $a+ib$ on $\mathbb C$, identified with $\mathbb R^2$, as \cref{sec:complex-matrices} observed; its eigenvalues are $a\pm ib$, by \eqref{eq:eigen-2x2}.

\begin{proposition}[Real form of a complex eigenvalue]
\label{prop:real-2x2-form}
Let $A\in\mathbb R^{n\times n}$ have the eigenvalue $\lambda=a+ib$ with $b\ne0$, and let $\mathbf x=\mathbf u+i\mathbf v$, with $\mathbf u,\mathbf v\in\mathbb R^n$, be an eigenvector for the eigenvalue $\bar\lambda=a-ib$. Then $\mathbf u$ and $\mathbf v$ are linearly independent, their span is invariant under $A$, and
\[
A\mathbf u=a\mathbf u+b\mathbf v,\qquad A\mathbf v=-b\mathbf u+a\mathbf v,\qquad\text{that is,}\qquad A\begin{bmatrix}\mathbf u&\mathbf v\end{bmatrix}=\begin{bmatrix}\mathbf u&\mathbf v\end{bmatrix}C_{a,b}.
\]
In particular, a real $2\times2$ matrix with the eigenvalues $a\pm ib$, $b\ne0$, is similar over $\mathbb R$ to $C_{a,b}$, with the real matrix $S$ of the columns $\mathbf u$ and $\mathbf v$.
\end{proposition}

\begin{proof}
By \cref{thm:eigenvalues-exist}\,(ii), $\bar\lambda$ is an eigenvalue whenever $\lambda$ is. Separating real and imaginary parts in
\[
A\mathbf u+iA\mathbf v=A\mathbf x=(a-ib)(\mathbf u+i\mathbf v)=(a\mathbf u+b\mathbf v)+i(-b\mathbf u+a\mathbf v),
\]
in which $A\mathbf u$ and $A\mathbf v$ are real, gives the two equations. If $\mathbf u$ and $\mathbf v$ were dependent, there would be a real vector $\mathbf w\ne\mathbf 0$ and a complex number $c\ne0$ with $\mathbf x=c\mathbf w$, since $\mathbf x\ne\mathbf 0$; then $A\mathbf w=\bar\lambda\mathbf w$, where the left side is real and the right side has the imaginary part $-b\mathbf w\ne\mathbf 0$, a contradiction. For $n=2$, $\mathbf u$ and $\mathbf v$ form a basis of $\mathbb R^2$, so $S$ is invertible and $S^{-1}AS=C_{a,b}$.
\end{proof}

In the basis $\mathbf u,\mathbf v$, the matrix acts as the rotation by $\theta$ followed by the scaling by $r=|\lambda|$: a real matrix of order $2$ with complex eigenvalues is a complex number in disguise. The rotation $R_\varphi$ itself is $C_{\cos\varphi,\sin\varphi}$, with the eigenvalues $e^{\pm i\varphi}$, which explains its complex diagonalisation in \cref{sec:similarity}. For the matrix $A$ of \cref{fig:powers-plane}B, with the rows $(0.9,-0.5)$ and $(0.4,0.7)$, the trace is $1.6$ and the determinant $0.83$, so the eigenvalues are $0.8\pm i\sqrt{0.19}\approx0.8\pm0.436i$, of modulus $\sqrt{0.83}\approx0.911$ and argument $\theta\approx28.6$ degrees. The first row of $A-\bar\lambda I$ is $(0.1+i\sqrt{0.19},-0.5)$, so $\mathbf x=(0.5,\,0.1+i\sqrt{0.19})$ is an eigenvector for $\bar\lambda$, with $\mathbf u=(0.5,0.1)$ and $\mathbf v=(0,\sqrt{0.19})$. In these coordinates $A^k=SC_{a,b}^kS^{-1}=r^kSR_{k\theta}S^{-1}$: the iterate $\mathbf x_k$ is obtained by rotating the coordinates of $\mathbf x_0$ through $k\theta$ and shrinking them by $r^k$. In the standard coordinates the circles of the rotation become the ellipses $S(\text{circle})$, and the iterates spiral into the origin along shrinking ellipses, as the figure shows. For $r>1$ they would spiral outwards, and for $r=1$ they would stay on one ellipse.

\subsection{The Real Block Diagonal Form}

Pairs of complex eigenvalues of a real matrix can be handled in this way all at once.

\begin{theorem}[Real block diagonal form]
\label{thm:real-block-diagonal}
Let $A\in\mathbb R^{n\times n}$ be diagonalisable over $\mathbb C$, with the real eigenvalues $\lambda_1,\ldots,\lambda_r$ and the nonreal eigenvalues $a_1\pm ib_1,\ldots,a_s\pm ib_s$ ($b_j>0$), counted with multiplicity, so that $r+2s=n$. Then there is an invertible real matrix $S$ with
\[
S^{-1}AS=\begin{bmatrix}\operatorname{diag}(\lambda_1,\ldots,\lambda_r)&&&\\&C_{a_1,b_1}&&\\&&\ddots&\\&&&C_{a_s,b_s}\end{bmatrix},
\]
a block diagonal matrix with $r$ blocks of order $1$ and $s$ blocks of order $2$.
\end{theorem}

\begin{proof}
For each real eigenvalue $\mu$, choose a basis of the real eigenspace $\mathcal N(A-\mu I)\subseteq\mathbb R^n$; by the remark after \cref{def:multiplicities} it has $\operatorname{gm}(\mu)$ vectors $\mathbf w_j$, and these are also independent over $\mathbb C$: if $\sum_jc_j\mathbf w_j=\mathbf 0$ with $c_j=\alpha_j+i\beta_j$ and real $\alpha_j,\beta_j$, the real and the imaginary parts of this vector are $\sum_j\alpha_j\mathbf w_j=\mathbf 0$ and $\sum_j\beta_j\mathbf w_j=\mathbf 0$, so all $\alpha_j$ and $\beta_j$ vanish. Being $\operatorname{gm}(\mu)$ independent vectors of $E_\mu$, they form a basis of the complex eigenspace. For each nonreal eigenvalue $\mu=a+ib$ with $b>0$, choose a basis $\mathbf x_1,\ldots,\mathbf x_g$ of $E_{\bar\mu}$; by \cref{thm:eigenvalues-exist}\,(ii), conjugation maps $E_{\bar\mu}$ to $E_\mu$, and it maps independent vectors to independent vectors, so $\bar{\mathbf x}_1,\ldots,\bar{\mathbf x}_g$ is a basis of $E_\mu$. By \cref{thm:diagonalisable}\,(iv) all these vectors together form a basis of $\mathbb C^n$. Replace each pair $\mathbf x_j,\bar{\mathbf x}_j$ by $\mathbf u_j=(\mathbf x_j+\bar{\mathbf x}_j)/2$ and $\mathbf v_j=(\mathbf x_j-\bar{\mathbf x}_j)/(2i)$, the real and imaginary parts of $\mathbf x_j$. Since $\mathbf x_j=\mathbf u_j+i\mathbf v_j$ and $\bar{\mathbf x}_j=\mathbf u_j-i\mathbf v_j$, the complex span does not change, so the $n$ resulting real vectors still span $\mathbb C^n$ and are a basis of $\mathbb C^n$, and in particular they are independent over $\mathbb R$ and form a basis of $\mathbb R^n$. Let $S$ be the real matrix with these vectors as columns. The real eigenvectors give the diagonal part, and by \cref{prop:real-2x2-form} each pair $\mathbf u_j,\mathbf v_j$, coming from an eigenvector for $a-ib$, gives a block $C_{a,b}$.
\end{proof}

The theorem says that a real diagonalisable matrix is, in suitable real coordinates, a collection of independent scalings along lines and rotation-scalings in planes. Without the assumption of diagonalisability there is still a real block triangular form with these blocks on the diagonal, the real Schur form, which Chapter~22 and Volume~VIII develop and on which the computation of eigenvalues of real matrices is based.

\subsection{Unitary and Orthogonal Matrices}

\Cref{sec:unitary-matrices} found that a rotation of space has an axis, a vector that it fixes, and that an orthogonal matrix with determinant $-1$ reverses some vector. In the language of eigenvalues, $1$ and $-1$ are eigenvalues in these cases, and all eigenvalues of such matrices lie on the unit circle.

\begin{proposition}[Eigenvalues of unitary matrices]
\label{prop:unitary-eigen}
Let $U\in\mathbb C^{n\times n}$ be unitary, for instance a real orthogonal matrix.
\begin{enumerate}
\item Every eigenvalue of $U$ has modulus $1$.
\item Eigenvectors of $U$ for different eigenvalues are orthogonal.
\item If $U=Q$ is real orthogonal, its eigenvalues are $1$, $-1$ and pairs $e^{\pm i\varphi}$, and $\det Q=(-1)^{\operatorname{am}(-1)}$. If moreover $n$ is odd and $\det Q=1$, then $1$ is an eigenvalue.
\end{enumerate}
\end{proposition}

\begin{proof}
(i) If $U\mathbf x=\lambda\mathbf x$ with $\mathbf x\ne\mathbf 0$, then $\|\mathbf x\|=\|U\mathbf x\|=|\lambda|\,\|\mathbf x\|$ (\cref{prop:hermitian-unitary}\,(ii)). (ii) If also $U\mathbf y=\mu\mathbf y$, then $\langle\mathbf x,\mathbf y\rangle=\langle U\mathbf x,U\mathbf y\rangle=\bar\lambda\mu\langle\mathbf x,\mathbf y\rangle$, and $\bar\lambda=1/\lambda$ by (i), so $\bar\lambda\mu=\mu/\lambda\ne1$ when $\mu\ne\lambda$; hence $\langle\mathbf x,\mathbf y\rangle=0$. (iii) The real eigenvalues have modulus $1$ by (i), and the nonreal ones come in conjugate pairs $e^{\pm i\varphi}$ (\cref{thm:eigenvalues-exist}), whose products are $1$; by \cref{thm:trace-det-eigen}, $\det Q$ is the product of the real eigenvalues, which is $(-1)^{\operatorname{am}(-1)}$. If $n$ is odd, the number of real eigenvalues is odd; if none of them were $1$, all would be $-1$, and their product would be $-1\ne\det Q$.
\end{proof}

Part (iii) is \cref{prop:rotation-axis}\,(i) again, now as a statement about eigenvalues, and (ii) of that proposition says that $-1$ is an eigenvalue when $\det Q=-1$, which also follows from $\det Q=(-1)^{\operatorname{am}(-1)}$. For rotations of space the eigenvalues determine the angle.

\begin{proposition}[Eigenvalues of rotations of space]
\label{prop:rotation-eigen}
Let $Q\in\mathbb R^{3\times3}$ be orthogonal with $\det Q=1$, the rotation by $\varphi$ about the axis spanned by the unit vector $\mathbf w$ (\cref{prop:rotation-axis}\,(iii)). Then the eigenvalues of $Q$ are $1$, $e^{i\varphi}$ and $e^{-i\varphi}$, so that $\operatorname{tr}Q=1+2\cos\varphi$. If $\varphi\notin2\pi\mathbb Z$, that is, if $Q\ne I$, the eigenspace $E_1(Q)$ is the axis $\operatorname{span}\{\mathbf w\}$.
\end{proposition}

\begin{proof}
By \cref{prop:rotation-axis}\,(iii), $U^{\mathsf T}QU$ is block diagonal with the blocks $1$ and $R_\varphi$ for an orthogonal $U$ whose first column is $\mathbf w$. By \cref{prop:charpoly-similar} and \cref{thm:det-block-triangular}, $p_Q(x)=(x-1)p_{R_\varphi}(x)$, whose roots are $1$ and $e^{\pm i\varphi}$ (\cref{sec:characteristic-polynomial}); the trace is their sum, by \cref{thm:trace-det-eigen}. If $\varphi\notin2\pi\mathbb Z$, then $e^{\pm i\varphi}\ne1$, the eigenvalue $1$ has algebraic multiplicity $1$, hence geometric multiplicity $1$ by \cref{thm:gm-am}, and $E_1(Q)$ is the line spanned by $\mathbf w$, which $Q$ fixes.
\end{proof}

The cyclic permutation matrix that maps $\mathbf e_1\mapsto\mathbf e_2\mapsto\mathbf e_3\mapsto\mathbf e_1$, discussed after \cref{prop:rotation-axis}, has $p(x)=x^3-1$, as \cref{sec:structured-eigen} shows, so its eigenvalues are the cube roots of unity $1$, $e^{2\pi i/3}$ and $e^{-2\pi i/3}$, and the angle $2\pi/3$ found there appears as the argument of the nonreal eigenvalues; its axis is spanned by $(1,1,1)$. A rotation by $\pi$ has the eigenvalues $1$, $-1$ and $-1$: it reverses every vector of the plane perpendicular to its axis.

\subsection{Hermitian and Skew-Hermitian Matrices}

\begin{proposition}[Eigenvalues of Hermitian matrices]
\label{prop:hermitian-eigen}
Let $A\in\mathbb C^{n\times n}$ be Hermitian, $A^{\mathsf H}=A$.
\begin{enumerate}
\item All eigenvalues of $A$ are real.
\item Eigenvectors of $A$ for different eigenvalues are orthogonal.
\item If $A$ is real and symmetric, all its eigenvalues are real and have real eigenvectors, and real eigenvectors for different eigenvalues are orthogonal.
\item If $A$ is skew-Hermitian, $A^{\mathsf H}=-A$, for instance real and antisymmetric, all its eigenvalues lie on the imaginary axis.
\end{enumerate}
\end{proposition}

\begin{proof}
(i) If $A\mathbf x=\lambda\mathbf x$ with $\mathbf x\ne\mathbf 0$, then by \cref{prop:hermitian-inner-product}\,(v)
\[
\lambda\|\mathbf x\|^2=\langle\mathbf x,A\mathbf x\rangle=\langle A^{\mathsf H}\mathbf x,\mathbf x\rangle=\langle A\mathbf x,\mathbf x\rangle=\bar\lambda\|\mathbf x\|^2 ,
\]
so $\lambda=\bar\lambda$. (ii) If moreover $A\mathbf y=\mu\mathbf y$ with $\mu\ne\lambda$, then $\mu\langle\mathbf x,\mathbf y\rangle=\langle\mathbf x,A\mathbf y\rangle=\langle A\mathbf x,\mathbf y\rangle=\bar\lambda\langle\mathbf x,\mathbf y\rangle=\lambda\langle\mathbf x,\mathbf y\rangle$, by (i), so $\langle\mathbf x,\mathbf y\rangle=0$. (iii) A real symmetric matrix is Hermitian, so (i) and (ii) apply, and real eigenvalues of real matrices have real eigenvectors (\cref{thm:eigenvalues-exist}\,(iii)). (iv) The matrix $iA$ is Hermitian, since $(iA)^{\mathsf H}=-iA^{\mathsf H}=iA$, and its eigenvalues are $i$ times those of $A$; by (i) they are real, so those of $A$ are imaginary.
\end{proof}

Part (iii) proves the first assertion of \cref{thm:definiteness-criteria}\,(i), that the eigenvalues of a real symmetric matrix are real. The matrix with the rows $(2,i)$ and $(-i,3)$ of \cref{sec:complex-matrices} has the trace $5$ and the determinant $6-i(-i)=5$, so its eigenvalues are $(5\pm\sqrt5)/2$, real as they must be. The rotation $R_{\pi/2}$, with the rows $(0,-1)$ and $(1,0)$, is antisymmetric and orthogonal, and its eigenvalues $\pm i$ lie both on the imaginary axis and on the unit circle. The antisymmetric matrix $W$ with $W\mathbf x=\mathbf w\times\mathbf x$ for a fixed nonzero $\mathbf w\in\mathbb R^3$ has the trace $0$, the principal minors $w_3^2$, $w_2^2$ and $w_1^2$ of order $2$, and the determinant $0$, so $p_W(\lambda)=\lambda^3+\|\mathbf w\|^2\lambda$, with the eigenvalues $0$ and $\pm i\|\mathbf w\|$; the eigenspace of $0$ is the line spanned by $\mathbf w$, the vectors parallel to $\mathbf w$.

If a Hermitian matrix is diagonalisable, it has an orthonormal basis of eigenvectors: the eigenspaces of different eigenvalues are orthogonal by (ii), and within each eigenspace the Gram--Schmidt process (\cref{thm:gram-schmidt}) produces an orthonormal basis, whose vectors are still eigenvectors. With these as columns, $A=Q\Lambda Q^{\mathsf H}$ with a unitary $Q$ and a real diagonal $\Lambda$. The spectral theorem of \cref{ch:symmetric-matrices} shows that every Hermitian matrix is diagonalisable, so that this factorisation always exists; it is the central theorem about symmetric matrices. Hermitian, skew-Hermitian and unitary matrices are normal, and \cref{thm:normal-unitary} shows that the normal matrices are exactly the matrices with an orthonormal basis of eigenvectors.

\section{Structured Examples}
\label{sec:structured-eigen}

For matrices with structure, the eigenvalues and eigenvectors can often be found exactly. The results of this section are collected in \cref{tab:structured-eigen} at its end.

\subsection{The Second Difference Matrix}

\Cref{sec:special-matrices} observed that the second difference matrix $T_n$, with $2$ on the diagonal and $-1$ next to it, maps sine vectors to multiples of themselves. The next theorem gives all its eigenvalues and eigenvectors.

\begin{theorem}[Eigenvalues of the second difference matrix]
\label{thm:second-difference-eigen}
Let $n\ge1$ and, for $j=1,\ldots,n$, let $\theta_j\coloneqq j\pi/(n+1)$ and let $\mathbf s_j\in\mathbb R^n$ have the components $\sin(i\theta_j)$, $i=1,\ldots,n$. Then
\[
T_n\mathbf s_j=\lambda_j\mathbf s_j,\qquad\lambda_j=2-2\cos\theta_j=4\sin^2\frac{j\pi}{2(n+1)} .
\]
The eigenvalues satisfy $0<\lambda_1<\lambda_2<\cdots<\lambda_n<4$, the eigenvectors are orthogonal, and $\|\mathbf s_j\|^2=(n+1)/2$. Consequently $T_n=Q\Lambda Q$ with $\Lambda=\operatorname{diag}(\lambda_1,\ldots,\lambda_n)$ and the symmetric orthogonal matrix $Q$ with the entries
\[
q_{ij}=\sqrt{\frac2{n+1}}\,\sin\frac{ij\pi}{n+1}\qquad(1\le i,j\le n).
\]
\end{theorem}

\begin{proof}
Write $s_i\coloneqq\sin(i\theta_j)$ for $i=0,\ldots,n+1$; then $s_0=\sin0=0$ and $s_{n+1}=\sin(j\pi)=0$. The $i$th component of $T_n\mathbf s_j$ is $-s_{i-1}+2s_i-s_{i+1}$, also for $i=1$ and $i=n$, because the boundary values vanish. By the addition formulas (\cref{thm:addition}), $\sin((i-1)\theta)+\sin((i+1)\theta)=2\sin(i\theta)\cos\theta$, so the component is $(2-2\cos\theta_j)s_i$. The vector $\mathbf s_j$ is not zero, since $s_1=\sin\theta_j>0$ for $0<\theta_j<\pi$. The cosine decreases strictly on $[0,\pi]$ (\cref{thm:trig-signs}), and $0<\theta_1<\cdots<\theta_n<\pi$, so $\cos\theta_j$ decreases strictly from a value below $1$ to a value above $-1$, which gives the inequalities; the double angle formula $\cos\theta=1-2\sin^2(\theta/2)$ gives the second expression for $\lambda_j$. $T_n$ is symmetric, so eigenvectors for different eigenvalues are orthogonal (\cref{prop:hermitian-eigen}). For the length, $\sin^2x=(1-\cos2x)/2$, and since $s_0=s_{n+1}=0$,
\[
\|\mathbf s_j\|^2=\sum_{i=0}^{n}\sin^2(i\theta_j)=\frac{n+1}2-\frac12\sum_{i=0}^{n}\cos\frac{2ij\pi}{n+1}=\frac{n+1}2,
\]
because the last sum is the real part of $\sum_{i=0}^n\omega^{ij}$ with $\omega=e^{2\pi i/(n+1)}$, which vanishes by \cref{prop:roots-of-unity}, as $n+1$ does not divide $j$. So the columns $\mathbf s_j/\|\mathbf s_j\|$ of $Q$ are orthonormal, $Q$ is orthogonal, and it is symmetric because $q_{ij}=q_{ji}$. Then $T_nQ=Q\Lambda$, and $Q^{-1}=Q^{\mathsf T}=Q$.
\end{proof}

The theorem has several consequences that later volumes use. The smallest and the largest eigenvalues are
\[
\lambda_1=4\sin^2\frac{\pi}{2(n+1)}\approx\frac{\pi^2}{(n+1)^2},\qquad\lambda_n=4\cos^2\frac{\pi}{2(n+1)}\approx4,
\]
so their ratio grows like $4(n+1)^2/\pi^2$; this ratio is the condition number $\kappa_2(T_n)$ of \cref{ch:symmetric-matrices} and \cref{ch:svd}, and it governs the accuracy and the speed of methods for the systems $T_n\mathbf u=h^2\mathbf f$ of \cref{sec:special-matrices}. The eigenvectors are samples of the functions $\sin(j\pi x)$ at the grid points $x_i=ih$, $h=1/(n+1)$, and these functions are the eigenfunctions of the operator $u\mapsto-u''$ on $[0,1]$ with $u(0)=u(1)=0$, with the eigenvalues $(j\pi)^2$. The matrix $T_n/h^2$, which approximates this operator, has the eigenvalues
\[
\frac{2-2\cos(j\pi h)}{h^2}=(j\pi)^2\Bigl(1-\frac{(j\pi h)^2}{12}+\cdots\Bigr),
\]
by the cosine series (\cref{thm:sin-cos-series}). The smooth modes, with $jh$ small, are reproduced accurately, while the oscillating modes, with $j$ near $n$, have eigenvalues near $4/h^2$ instead of $(j\pi)^2\approx\pi^2/h^2$. \Cref{fig:second-difference} shows both effects. The orthogonal matrix $Q$ is the matrix of the \term{discrete sine transform}, and the theorem says that this transform diagonalises $T_n$: in the coordinates of the sine basis, solving $T_n\mathbf u=\mathbf b$ amounts to dividing by the $\lambda_j$. Volume~XI computes the transform with the fast Fourier transform, and Strang~\cite{strang1999dct} explains the family of such transforms through the eigenvectors of second difference matrices with other boundary conditions.

\begin{figure}
\sdpanel{A}{\begin{tikzpicture}
\begin{axis}[sdplot, width=0.47\textwidth, height=0.4\textwidth, xmin=0, xmax=1, ymin=-0.5, ymax=0.62,
  xtick={0,0.5,1}, ytick={-0.4,0,0.4}, xlabel={$x$}, legend style={at={(0.5,1.0)}, anchor=south}, legend columns=-1]
\addplot[black, line width=0.6pt, domain=0:1, samples=80, forget plot] {sqrt(2/11)*sin(deg(pi*x))};
\addplot[black, dashed, line width=0.6pt, domain=0:1, samples=80, forget plot] {sqrt(2/11)*sin(deg(2*pi*x))};
\addplot[black!50, line width=0.6pt, domain=0:1, samples=120, forget plot] {sqrt(2/11)*sin(deg(3*pi*x))};
\addplot[black, only marks, mark=*, mark size=1.4pt, samples at={0,1,...,11}] ({x/11},{sqrt(2/11)*sin(deg(pi*x/11))});
\addlegendentry{$j=1$}
\addplot[black, only marks, mark=square*, mark size=1.3pt, samples at={0,1,...,11}] ({x/11},{sqrt(2/11)*sin(deg(2*pi*x/11))});
\addlegendentry{$j=2$}
\addplot[black!50, only marks, mark=triangle*, mark size=1.6pt, samples at={0,1,...,11}] ({x/11},{sqrt(2/11)*sin(deg(3*pi*x/11))});
\addlegendentry{$j=3$}
\end{axis}
\end{tikzpicture}}\hfill
\sdpanel{B}{\begin{tikzpicture}
\begin{axis}[sdplot, width=0.47\textwidth, height=0.4\textwidth, xmin=0, xmax=1, ymin=0, ymax=4.4,
  xtick={0,0.5,1}, ytick={0,1,2,3,4}, xlabel={$j/(n+1)$}, legend pos=north west]
\addplot[black, line width=0.6pt, domain=0:1, samples=80] {2-2*cos(deg(pi*x))};
\addlegendentry{$2-2\cos(\pi t)$}
\addplot[black, dashed, line width=0.6pt, domain=0:0.668, samples=60] {(pi*x)^2};
\addlegendentry{$(\pi t)^2$}
\addplot[black, only marks, mark=*, mark size=1.4pt, samples at={1,2,...,10}, forget plot] ({x/11},{2-2*cos(deg(pi*x/11))});
\end{axis}
\end{tikzpicture}}
\caption{Eigenvectors and eigenvalues of the second difference matrix $T_{10}$ (\cref{thm:second-difference-eigen}). (A) The components of the normalised eigenvectors for $j=1,2,3$, plotted at the grid points $x_i=i/11$ together with the boundary values $0$, lie on the curves $\sqrt{2/11}\sin(j\pi x)$. (B) The eigenvalues $\lambda_j$ (dots) on the curve $2-2\cos(\pi t)$, $t=j/(n+1)$, and the eigenvalues $(j\pi h)^2$ of $-u''$, scaled by $h^2=1/121$ (dashed); they agree for small $j$ and part for large $j$.\label{fig:second-difference}}
\end{figure}

\subsection{Circulant Matrices}

A \term{circulant matrix} is constant along each of its diagonals, wrapping around: for $\mathbf c=(c_0,\ldots,c_{n-1})\in\mathbb C^n$, with rows and columns numbered from $0$, the circulant $C$ with the first column $\mathbf c$ has the entries $C_{lm}\coloneqq c_{(l-m)\bmod n}$. For $n=4$,
\[
C=\begin{bmatrix}c_0&c_3&c_2&c_1\\c_1&c_0&c_3&c_2\\c_2&c_1&c_0&c_3\\c_3&c_2&c_1&c_0\end{bmatrix}.
\]
Each column is the one before it shifted down by one position, with the last entry moved to the top. The basic circulant is the \term{cyclic shift} $Z$ with the first column $(0,1,0,\ldots,0)$; with the standard basis vectors numbered $\mathbf e_0,\ldots,\mathbf e_{n-1}$ like the components, it maps $\mathbf e_m$ to $\mathbf e_{(m+1)\bmod n}$; its powers $Z^m$ have the entry $1$ in the positions $(l,k)$ with $l\equiv k+m\pmod n$, so every circulant is a polynomial in $Z$,
\[
C=c_0I+c_1Z+c_2Z^2+\cdots+c_{n-1}Z^{n-1}.
\]
The Fourier vectors $\mathbf f_k$ of \cref{prop:fourier-basis}, with the components $\omega^{lk}/\sqrt n$, $\omega=e^{2\pi i/n}$, are eigenvectors of every circulant.

\begin{theorem}[Eigenvalues of circulant matrices]
\label{thm:circulant}
Let $C$ be the circulant with the first column $\mathbf c$, and let $F$ be the unitary matrix with the columns $\mathbf f_0,\ldots,\mathbf f_{n-1}$, the matrix $F_n/\sqrt n$ of \cref{sec:complex-matrices}. Then
\[
C\mathbf f_k=\gamma_k\mathbf f_k,\qquad\gamma_k\coloneqq\sum_{m=0}^{n-1}c_m\omega^{-mk}\qquad(k=0,\ldots,n-1),
\]
so that $C=F\operatorname{diag}(\gamma_0,\ldots,\gamma_{n-1})F^{\mathsf H}$. Conversely, every matrix $F\operatorname{diag}(d_0,\ldots,d_{n-1})F^{\mathsf H}$ is a circulant. In particular, all circulants of order $n$ commute, and sums, products and inverses of circulants are circulants.
\end{theorem}

\begin{proof}
The shift moves the components down cyclically, $(Z\mathbf x)_l=x_{(l-1)\bmod n}$, so $(Z\mathbf f_k)_l=\omega^{(l-1)k}/\sqrt n=\omega^{-k}(\mathbf f_k)_l$, using $\omega^n=1$ for $l=0$: $Z\mathbf f_k=\omega^{-k}\mathbf f_k$. By \cref{prop:eigen-basic}\,(iv), $C\mathbf f_k=\sum_mc_mZ^m\mathbf f_k=\sum_mc_m\omega^{-mk}\mathbf f_k$. The matrix $F$ is unitary by \cref{prop:fourier-basis}, so $CF=F\operatorname{diag}(\gamma_k)$ gives the factorisation. For the converse, the numbers $\omega^{-k}$, $k=0,\ldots,n-1$, are pairwise different, so there is a polynomial $q$ of degree less than $n$ with $q(\omega^{-k})=d_k$ (\cref{thm:interpolation}), and then $F\operatorname{diag}(d_k)F^{\mathsf H}=Fq(\operatorname{diag}(\omega^{-k}))F^{\mathsf H}=q(Z)$, a polynomial in $Z$, hence a circulant. Diagonal matrices commute, and their sums, products and inverses are diagonal, which gives the last statement.
\end{proof}

In the notation of \cref{prop:fourier-basis}, $\gamma_k=\sqrt n\,\hat c_k$: the eigenvalues of a circulant are, up to the factor $\sqrt n$, the discrete Fourier coefficients of its first column. Multiplying by a circulant is therefore the same as transforming to Fourier coefficients, multiplying each by $\gamma_k$ and transforming back, and with the fast Fourier transform of Volume~XI each of these steps costs a multiple of $n\log n$ operations instead of $n^2$. The cyclic shift itself has the eigenvalues $\omega^{-k}$, all $n$th roots of unity; it is the companion matrix of $x^n-1$ in the sense of \cref{prop:companion}, so $p_Z(x)=x^n-1$. The periodic second difference matrix, the circulant with the first column $(2,-1,0,\ldots,0,-1)$, which arises from $-u''$ on a circle, has the eigenvalues
\[
\gamma_k=2-\omega^{-k}-\omega^{k}=2-2\cos\frac{2\pi k}n\qquad(k=0,\ldots,n-1),
\]
which include $\gamma_0=0$, for the constant vector, and in which $\gamma_k=\gamma_{n-k}$, so that most eigenvalues are double, with the eigenvectors $\mathbf f_k$ and $\mathbf f_{n-k}$, or equivalently the real cosine and sine vectors.

\subsection{Permutation Matrices}

\begin{proposition}[Eigenvalues of permutation matrices]
\label{prop:permutation-eigen}
Let $\sigma$ be a permutation of $\{1,\ldots,n\}$ whose cycles, including the fixed points as cycles of length $1$, have the lengths $\ell_1,\ldots,\ell_r$. Then
\[
p_{P_\sigma}(x)=(x^{\ell_1}-1)(x^{\ell_2}-1)\cdots(x^{\ell_r}-1),
\]
so the eigenvalues of $P_\sigma$ are roots of unity, the eigenvalue $1$ has algebraic multiplicity $r$, and $P_\sigma$ is diagonalisable over $\mathbb C$.
\end{proposition}

\begin{proof}
Number the elements anew so that each cycle consists of consecutive numbers, in the order in which the cycle visits them: if $\tau$ is this renumbering, then $P_\tau P_\sigma P_\tau^{-1}=P_{\tau\sigma\tau^{-1}}$ (\cref{prop:permutation-matrices}), and $\tau\sigma\tau^{-1}$ maps each block of consecutive numbers cyclically onto itself. Its matrix is block diagonal, with the cyclic shift of order $\ell_i$ as the $i$th block, whose characteristic polynomial is $x^{\ell_i}-1$. Similar matrices have the same characteristic polynomial, and that of a block diagonal matrix is the product of those of the blocks (\cref{thm:det-block-triangular}). Each factor $x^{\ell}-1$ has the root $1$ exactly once, by \cref{prop:roots-of-unity}. Finally $\sigma^m=\operatorname{id}$ for some $m\ge1$, since the permutations of a finite set form a finite group, so $P_\sigma^m=I$, and the matrix is diagonalisable by \cref{thm:diagonalisable-minimal} and the remark after it.
\end{proof}

The trace of $P_\sigma$ is the number of fixed points of $\sigma$, and by \cref{thm:trace-det-eigen} it is the sum of the eigenvalues; for a single cycle of length $n$ the eigenvalues are all the $n$th roots of unity, whose sum is $0$. The determinant is the product of the eigenvalues; since $x^\ell-1$ has the roots $\omega^k$ with product $(-1)^{\ell-1}$, it is $\prod_i(-1)^{\ell_i-1}$, and since $\det P_\sigma=\operatorname{sgn}\sigma$ (\cref{prop:det-triangular}\,(ii)), this is a formula for the sign of a permutation in terms of its cycles.

\subsection{Rank-One Matrices and Block Triangular Matrices}

The characteristic polynomial of a rank-one matrix was found after \cref{thm:ab-ba}. The eigenvectors complete the picture.

\begin{proposition}[Rank-one matrices]
\label{prop:rank-one-eigen}
Let $\mathbf u,\mathbf v\in\mathbb C^n$ be nonzero and $A\coloneqq\mathbf u\mathbf v^{\mathsf H}$. Then $p_A(\lambda)=\lambda^{n-1}(\lambda-\mathbf v^{\mathsf H}\mathbf u)$, and $E_0(A)=\{\mathbf x:\mathbf v^{\mathsf H}\mathbf x=0\}$ has dimension $n-1$. If $\mathbf v^{\mathsf H}\mathbf u\ne0$, then $\mathbf u$ is an eigenvector for $\mathbf v^{\mathsf H}\mathbf u$ and $A$ is diagonalisable; if $\mathbf v^{\mathsf H}\mathbf u=0$, then $A^2=0$ and $A$ is not diagonalisable.
\end{proposition}

\begin{proof}
The characteristic polynomial is \cref{thm:ab-ba} with the $n\times1$ matrix $\mathbf u$ and the $1\times n$ matrix $\mathbf v^{\mathsf H}$. Since $A\mathbf x=\mathbf u(\mathbf v^{\mathsf H}\mathbf x)$ and $\mathbf u\ne\mathbf 0$, $A\mathbf x=\mathbf 0$ if and only if $\mathbf v^{\mathsf H}\mathbf x=0$, a single nontrivial linear equation, whose solutions form a subspace of dimension $n-1$. If $\mathbf v^{\mathsf H}\mathbf u\ne0$, then $A\mathbf u=(\mathbf v^{\mathsf H}\mathbf u)\mathbf u$, and the geometric multiplicities add up to $(n-1)+1=n$. If $\mathbf v^{\mathsf H}\mathbf u=0$, then $A^2=\mathbf u(\mathbf v^{\mathsf H}\mathbf u)\mathbf v^{\mathsf H}=0$, the only eigenvalue is $0$, and its geometric multiplicity $n-1$ is less than $n$.
\end{proof}

By the spectral mapping theorem with $q(x)=a+x$, the matrix $aI+b\mathbf 1\mathbf 1^{\mathsf T}$ of \cref{sec:det-block}, with $a$ on the diagonal raised by $b$ and $b$ elsewhere, has the eigenvalues $a+nb$, for the eigenvector $\mathbf 1=(1,\ldots,1)$, and $a$, with multiplicity $n-1$, for the vectors whose components add up to $0$; the product $a^{n-1}(a+nb)$ is the determinant found there. The Householder reflection $I-2\mathbf u\mathbf u^{\mathsf T}$ with a unit vector $\mathbf u$ has the eigenvalue $-1$ for $\mathbf u$ and $1$, with multiplicity $n-1$, on the hyperplane $\mathbf u^\perp$.

\begin{proposition}[Block triangular matrices]
\label{prop:block-triangular-eigen}
Let
\[
M=\begin{bmatrix}A&B\\0&D\end{bmatrix}
\]
with square blocks $A$ and $D$. Then $p_M=p_Ap_D$: the eigenvalues of $M$, counted with multiplicity, are those of $A$ together with those of $D$. If $A\mathbf x=\lambda\mathbf x$, then $(\mathbf x,\mathbf 0)$ is an eigenvector of $M$.
\end{proposition}

\begin{proof}
$\lambda I-M$ is block triangular with the diagonal blocks $\lambda I-A$ and $\lambda I-D$, and \cref{thm:det-block-triangular} gives the first statement; the second follows by multiplying out.
\end{proof}

By induction the same holds for block triangular matrices with any number of square diagonal blocks, and \cref{prop:triangular-eigen} is the case of blocks of order $1$. Together with \cref{prop:block-triangular}, the proposition says that if an operator has an invariant subspace, its eigenvalues are those of its restriction to the subspace together with those of the lower right block, which Chapter~22 uses repeatedly.

\subsection{Graphs}

The \term{adjacency matrix} of a graph with the vertices $1,\ldots,n$ (\cref{sec:graphs}) has the entry $a_{ij}=1$ if $\{i,j\}$ is an edge and $0$ otherwise. It is real and symmetric, so its eigenvalues are real, and its powers count walks: $(A^k)_{ij}$ is the number of walks of length $k$ from $i$ to $j$, by induction on $k$, since a walk of length $k+1$ from $i$ to $j$ is a walk of length $k$ from $i$ to some $l$ followed by an edge $\{l,j\}$, and $(A^{k+1})_{ij}=\sum_l(A^k)_{il}a_{lj}$. The row sums of $A$ are the degrees of the vertices, so $\rho(A)\le\|A\|_\infty$ is at most the largest degree (\cref{prop:spectral-radius-norm}), with equality for a regular graph, for which $\mathbf 1$ is an eigenvector. The path with $n$ vertices has the adjacency matrix $2I-T_n$, with the eigenvalues $2\cos(j\pi/(n+1))$ by \cref{thm:second-difference-eigen,thm:spectral-mapping}, and the cycle with $n$ vertices has the circulant adjacency matrix $Z+Z^{-1}$, with the eigenvalues $2\cos(2\pi k/n)$ by \cref{thm:circulant}. The eigenvalues of graphs carry a surprising amount of information about their structure, connectivity and expansion, a subject called spectral graph theory~\cite{biggs1993,godsil2001} that Volume~IX takes up together with the computational side of graphs.

\begin{table}
\sdfloat{%
\begin{sdtable}{colspec={lll}, row{1}={font=\small\bfseries}}
\toprule
matrix & eigenvalues & eigenvectors \\
\midrule
$T_n$ (second difference) & $2-2\cos(j\pi/(n+1))$, $j=1,\ldots,n$ & $\bigl(\sin(ij\pi/(n+1))\bigr)_{i=1}^n$ \\
circulant, first column $\mathbf c$ & $\gamma_k=\sum_mc_m\omega^{-mk}$ & Fourier vectors $\mathbf f_k$ \\
cyclic shift $Z$ & $\omega^{-k}$, all $n$th roots of unity & $\mathbf f_k$ \\
permutation $P_\sigma$ & roots of $\prod_i(x^{\ell_i}-1)$ & from the cycles \\
$\mathbf u\mathbf v^{\mathsf H}$ & $\mathbf v^{\mathsf H}\mathbf u$; $0$ ($n-1$ times) & $\mathbf u$; $\mathbf v^\perp$ \\
$aI+b\mathbf 1\mathbf 1^{\mathsf T}$ & $a+nb$; $a$ ($n-1$ times) & $\mathbf 1$; $\mathbf 1^\perp$ \\
Householder $I-2\mathbf u\mathbf u^{\mathsf T}$ & $-1$; $1$ ($n-1$ times) & $\mathbf u$; $\mathbf u^\perp$ \\
block triangular & those of the diagonal blocks & \\
path, cycle (adjacency) & $2\cos(j\pi/(n+1))$; $2\cos(2\pi k/n)$ & as for $T_n$; $\mathbf f_k$ \\
\bottomrule
\end{sdtable}}{Eigenvalues and eigenvectors of structured matrices (\cref{sec:structured-eigen}); $\omega=e^{2\pi i/n}$, and $\ell_1,\ldots,\ell_r$ are the lengths of the cycles of $\sigma$.\label{tab:structured-eigen}}
\end{table}

\section{Markov Matrices and Gershgorin's Theorem}
\label{sec:markov}

\subsection{Markov Matrices}

A system that is in one of $n$ states at each time step, and moves from state $j$ to state $i$ with a fixed probability $m_{ij}$, is a \term{Markov chain}. Volume~III develops the probability theory; here we need only the matrix. If $\mathbf p_k$ is the vector of the probabilities of the $n$ states at time $k$, then the probability of state $i$ at time $k+1$ is $\sum_jm_{ij}(\mathbf p_k)_j$, so that $\mathbf p_{k+1}=M\mathbf p_k$ and $\mathbf p_k=M^k\mathbf p_0$.

\begin{definition}[Markov matrix]
\label{def:markov}
A real matrix $M\in\mathbb R^{n\times n}$ is a \term{Markov matrix} if its entries are nonnegative and each of its columns sums to $1$, that is, $\mathbf 1^{\mathsf T}M=\mathbf 1^{\mathsf T}$ for $\mathbf 1=(1,\ldots,1)$. A vector $\mathbf p\in\mathbb R^n$ with $p_i\ge0$ and $\sum_ip_i=1$ is a \term{probability vector}, and a probability vector with $M\boldsymbol\pi=\boldsymbol\pi$ is a \term{stationary distribution} of $M$.
\end{definition}

Many authors write probability vectors as rows and use matrices whose rows sum to $1$; the two conventions are transposes of each other. The simplest example has two states and the matrix
\[
M=\begin{bmatrix}1-a&b\\a&1-b\end{bmatrix},\qquad0\le a,b\le1 ,
\]
in which $a$ is the probability of leaving state $1$ and $b$ that of leaving state $2$. Its trace is $2-a-b$, and one eigenvalue is $1$, as the next proposition shows, so the other is $1-a-b$, with the eigenvector $(1,-1)$. If $a+b>0$, the stationary distribution is $\boldsymbol\pi=(b,a)/(a+b)$. A probability vector $\mathbf p_0$ differs from $\boldsymbol\pi$ by a vector whose components add up to $0$, a multiple of $(1,-1)$, so
\[
\mathbf p_k=\boldsymbol\pi+(1-a-b)^k(\mathbf p_0-\boldsymbol\pi).
\]
For $a=0.1$ and $b=0.3$, $\boldsymbol\pi=(0.75,0.25)$ and the distance to $\boldsymbol\pi$ shrinks by the factor $0.6$ per step, whatever the initial state. The convergence to $\boldsymbol\pi$ from every initial state fails only in two cases: for $a=b=0$, when $M=I$ and nothing moves, and for $a=b=1$, when the chain jumps back and forth and $M$, the permutation matrix with the rows $(0,1)$ and $(1,0)$, has the eigenvalues $1$ and $-1$.

\begin{proposition}[Eigenvalues of Markov matrices]
\label{prop:markov-basic}
Let $M$ be a Markov matrix. Then $1$ is an eigenvalue of $M$, every eigenvalue satisfies $|\lambda|\le1$, $\|M\mathbf x\|_1\le\|\mathbf x\|_1$ for every $\mathbf x\in\mathbb C^n$, and $M$ maps probability vectors to probability vectors and vectors with component sum $0$ to such vectors.
\end{proposition}

\begin{proof}
$M^{\mathsf T}\mathbf 1=\mathbf 1$, so $1$ is an eigenvalue of $M^{\mathsf T}$ and therefore of $M$ (\cref{prop:charpoly-similar}). The column sums of $|m_{ij}|=m_{ij}$ are $1$, so $\|M\|_1=1$ by \cref{thm:norm-formulas}, and $\rho(M)\le1$ by \cref{prop:spectral-radius-norm}. If $p_i\ge0$, then $(M\mathbf p)_i=\sum_jm_{ij}p_j\ge0$, and $\mathbf 1^{\mathsf T}M\mathbf x=\mathbf 1^{\mathsf T}\mathbf x$ for every $\mathbf x$ shows that $M$ preserves the sum of the components.
\end{proof}

When all transition probabilities are positive, the chain forgets its initial state at a geometric rate.

\begin{theorem}[Perron's theorem for positive Markov matrices]
\label{thm:perron-positive}
Let $M$ be a Markov matrix with positive entries, let $\delta\coloneqq\min_{i,j}m_{ij}>0$ and $q\coloneqq1-n\delta$. Then $0\le q<1$, and:
\begin{enumerate}
\item $M$ has exactly one stationary distribution $\boldsymbol\pi$, and $\pi_i\ge\delta$ for all $i$;
\item $\|M^k\mathbf p_0-\boldsymbol\pi\|_1\le2q^k$ for every probability vector $\mathbf p_0$ and every $k\in\mathbb N$;
\item the eigenvalue $1$ has algebraic multiplicity $1$, and every other eigenvalue satisfies $|\lambda|\le q$.
\end{enumerate}
\end{theorem}

\begin{proof}
Each column of $M$ has $n$ entries of at least $\delta$ and sums to $1$, so $n\delta\le1$. The proof rests on a contraction: \emph{if $\mathbf x\in\mathbb C^n$ has $\sum_ix_i=0$, then $\|M\mathbf x\|_1\le q\|\mathbf x\|_1$}. Indeed, let $E\coloneqq\mathbf 1\mathbf 1^{\mathsf T}$, the matrix of ones. Then $E\mathbf x=\mathbf 1(\mathbf 1^{\mathsf T}\mathbf x)=\mathbf 0$, so $M\mathbf x=(M-\delta E)\mathbf x$, and $M-\delta E$ has nonnegative entries $m_{ij}-\delta$ and column sums $1-n\delta=q$, so that $\|M-\delta E\|_1=q$ by \cref{thm:norm-formulas}, and $\|M\mathbf x\|_1\le q\|\mathbf x\|_1$.

(i) Let $\mathbf p_0$ be any probability vector and $\mathbf p_k\coloneqq M^k\mathbf p_0$. The difference $\mathbf p_1-\mathbf p_0$ has component sum $0$, and so has $M^k(\mathbf p_1-\mathbf p_0)=\mathbf p_{k+1}-\mathbf p_k$ for every $k$, by \cref{prop:markov-basic}; by the contraction, $\|\mathbf p_{k+1}-\mathbf p_k\|_1\le q^k\|\mathbf p_1-\mathbf p_0\|_1\le2q^k$. For $l>k$, the triangle inequality gives $\|\mathbf p_l-\mathbf p_k\|_1\le2\sum_{j\ge k}q^j=2q^k/(1-q)$, so each component of $\mathbf p_k$ is a Cauchy sequence and converges (\cref{thm:cauchy-criterion}). The limit $\boldsymbol\pi$ has nonnegative components with sum $1$, and $M\boldsymbol\pi=\lim M\mathbf p_k=\lim\mathbf p_{k+1}=\boldsymbol\pi$, since each component of $M\mathbf p_k$ is a fixed combination of components of $\mathbf p_k$. Moreover $\pi_i=\sum_jm_{ij}\pi_j\ge\delta\sum_j\pi_j=\delta$. If $\boldsymbol\pi'$ is another stationary distribution, then $\boldsymbol\pi-\boldsymbol\pi'$ has component sum $0$, and $\|\boldsymbol\pi-\boldsymbol\pi'\|_1=\|M(\boldsymbol\pi-\boldsymbol\pi')\|_1\le q\|\boldsymbol\pi-\boldsymbol\pi'\|_1$, which forces $\boldsymbol\pi'=\boldsymbol\pi$.

(ii) $\mathbf p_0-\boldsymbol\pi$ has component sum $0$ and $\|\mathbf p_0-\boldsymbol\pi\|_1\le\|\mathbf p_0\|_1+\|\boldsymbol\pi\|_1=2$, so $\|M^k\mathbf p_0-\boldsymbol\pi\|_1=\|M^k(\mathbf p_0-\boldsymbol\pi)\|_1\le2q^k$.

(iii) The subspace $W\coloneqq\{\mathbf x\in\mathbb C^n:\sum_ix_i=0\}$ has dimension $n-1$ and is invariant under $M$, and so is the line spanned by $\boldsymbol\pi$, which is not contained in $W$. So $\mathbb C^n=\operatorname{span}\{\boldsymbol\pi\}\oplus W$, and in a basis consisting of $\boldsymbol\pi$ and a basis of $W$ the matrix of $M$ is block diagonal with the blocks $1$ and the matrix $M_W$ of the restriction to $W$ (\cref{prop:block-triangular}). By \cref{prop:block-triangular-eigen}, $p_M(\lambda)=(\lambda-1)p_{M_W}(\lambda)$. An eigenvalue $\mu$ of $M_W$ has an eigenvector $\mathbf w\in W$, and the contraction gives $|\mu|\,\|\mathbf w\|_1=\|M\mathbf w\|_1\le q\|\mathbf w\|_1$, so $|\mu|\le q<1$. Hence $1$ is a simple root of $p_M$, and the other eigenvalues are those of $M_W$.
\end{proof}

The bound $q$ is crude, and the true rate is the largest modulus of the eigenvalues other than $1$. For
\[
M=\begin{bmatrix}0.5&0.3&0.2\\0.3&0.5&0.3\\0.2&0.2&0.5\end{bmatrix},
\]
whose columns sum to $1$, $\delta=0.2$ and $q=0.4$. The characteristic polynomial is $\lambda^3-1.5\lambda^2+0.56\lambda-0.06=(\lambda-1)(\lambda-0.3)(\lambda-0.2)$, and the stationary distribution, the solution of $(M-I)\boldsymbol\pi=\mathbf 0$ with component sum $1$, is $\boldsymbol\pi=(19,21,16)/56$, as one checks by multiplying. The eigenvalues are different, so $M$ is diagonalisable, and by \eqref{eq:power-expansion} the difference $M^k\mathbf p_0-\boldsymbol\pi$ is a combination of $0.3^k$ and $0.2^k$: it shrinks asymptotically by the factor $0.3$ per step, faster than the bound $0.4$ guarantees.

Positivity of the entries is essential in this proof but not in the conclusion. The conclusions hold as well if some power $M^m$ has positive entries, since the theorem applies to $M^m$; they fail for the permutation matrix with the rows $(0,1)$ and $(1,0)$, no power of which is positive. The general theory of nonnegative matrices, due to Perron and Frobenius, is summarised in the note below, and Volume~III treats Markov chains with the methods of probability.

A Markov matrix of enormous size underlies Google's \term{PageRank}~\cite{brin1998}. A random surfer on the $n$ pages of the web follows, from page $j$, one of its $d_j$ links at random, so that the matrix $L$ of this walk has the entries $1/d_j$ in column $j$ in the rows of the pages that $j$ links to; a page without links is usually given the column $\mathbf 1/n$~\cite{langville2006}. With probability $1-\alpha$ the surfer instead jumps to a page chosen uniformly at random, which gives the Markov matrix
\[
G=\alpha L+\frac{1-\alpha}n\,\mathbf 1\mathbf 1^{\mathsf T},
\]
with $\alpha=0.85$ in the original proposal. All entries of $G$ are at least $(1-\alpha)/n$, so by \cref{thm:perron-positive} it has a unique stationary distribution $\boldsymbol\pi$, whose components rank the pages, and since $q\le1-n(1-\alpha)/n=\alpha$, every eigenvalue other than $1$ has modulus at most $\alpha$, and the distance of the iterates $\mathbf p_{k+1}=G\mathbf p_k$ from $\boldsymbol\pi$ shrinks at least like $\alpha^k$, by \cref{thm:perron-positive}\,(ii). This iteration is the power method of \cref{sec:computing-eigenvalues}. The matrix $G$ is never formed, since it is full: for a probability vector, $G\mathbf p=\alpha L\mathbf p+(1-\alpha)\mathbf 1/n$, and $L$ is sparse (Volume~IX).

\begin{note}{The Perron--Frobenius theorem}
Perron proved in 1907 that a square matrix with positive entries has a positive eigenvalue equal to its spectral radius, that this eigenvalue is a simple root of the characteristic polynomial, that it has an eigenvector with positive components, and that all other eigenvalues have smaller modulus~\cite{perron1907}. Frobenius extended the theory in 1912 to nonnegative matrices~\cite{frobenius1912}. A nonnegative matrix $A$ is \emph{irreducible} if the directed graph with an edge from $j$ to $i$ whenever $a_{ij}>0$ is strongly connected, that is, if every state can be reached from every other. For an irreducible nonnegative matrix, $\rho(A)>0$ is a simple eigenvalue with a positive eigenvector, unique up to scaling, and no other eigenvalue has a nonnegative eigenvector; other eigenvalues of modulus $\rho(A)$ can occur, as for the permutation matrix of a cycle, and they are then $\rho(A)$ times roots of unity. If some power of $A$ is positive, $\rho(A)$ is strictly larger than the modulus of every other eigenvalue. Horn and Johnson~\cite[Chap.~8]{horn2013} prove all of this; for a Markov matrix, $\rho=1$ and the positive eigenvector is the stationary distribution. Nonnegative matrices arise in Markov chains, in population models and in the ranking of web pages, and for the irreducible ones the theorem is the reason why the dominant eigenvalue is real, positive and simple, with a positive eigenvector.
\end{note}

\subsection{Gershgorin's Theorem}

The eigenvalues of a triangular matrix are its diagonal entries. If the entries off the diagonal are not zero but small, the eigenvalues stay near the diagonal entries, and a theorem of Gershgorin~\cite{gershgorin1931} makes this precise without computing anything.

\begin{theorem}[Gershgorin's theorem]
\label{thm:gershgorin}
Let $A\in\mathbb C^{n\times n}$, and let $D_i$ be the closed disc in $\mathbb C$ with the centre $a_{ii}$ and the radius $R_i\coloneqq\sum_{j\ne i}|a_{ij}|$, the sum of the moduli of the other entries of row $i$. Then every eigenvalue of $A$ lies in $D_1\cup D_2\cup\cdots\cup D_n$. The same holds for the discs whose radii are the sums over the other entries of the columns.
\end{theorem}

\begin{proof}
Let $A\mathbf x=\lambda\mathbf x$ with $\mathbf x\ne\mathbf 0$, and choose $i$ with $|x_i|=\|\mathbf x\|_\infty>0$. The $i$th component of the equation is $(\lambda-a_{ii})x_i=\sum_{j\ne i}a_{ij}x_j$, so
\[
|\lambda-a_{ii}|\,|x_i|\le\sum_{j\ne i}|a_{ij}|\,|x_j|\le R_i|x_i| ,
\]
and $\lambda\in D_i$. For the columns, apply this to $A^{\mathsf T}$, which has the same eigenvalues.
\end{proof}

The proof is a refinement of $\rho(A)\le\|A\|_\infty$: it uses the largest component of the eigenvector and keeps the diagonal entry apart. A matrix that is strictly diagonally dominant by rows (\cref{def:diagonal-dominance}), $|a_{ii}|>R_i$ for every $i$, has no disc containing $0$, so $0$ is not an eigenvalue and the matrix is invertible. This is the invertibility in \cref{thm:diagonal-dominance}, now for complex matrices and with a proof that does not use elimination; the result goes back to L\'evy and Desplanques at the end of the nineteenth century and has been rediscovered many times, as Taussky~\cite{taussky1949} recounts. When the discs fall into separate groups, the theorem says more.

\begin{theorem}[Separated discs]
\label{thm:gershgorin-separated}
If the union of $k$ of the Gershgorin discs of $A$ is disjoint from the union of the other $n-k$, then it contains exactly $k$ eigenvalues of $A$, counted with multiplicity. In particular a disc that meets no other disc contains exactly one eigenvalue, and if $A$ is real, this eigenvalue is real.
\end{theorem}

\begin{proof}
Let $A_0$ be the diagonal part of $A$ and $A_t\coloneqq A_0+t(A-A_0)$ for $0\le t\le1$. The discs of $A_t$ have the same centres as those of $A$ and the radii $tR_i$, so for every $t$ the union of the first group of discs lies inside the union for $A$ and is disjoint from the union of the others. The eigenvalues of $A_t$ are the roots of $p_{A_t}$, whose coefficients are polynomials in $t$, and the roots of a polynomial depend continuously on its coefficients, a consequence of the argument principle of complex analysis (Volume~II). As $t$ runs from $0$ to $1$, no eigenvalue can pass from one union to the other without leaving both, which Gershgorin's theorem forbids; so the number of eigenvalues in the first union is the same for $t=1$ as for $t=0$, where the eigenvalues are the diagonal entries and exactly $k$ of them are centres of the discs of the first group. Horn and Johnson~\cite[Sect.~6.1]{horn2013} and Varga~\cite{varga2004} give the details. For a real matrix, the discs are symmetric about the real axis, so the conjugate of an eigenvalue in an isolated disc lies in the same disc, and the disc contains only one eigenvalue, which is therefore real.
\end{proof}

\Cref{fig:gershgorin} shows the discs of
\[
A=\begin{bmatrix}4&-1.5&0\\1&2&0.5\\0.5&0.5&-2\end{bmatrix}.
\]
The row discs have the centres $4$, $2$, $-2$ and the radii $1.5$, $1.5$, $1$; the first two overlap, and the third is separate. By \cref{thm:gershgorin-separated}, the third contains exactly one eigenvalue, which is real, and the union of the first two contains the other two. The characteristic polynomial is $\lambda^3-4\lambda^2-11\lambda/4+163/8$, with the roots $-2.0722$ and $3.0361\pm0.7841i$ to four decimals. The column discs have the radii $1.5$, $2$ and $0.5$, and since every eigenvalue lies in both unions, the real eigenvalue is confined to the interval $[-2.5,-1.5]$.

\begin{figure}
\sdfloat{%
\begin{tikzpicture}[scale=0.82, font=\scriptsize]
\fill[black!10] (4,0) circle (1.5);
\fill[black!10] (2,0) circle (1.5);
\fill[black!10] (-2,0) circle (1);
\draw[line width=0.6pt] (4,0) circle (1.5);
\draw[line width=0.6pt] (2,0) circle (1.5);
\draw[line width=0.6pt] (-2,0) circle (1);
\draw[dashed, line width=0.5pt] (4,0) circle (1.5);
\draw[dashed, line width=0.5pt] (2,0) circle (2);
\draw[dashed, line width=0.5pt] (-2,0) circle (0.5);
\foreach \c in {4,2,-2} {\draw[line width=0.5pt] (\c,-0.08) -- (\c,0.08);}
\fill (-2.07218,0) circle (2pt);
\fill (3.03609,0.78409) circle (2pt);
\fill (3.03609,-0.78409) circle (2pt);
\draw[line width=0.4pt, -{Latex[length=1.5mm]}] (-3.6,0) -- (6.3,0) node[below] {$\operatorname{Re}$};
\draw[line width=0.4pt, -{Latex[length=1.5mm]}] (0,-2.3) -- (0,2.4) node[left] {$\operatorname{Im}$};
\node[below right, inner sep=1pt] at (4.1,-0.04) {$4$};
\node[below] at (2.15,-0.1) {$2$};
\node at (-2,-0.75) {$-2$};
\end{tikzpicture}}{Gershgorin discs of the matrix $A$ in the text. The row discs (shaded, solid boundaries) have the centres $4$, $2$, $-2$ and the radii $1.5$, $1.5$, $1$; the column discs (dashed) have the same centres and the radii $1.5$, $2$, $0.5$, so that the two discs about $4$ coincide. The eigenvalues (dots) are $-2.0722$, in the isolated discs, and $3.0361\pm0.7841i$, in the union of the two overlapping row discs (\cref{thm:gershgorin,thm:gershgorin-separated}).\label{fig:gershgorin}}
\end{figure}

\section{Computing Eigenvalues: A First Look}
\label{sec:computing-eigenvalues}

\subsection{The Power Method}

\Cref{prop:dominant} suggests a method for the dominant eigenvalue and its eigenvector: apply the matrix again and again. To keep the numbers in range, the vector is normalised after each step, and the eigenvalue is estimated by the \term{Rayleigh quotient}
\[
r_A(\mathbf x)\coloneqq\frac{\mathbf x^{\mathsf H}A\mathbf x}{\mathbf x^{\mathsf H}\mathbf x}\qquad(\mathbf x\ne\mathbf 0),
\]
which \cref{sec:lagrange} met for symmetric matrices; if $\mathbf x$ is an eigenvector for $\lambda$, then $r_A(\mathbf x)=\lambda$. \Cref{alg:power-method} states the \term{power method}.

\begin{algorithm}
\begin{algorithmic}[1]
\Require $A\in\mathbb C^{n\times n}$, a starting vector $\mathbf x_0\ne\mathbf 0$, a number of steps $K$
\Ensure an approximation $\mu$ of the dominant eigenvalue and a unit vector $\mathbf x$ approximately in its eigenspace
\Name{PowerMethod}{A,\mathbf x_0,K}
\State $\mathbf x\gets\mathbf x_0/\|\mathbf x_0\|$
\For{$k=1,\ldots,K$}
  \State $\mathbf y\gets A\mathbf x$ \Comment{the only use of $A$}
  \State $\mu\gets\mathbf x^{\mathsf H}\mathbf y$ \Comment{Rayleigh quotient, since $\|\mathbf x\|=1$}
  \State $\mathbf x\gets\mathbf y/\|\mathbf y\|$
\EndFor
\State \Return $\mu$ and $\mathbf x$
\end{algorithmic}
\caption{The power method (\cref{thm:power-method}). Each step costs one product of $A$ with a vector, $2n^2-n$ operations for a full real matrix, and about $5n$ further operations; in practice the loop stops when $\mu$ changes little. If $\mathbf y=\mathbf 0$ occurs, $\mathbf x$ is an eigenvector for $0$.}
\label{alg:power-method}
\end{algorithm}

The method uses $A$ only through products $A\mathbf x$, so it applies to matrices that are too large to be stored, as long as they can be multiplied with vectors, which is how the PageRank vector of \cref{sec:markov} has been computed. Its accuracy is measured by the angle between the computed vector and the eigenvector. For nonzero $\mathbf x$ and $\mathbf y$ in $\mathbb C^n$, let $\hat{\mathbf y}=\mathbf y/\|\mathbf y\|$ and
\[
\sin\angle(\mathbf x,\mathbf y)\coloneqq\frac{\|\mathbf x-\langle\hat{\mathbf y},\mathbf x\rangle\hat{\mathbf y}\|}{\|\mathbf x\|},
\]
the length of the component of $\mathbf x/\|\mathbf x\|$ orthogonal to $\mathbf y$ (\cref{prop:projection-line}); for real vectors it is the sine of the angle between the two lines. It does not change when $\mathbf x$ or $\mathbf y$ is multiplied by a nonzero number, so it measures the distance between the directions only. The Rayleigh quotient of a vector close in direction to an eigenvector is close to the eigenvalue, and much closer for Hermitian matrices.

\begin{proposition}[Accuracy of the Rayleigh quotient]
\label{prop:rayleigh-error}
Let $A\mathbf x_1=\lambda_1\mathbf x_1$ with $\mathbf x_1\ne\mathbf 0$, let $\mathbf x\ne\mathbf 0$ and $s\coloneqq\sin\angle(\mathbf x,\mathbf x_1)$. Then
\[
|r_A(\mathbf x)-\lambda_1|\le\|A-\lambda_1I\|_2\,s ,
\]
and if $A$ is Hermitian, even
\[
|r_A(\mathbf x)-\lambda_1|\le\|A-\lambda_1I\|_2\,s^2 .
\]
\end{proposition}

\begin{proof}
Both sides are unchanged if $\mathbf x$ and $\mathbf x_1$ are scaled, so let $\|\mathbf x\|=\|\mathbf x_1\|=1$, and write $\mathbf x=c\mathbf x_1+\mathbf w$ with $c=\langle\mathbf x_1,\mathbf x\rangle$ and $\mathbf w\perp\mathbf x_1$; then $\|\mathbf w\|=s$. Since $\mathbf x^{\mathsf H}\mathbf x=1$ and $(A-\lambda_1I)\mathbf x_1=\mathbf 0$,
\[
r_A(\mathbf x)-\lambda_1=\mathbf x^{\mathsf H}(A-\lambda_1I)\mathbf x=\mathbf x^{\mathsf H}(A-\lambda_1I)\mathbf w ,
\]
whose modulus is at most $\|\mathbf x\|\,\|A-\lambda_1I\|_2\|\mathbf w\|=\|A-\lambda_1I\|_2\,s$ by the Cauchy--Schwarz inequality. If $A$ is Hermitian, $\lambda_1$ is real (\cref{prop:hermitian-eigen}), so $A-\lambda_1I$ is Hermitian and $\mathbf x_1^{\mathsf H}(A-\lambda_1I)=\bigl((A-\lambda_1I)\mathbf x_1\bigr)^{\mathsf H}=\mathbf 0^{\mathsf T}$; hence $\mathbf x^{\mathsf H}(A-\lambda_1I)\mathbf w=\mathbf w^{\mathsf H}(A-\lambda_1I)\mathbf w$, whose modulus is at most $\|A-\lambda_1I\|_2\|\mathbf w\|^2$.
\end{proof}

\begin{theorem}[Convergence of the power method]
\label{thm:power-method}
Let $A\in\mathbb C^{n\times n}$ be diagonalisable with the eigenvalues $|\lambda_1|>|\lambda_2|\ge\cdots\ge|\lambda_n|$ and the eigenvector $\mathbf x_1$ for $\lambda_1$, and let the starting vector have the coordinate $c_1\ne0$ in \eqref{eq:power-expansion}. Let $\mathbf x_k$ be the vector after $k$ steps of \cref{alg:power-method} and $\mu_k\coloneqq r_A(\mathbf x_k)$ its Rayleigh quotient, the value of $\mu$ computed in step $k+1$. There are constants $C_1$ and $C_2$, independent of $k$, with
\[
\sin\angle(\mathbf x_k,\mathbf x_1)\le C_1\Bigl|\frac{\lambda_2}{\lambda_1}\Bigr|^k,\qquad|\mu_k-\lambda_1|\le C_2\Bigl|\frac{\lambda_2}{\lambda_1}\Bigr|^k ,
\]
and if $A$ is Hermitian, $|\mu_k-\lambda_1|\le C_2|\lambda_2/\lambda_1|^{2k}$.
\end{theorem}

\begin{proof}
The vector $\mathbf x_k$ is a positive multiple of $A^k\mathbf x_0$, which is $\lambda_1^k(c_1\mathbf x_1+\mathbf e_k)$ with $\|\mathbf e_k\|\le C|\lambda_2/\lambda_1|^k$ by \cref{prop:dominant}; here $\lambda_1\ne0$ because $|\lambda_1|>|\lambda_2|\ge0$, and $A^k\mathbf x_0\ne\mathbf 0$ for large $k$. For $\mathbf z\coloneqq c_1\mathbf x_1+\mathbf e_k$, the component orthogonal to $\mathbf x_1$ is that of $\mathbf e_k$, whose length is at most $\|\mathbf e_k\|$, and $\|\mathbf z\|\ge|c_1|\|\mathbf x_1\|-\|\mathbf e_k\|\ge|c_1|\|\mathbf x_1\|/2$ as soon as $\|\mathbf e_k\|\le|c_1|\|\mathbf x_1\|/2$. So $\sin\angle(\mathbf x_k,\mathbf x_1)\le2\|\mathbf e_k\|/(|c_1|\|\mathbf x_1\|)$ for all large $k$, and since a sine is at most $1$, a constant $C_1$ covers the finitely many other $k$. The bounds for $\mu_k$ follow from \cref{prop:rayleigh-error} with $C_2=\|A-\lambda_1I\|_2\max(C_1,C_1^2)$.
\end{proof}

The convergence is linear, with the factor $|\lambda_2/\lambda_1|$ per step for the vector, and it is slow when the two largest eigenvalues are close in modulus. \Cref{tab:power-method} shows the method at work. We ran \cref{alg:power-method} in binary64 arithmetic, as a Python program, on the second difference matrix $T_3$, with the eigenvalues $2+\sqrt2$, $2$ and $2-\sqrt2$, and on the matrix
\[
B=ST_3S^{-1}=\begin{bmatrix}1&0&-1\\-1&2&-1\\0&-1&3\end{bmatrix},\qquad S=\begin{bmatrix}1&1&0\\0&1&1\\0&0&1\end{bmatrix},
\]
which has the same eigenvalues but is not symmetric, both from $\mathbf x_0=\mathbf e_1$. The errors were measured against the exact eigenvalue $2+\sqrt2$ and the exact eigenvectors, evaluated in the arbitrary-precision ball arithmetic of the python-flint package. The ratio $|\lambda_2/\lambda_1|=2/(2+\sqrt2)=2-\sqrt2\approx0.5858$ is the same for both matrices. In both, the angle decreases by this factor per step, as the ratios in the last row show. The eigenvalue error of the nonsymmetric matrix decreases at the same rate, while that of the symmetric matrix decreases with the factor $(2-\sqrt2)^2\approx0.3431$ and has reached the level of the rounding errors, about $5.7\times10^{-16}$ or about one unit in the last place of $\lambda_1$, after $35$ steps; the nonsymmetric matrix needs about $60$ steps for an error of $2.4\times10^{-14}$.

\begin{table}
\sdfloat{%
\begin{sdtable}{colspec={rlllll}, row{1-2}={font=\small\bfseries}}
\toprule
 & & \SetCell[c=2]{l} symmetric $T_3$ & & \SetCell[c=2]{l} nonsymmetric $B$ & \\
\cmidrule[lr]{3-4}\cmidrule[lr]{5-6}
$k$ & $(2-\sqrt2)^k$ & angle & eigenvalue & angle & eigenvalue \\
\midrule
$0$ & $1$ & $8.7\times10^{-1}$ & $1.4$ & $9.3\times10^{-1}$ & $2.4$ \\
$5$ & $6.9\times10^{-2}$ & $9.7\times10^{-2}$ & $1.3\times10^{-2}$ & $2.0\times10^{-1}$ & $1.0\times10^{-1}$ \\
$10$ & $4.8\times10^{-3}$ & $6.7\times10^{-3}$ & $6.4\times10^{-5}$ & $1.2\times10^{-2}$ & $9.9\times10^{-3}$ \\
$15$ & $3.3\times10^{-4}$ & $4.6\times10^{-4}$ & $3.0\times10^{-7}$ & $8.5\times10^{-4}$ & $6.9\times10^{-4}$ \\
$20$ & $2.3\times10^{-5}$ & $3.2\times10^{-5}$ & $1.4\times10^{-9}$ & $5.9\times10^{-5}$ & $4.8\times10^{-5}$ \\
$25$ & $1.6\times10^{-6}$ & $2.2\times10^{-6}$ & $6.9\times10^{-12}$ & $4.0\times10^{-6}$ & $3.3\times10^{-6}$ \\
$30$ & $1.1\times10^{-7}$ & $1.5\times10^{-7}$ & $3.3\times10^{-14}$ & $2.8\times10^{-7}$ & $2.3\times10^{-7}$ \\
$35$ & $7.4\times10^{-9}$ & $1.1\times10^{-8}$ & $5.7\times10^{-16}$ & $1.9\times10^{-8}$ & $1.6\times10^{-8}$ \\
$40$ & $5.1\times10^{-10}$ & $7.2\times10^{-10}$ & $5.7\times10^{-16}$ & $1.3\times10^{-9}$ & $1.1\times10^{-9}$ \\
\midrule
ratio at $20$ & $0.5858$ & $0.5858$ & $0.3431$ & $0.5858$ & $0.5858$ \\
\bottomrule
\end{sdtable}}{The power method in binary64 arithmetic for the symmetric matrix $T_3$ and the nonsymmetric matrix $B$ with the same eigenvalues, from $\mathbf x_0=\mathbf e_1$: the sine of the angle between $\mathbf x_k$ and the dominant eigenvector, and the error $|\mu_k-(2+\sqrt2)|$ of its Rayleigh quotient $\mu_k=r_A(\mathbf x_k)$. The last row gives the ratio of the values at $k=20$ and $k=19$.\label{tab:power-method}}
\end{table}

Two refinements make the method practical. By \cref{thm:spectral-mapping}, the matrix $(A-\sigma I)^{-1}$ has the eigenvalues $1/(\lambda_i-\sigma)$, and the largest of them in modulus belongs to the eigenvalue $\lambda_i$ nearest to $\sigma$. The power method applied to it, \term{inverse iteration with the shift $\sigma$}, converges to the eigenvector of that eigenvalue with the factor $|\lambda_i-\sigma|/|\lambda_j-\sigma|$ per step, where $\lambda_j$ is the next nearest eigenvalue, and this factor is small when $\sigma$ is a good approximation of $\lambda_i$. Each step solves a linear system $(A-\sigma I)\mathbf y=\mathbf x$, with one LU factorisation computed once (\cref{sec:lu}) and two triangular solves per step. Replacing $\sigma$ in every step by the current Rayleigh quotient gives the \term{Rayleigh quotient iteration}, which for Hermitian matrices converges cubically~\cite{parlett1998}. The methods that libraries use to compute all eigenvalues of a matrix, the QR algorithm for full matrices and the Lanczos and Arnoldi methods for large sparse ones, are developments of the power method and of inverse iteration~\cite{golub2013,trefethen1997}; Volumes~VIII and~IX treat them, and with them the routines \code{linalg.eig}, \code{linalg.eigvals}, \code{linalg.eigh} and \code{linalg.eigvalsh} of the stardust library and the sparse routines of Volume~IX.

\subsection{Why Not the Characteristic Polynomial}

The eigenvalues are the roots of the characteristic polynomial, and one might compute its coefficients and then find the roots. This is a bad idea, because the roots of a polynomial can depend very sensitively on its coefficients. Wilkinson's famous example~\cite{wilkinson1959,wilkinson1965} is the polynomial
\[
W(x)=(x-1)(x-2)\cdots(x-20)=x^{20}-210x^{19}+\cdots+20!,
\]
whose roots $1,2,\ldots,20$ are well separated; it is the characteristic polynomial of the matrix $\operatorname{diag}(1,\ldots,20)$. Its coefficients are integers of up to $20$ digits, the largest about $1.38\times10^{19}$. Wilkinson changed the coefficient $-210$ of $x^{19}$ to $-210-2^{-23}$, a relative change of about $6\times10^{-10}$. We computed the roots of the changed polynomial with the certified root finder of python-flint, which encloses each root in a disc of known small radius, and give them to the digits shown; \cref{fig:sensitivity}A shows them. The roots $1$ to $5$ move by less than $10^{-6}$, the roots $6$, $7$ and $8$ become $6.000007$, $6.999697$ and $8.007268$, and from $10$ on the roots pair up and leave the real line: the changed polynomial has the ten pairs of complex roots $10.095266\pm0.643501i$, $11.793634\pm1.652330i$, $13.992358\pm2.518830i$, $16.730737\pm2.812625i$ and $19.502439\pm1.940330i$, and the real roots $8.917250$ and $20.846908$.

The sensitivity can be computed. If $r$ is a simple root of $p$, so that $p'(r)\ne0$, the implicit function theorem (\cref{thm:implicit-scalar}) applied to $F(\varepsilon,y)=p(y)+\varepsilon q(y)$ gives a root $r(\varepsilon)$ of $p+\varepsilon q$ that depends differentiably on $\varepsilon$ near $0$, with $r(0)=r$, and differentiating $p(r(\varepsilon))+\varepsilon q(r(\varepsilon))=0$ at $\varepsilon=0$ gives
\[
r'(0)=-\frac{q(r)}{p'(r)} .
\]
For $p=W$ and $q(x)=-x^{19}$, the root $j$ moves to first order by $\varepsilon j^{19}/|W'(j)|$ with $|W'(j)|=(j-1)!\,(20-j)!$. For $j=6$, $7$ and $8$ and $\varepsilon=2^{-23}$ this gives $6.9\times10^{-6}$, $3.0\times10^{-4}$ and $7.1\times10^{-3}$, in agreement with the changes observed; for $j=16$ the factor $16^{19}/(15!\,4!)\approx2.4\times10^9$ predicts a change of about $290$, far beyond the range in which the first-order approximation holds, and the roots leave the real line instead.

Even the rounding of the coefficients to binary64 numbers matters. Of the $21$ integer coefficients of $W$, exactly five, those of $x^3$ to $x^7$, are not binary64 numbers; rounded to the nearest binary64 numbers, they change by between $16$ and $512$, each by a relative amount below $2^{-53}$, the largest change occurring in the coefficient $-1.29\times10^{19}$ of $x^3$. The exact roots of the rounded polynomial, again computed with python-flint, are all real, but they differ from the integers $1,\ldots,20$ by up to $6.2\times10^{-4}$, at the root $13$, and by more than $5\times10^{-4}$ also at $12$ and $14$. The eigenvalues of the companion matrix of the rounded polynomial, computed by \code{linalg.eigvals} in binary64, differ from the integers by up to $0.53$ and include the complex pair $13.4899\pm0.1959i$. In contrast, for the matrix $H\operatorname{diag}(1,\ldots,20)H$, where $H$ is the Householder reflection $I-2\mathbf v\mathbf v^{\mathsf T}/\mathbf v^{\mathsf T}\mathbf v$ with $\mathbf v=(1,2,\ldots,20)$ and the product was formed in binary64 with the matrix product of the library, from left to right, the routines \code{linalg.eigvals} and \code{linalg.eigvalsh} return eigenvalues within $5.7\times10^{-14}$ and $7.1\times10^{-15}$ of the integers. The eigenvalues of this matrix are insensitive to perturbations of its entries, but the passage to the coefficients of the characteristic polynomial makes them extremely sensitive. Eigenvalues are therefore computed from the matrix, by transformations that preserve them, and never from the characteristic polynomial; conversely, the eigenvalues of the companion matrix are a good way to find the roots of a polynomial when the polynomial, and not a matrix, is the given data~\cite{edelman1995}.

\subsection{The Sensitivity of Eigenvalues}

Eigenvalues can be sensitive to perturbations of the matrix itself, and the defective matrices are the extreme case. Let $\varepsilon>0$ and let $J_\varepsilon$ be the Jordan block $J_n(0)$ with $\varepsilon$ added in the lower left corner. It is the transposed companion matrix $C_q^{\mathsf T}$ of $q(x)=x^n-\varepsilon$, so by \cref{prop:companion,prop:charpoly-similar}
\[
p_{J_\varepsilon}(x)=x^n-\varepsilon,
\]
and the eigenvalues are the $n$ numbers $\varepsilon^{1/n}e^{2\pi ik/n}$, on the circle of radius $\varepsilon^{1/n}$. A change of size $\varepsilon$ in one entry moves the eigenvalue $0$ by $\varepsilon^{1/n}$, which is enormous compared with $\varepsilon$: for $n=16$ and $\varepsilon=10^{-16}$, near the unit roundoff, the eigenvalues move to the circle of radius $0.1$. \Cref{tab:jordan-sensitivity} shows the radii for $n=2,4,8,16$, and the eigenvalues of the stored matrices $J_\varepsilon$ computed by \code{linalg.eigvals}, whose moduli agree with $\varepsilon^{1/n}$ to a relative error of at most $1.0\times10^{-13}$. The perturbation does not have to be made on purpose. For the matrix $HJ_n(1)H$, with the Householder reflection $H$ of the previous subsection for $\mathbf v=(1,\ldots,n)$, formed in binary64 in the same way, which is similar to $J_n(1)$ but whose entries are rounded and which no longer looks triangular, the computed eigenvalues lie at distances up to $1.2\times10^{-8}$, $1.2\times10^{-4}$, $2.9\times10^{-3}$ and $7.6\times10^{-2}$ from the exact eigenvalue $1$, numbers of the order of $u^{1/n}$ or smaller. \Cref{fig:sensitivity}B shows the case $n=16$.

\begin{table}
\sdfloat{%
\begin{sdtable}{colspec={rllll}, row{1}={font=\small\bfseries}}
\toprule
$n$ & $\varepsilon^{1/n}$ & computed, rel.\ deviation & $HJ_n(1)H$: largest $|\hat\lambda-1|$ & $u^{1/n}$ \\
\midrule
$2$ & $1.0\times10^{-8}$ & $3.1\times10^{-17}$ & $1.2\times10^{-8}$ & $1.1\times10^{-8}$ \\
$4$ & $1.0\times10^{-4}$ & $8.7\times10^{-16}$ & $1.2\times10^{-4}$ & $1.0\times10^{-4}$ \\
$8$ & $1.0\times10^{-2}$ & $1.2\times10^{-15}$ & $2.9\times10^{-3}$ & $1.0\times10^{-2}$ \\
$16$ & $1.0\times10^{-1}$ & $1.0\times10^{-13}$ & $7.6\times10^{-2}$ & $1.0\times10^{-1}$ \\
\bottomrule
\end{sdtable}}{Sensitivity of a defective eigenvalue. Second column: the modulus $\varepsilon^{1/n}$ of the eigenvalues of the Jordan block $J_n(0)$ with $\varepsilon=10^{-16}$ in the lower left corner; third column: the largest relative deviation from it of the moduli of the eigenvalues computed by \code{linalg.eigvals} for the stored matrix, measured against the exact value for the binary64 number nearest to $10^{-16}$. Fourth column: the largest distance from $1$ of the computed eigenvalues of the matrix $HJ_n(1)H$ formed in binary64; last column: $u^{1/n}$ with $u=2^{-53}$.\label{tab:jordan-sensitivity}}
\end{table}

\begin{figure}
\sdpanel{A}{\begin{tikzpicture}
\begin{axis}[sdplot, width=0.55\textwidth, height=0.36\textwidth, xmin=0, xmax=22, ymin=-3.3, ymax=3.3,
  xtick={0,5,10,15,20}, ytick={-3,0,3}, xlabel={$\operatorname{Re}$}, ylabel={$\operatorname{Im}$}]
\addplot[black!55, only marks, mark=o, mark size=1.6pt, samples at={1,2,...,20}] (x,0);
\addplot[black, only marks, mark=*, mark size=1.3pt] coordinates {(1,0) (2,0) (3,0) (4,0) (5,0) (6.000007,0) (6.999697,0) (8.007268,0) (8.917250,0) (10.095266,-0.643501) (10.095266,0.643501) (11.793634,-1.652330) (11.793634,1.652330) (13.992358,-2.518830) (13.992358,2.518830) (16.730737,-2.812625) (16.730737,2.812625) (19.502439,-1.940330) (19.502439,1.940330) (20.846908,0)};
\end{axis}
\end{tikzpicture}}\hfill
\sdpanel{B}{\begin{tikzpicture}
\begin{axis}[sdplot, width=0.36\textwidth, height=0.36\textwidth, axis equal image, xmin=0.88, xmax=1.12, ymin=-0.12, ymax=0.12,
  xtick={0.9,1,1.1}, ytick={-0.1,0,0.1}, xlabel={$\operatorname{Re}$}, ylabel={$\operatorname{Im}$}]
\addplot[black, dashed, line width=0.5pt, domain=0:360, samples=120] ({1+0.1007*cos(x)},{0.1007*sin(x)});
\addplot[black, only marks, mark=*, mark size=1.3pt] coordinates {(0.9253,0.0000) (0.9325,-0.0322) (0.9325,0.0322) (0.9529,-0.0583) (0.9529,0.0583) (0.9827,-0.0732) (0.9827,0.0732) (1.0001,0.0000) (1.0162,-0.0737) (1.0162,0.0737) (1.0469,-0.0595) (1.0469,0.0595) (1.0682,-0.0331) (1.0682,0.0331) (1.0759,0.0000)};
\addplot[black!55, only marks, mark=o, mark size=2pt] coordinates {(1,0)};
\end{axis}
\end{tikzpicture}}
\caption{Sensitive eigenvalues. (A) The roots of Wilkinson's polynomial $W$ (circles) and of $W$ with the coefficient of $x^{19}$ changed by $-2^{-23}$ (dots), computed with certified error bounds. (B) The eigenvalues of $HJ_{16}(1)H$, a matrix similar to the Jordan block $J_{16}(1)$, computed by \code{linalg.eigvals} from the matrix formed in binary64 (dots; two of them nearly coincide at $1.0001$), with the exact eigenvalue $1$ (circle) and the circle of radius $u^{1/16}\approx0.101$ about it (dashed).\label{fig:sensitivity}}
\end{figure}

For diagonalisable matrices the sensitivity is governed by the eigenvector matrix, by a theorem of Bauer and Fike~\cite{bauerfike1960}.

\begin{theorem}[Bauer--Fike theorem]
\label{thm:bauer-fike}
Let $A=S\Lambda S^{-1}$ with $\Lambda=\operatorname{diag}(\lambda_1,\ldots,\lambda_n)$, let $E\in\mathbb C^{n\times n}$, and let $\|\cdot\|$ be one of the operator norms $\|\cdot\|_1$, $\|\cdot\|_2$, $\|\cdot\|_\infty$. Then every eigenvalue $\mu$ of $A+E$ satisfies
\[
\min_{1\le i\le n}|\mu-\lambda_i|\le\kappa(S)\,\|E\|,\qquad\kappa(S)\coloneqq\|S\|\,\|S^{-1}\|.
\]
\end{theorem}

\begin{proof}
If $\mu$ is an eigenvalue of $A$, the left side is $0$. Otherwise $\mu I-A=S(\mu I-\Lambda)S^{-1}$ is invertible. Let $(A+E)\mathbf x=\mu\mathbf x$ with $\mathbf x\ne\mathbf 0$. Then $(\mu I-A)\mathbf x=E\mathbf x$, so $\mathbf x=S(\mu I-\Lambda)^{-1}S^{-1}E\mathbf x$, and taking norms,
\[
\|\mathbf x\|\le\|S\|\,\|(\mu I-\Lambda)^{-1}\|\,\|S^{-1}\|\,\|E\|\,\|\mathbf x\|=\frac{\kappa(S)\|E\|}{\min_i|\mu-\lambda_i|}\,\|\mathbf x\|,
\]
because the norm of a diagonal matrix is the largest modulus of its entries, for the $1$- and the $\infty$-norm by \cref{thm:norm-formulas} and for the $2$-norm since $\|D\mathbf y\|^2=\sum_i|d_i|^2|y_i|^2\le\max_i|d_i|^2\|\mathbf y\|^2$, with equality for a unit vector $\mathbf e_i$. Dividing by $\|\mathbf x\|$ gives the claim.
\end{proof}

The number $\kappa(S)$ is the condition number of the eigenvector matrix, in the sense of \cref{sec:trace-norms}; it is at least $1$, and it is large when the eigenvectors are nearly dependent, as they are near a defective matrix, where they collapse. For a matrix that can be diagonalised by a unitary $S$, $\kappa_2(S)=1$, and a perturbation of norm $\|E\|_2$ moves every eigenvalue by at most $\|E\|_2$: the eigenvalues are perfectly conditioned. By the spectral theorem of \cref{ch:symmetric-matrices} this is the case for every Hermitian matrix; for the second difference matrices $T_n$ it is already shown by the orthogonal matrix $Q$ of \cref{thm:second-difference-eigen}, and \cref{tab:tn-eigen} shows the consequence for these matrices, whose eigenvalues are known exactly. We computed them with \code{linalg.eigvalsh}, the routine for Hermitian matrices, and with the general routine \code{linalg.eigvals}, and measured the errors against the exact values $2-2\cos(j\pi/(n+1))$, evaluated with python-flint; the machine and the version of the library are those described in \cref{sec:systems-practice}. The largest absolute errors are between $5\times10^{-16}$ and $8\times10^{-14}$, all of them small compared with the eigenvalues near $4$ and with $\|T_n\|_2<4$; those of the general routine, which does not use the symmetry, are larger by factors between $4$ and $21$, and its computed eigenvalues of these real matrices had zero imaginary parts. The smallest eigenvalue, about $\pi^2/(n+1)^2$, has a much larger relative error, up to $1.6\times10^{-10}$ for $n=1000$: an absolute error of size $u\|T_n\|_2$ is a large relative error for an eigenvalue much smaller than $\|T_n\|_2$.

\begin{table}
\sdfloat{%
\begin{sdtable}{colspec={rlllll}, row{1-2}={font=\small\bfseries}}
\toprule
 & & \SetCell[c=2]{l} largest abs.\ error & & \SetCell[c=2]{l} rel.\ error of $\lambda_1$ & \\
\cmidrule[lr]{3-4}\cmidrule[lr]{5-6}
$n$ & $\lambda_1$ & \code{eigvalsh} & \code{eigvals} & \code{eigvalsh} & \code{eigvals} \\
\midrule
$10$ & $8.1\times10^{-2}$ & $5.1\times10^{-16}$ & $2.0\times10^{-15}$ & $2.0\times10^{-15}$ & $8.4\times10^{-15}$ \\
$100$ & $9.7\times10^{-4}$ & $1.9\times10^{-15}$ & $2.1\times10^{-14}$ & $1.5\times10^{-13}$ & $3.1\times10^{-13}$ \\
$1000$ & $9.8\times10^{-6}$ & $3.6\times10^{-15}$ & $7.6\times10^{-14}$ & $1.6\times10^{-10}$ & $3.4\times10^{-11}$ \\
\bottomrule
\end{sdtable}}{Eigenvalues of the second difference matrices $T_n$ computed in binary64 by the routines \code{linalg.eigvalsh} and \code{linalg.eigvals}, compared with the exact values of \cref{thm:second-difference-eigen}: the smallest eigenvalue $\lambda_1$, the largest absolute error over all $n$ eigenvalues, and the relative error of $\lambda_1$.\label{tab:tn-eigen}}
\end{table}

\bigskip
An eigenvector of an operator is a nonzero vector that the operator merely stretches, and the stretching factor is the eigenvalue. The eigenvalues of a matrix are the roots of its characteristic polynomial $\det(\lambda I-A)$, a monic polynomial whose coefficients include the trace and the determinant, which are therefore the sum and the product of the eigenvalues; over $\mathbb C$ there are always $n$ of them, counted with multiplicity, while a real matrix may have complex eigenvalues in conjugate pairs. Every monic polynomial is the characteristic polynomial of its companion matrix, so eigenvalues cannot in general be computed by formulas, and $AB$ and $BA$ have the same nonzero eigenvalues. Eigenvectors for different eigenvalues are independent, and a matrix is diagonalisable exactly when the geometric multiplicity of every eigenvalue equals its algebraic multiplicity, or equivalently when its minimal polynomial has no repeated roots; commuting diagonalisable matrices are diagonalised together. Diagonalisation turns powers and polynomials of a matrix into powers and polynomials of numbers, and explains the growth of the Fibonacci numbers; the powers of a matrix tend to zero exactly when its spectral radius is less than $1$, and the spectral radius is the limit of $\|A^k\|^{1/k}$. Every matrix satisfies its characteristic polynomial. Real matrices with complex eigenvalues rotate and scale invariant planes, unitary matrices have their eigenvalues on the unit circle and Hermitian matrices on the real line, and rotations of space have the eigenvalues $1$ and $e^{\pm i\varphi}$. The second difference matrix, the circulant matrices and the permutation matrices have explicit eigenvalues, a positive Markov matrix has a unique stationary distribution approached at a geometric rate, and Gershgorin's discs locate eigenvalues without computation. In computation, the power method finds a dominant eigenvalue, eigenvalues are never computed from the characteristic polynomial, and defective eigenvalues react to a perturbation $\varepsilon$ like $\varepsilon^{1/n}$. \Cref{ch:symmetric-matrices} proves that every Hermitian matrix has an orthonormal basis of eigenvectors.

\chapternotes

Eigenvalue problems arose in the eighteenth century in mechanics and astronomy, long before matrices, as the problems of finding the principal axes of a quadratic form and the frequencies of small oscillations~\cite{hawkins1975cauchy}. Cauchy proved in 1829, in a memoir on the equation that determines the secular perturbations of the planets, that the eigenvalues of a real symmetric matrix are real, and Hawkins~\cite{hawkins1975cauchy} traces how this result and its relatives grew into a spectral theory of matrices. The German words \emph{Eigenfunktion} and \emph{Eigenwert} appear in the first of Hilbert's communications on integral equations~\cite{hilbert1904}. The Cayley--Hamilton theorem was stated by Cayley in his memoir of 1858, who verified it for matrices of orders $2$ and $3$~\cite{cayley1858}, and first proved in general by Frobenius in 1878~\cite{frobenius1878}; Hamilton had met an instance of it in his work on quaternions~\cite{hamilton1853}. Perron's theorem on positive matrices dates from 1907~\cite{perron1907} and Frobenius' extension to nonnegative matrices from 1912~\cite{frobenius1912}; Gershgorin published his discs in 1931~\cite{gershgorin1931}, and Varga's book~\cite{varga2004} develops the many results that grew from them, while Taussky~\cite{taussky1949} recounts the history of the theorem on diagonally dominant matrices. The power method is commonly credited to von Mises, who described it with Pollaczek-Geiringer in the second part of a paper of 1929~\cite{vonmises1929}. Wilkinson's analysis of his polynomial~\cite{wilkinson1959} and his book on the algebraic eigenvalue problem~\cite{wilkinson1965} established why eigenvalues must be computed from matrices and not from polynomials, and the Bauer--Fike theorem appeared in 1960~\cite{bauerfike1960}. Edelman and Murakami~\cite{edelman1995} analyse the computation of polynomial roots as eigenvalues of companion matrices. Brin and Page described PageRank in 1998~\cite{brin1998}, and Langville and Meyer~\cite{langville2006} treat its mathematics.

For the theory, Axler~\cite[Chaps.~5 and~8]{axler2024} develops eigenvalues, the minimal polynomial and diagonalisability from operators, without determinants, and Hoffman and Kunze~\cite[Chap.~6]{hoffman1971} the characteristic and minimal polynomials and simultaneous diagonalisation in the classical way. Strang~\cite[Chap.~6]{strang2016} and Meyer~\cite[Chap.~7]{meyer2000} present eigenvalues with many applications, among them difference equations, and Meyer~\cite[Chap.~8]{meyer2000} the Perron--Frobenius theory and Markov chains. Horn and Johnson~\cite{horn2013} is the standard reference for the matrix theory: eigenvalues and similarity in Chapter~1, the location and perturbation of eigenvalues, including Gershgorin's theorem and the Bauer--Fike theorem, in Chapter~6, and positive and nonnegative matrices in Chapter~8. On the computational side, Trefethen and Bau~\cite[Lectures~24--27]{trefethen1997} give a short and clear introduction to eigenvalue algorithms, from the power method and inverse iteration to the Rayleigh quotient iteration, Golub and Van Loan~\cite[Chaps.~7 and~8]{golub2013} treat the unsymmetric and the symmetric eigenvalue problem in depth, Parlett~\cite{parlett1998} the symmetric one, and Wilkinson's book~\cite{wilkinson1965} remains a classic on the sensitivity of eigenvalues and the effect of rounding errors.

\chapterreferences
